% ============================================================
% APUNTES INTEGRADOS DE DERECHO PROCESAL CIVIL Y COMERCIAL
% Documento único, autónomo y listo para compilar (pdflatex, dos pasadas).
% ============================================================
\documentclass[12pt,oneside]{book}

\newcommand{\universidad}{Universidad de Buenos Aires}
\newcommand{\facultad}{Facultad de Derecho}
\newcommand{\materia}{Derecho Procesal Civil y Comercial}
\newcommand{\titulo}{Derecho Procesal Civil y Comercial: apuntes integrados}
\newcommand{\titulocorto}{Derecho Procesal Civil y Comercial}
\newcommand{\autor}{Isaac Edgar Camacho Ocampo}
% TODO: verificar consistencia entre fuentes originales (un archivo consigna 2025 y otros, 2026).
\newcommand{\anio}{2026}

\usepackage[utf8]{inputenc}
\usepackage[T1]{fontenc}
\usepackage[spanish,es-nodecimaldot]{babel}
\usepackage[top=2.5cm,bottom=2.5cm,left=3cm,right=2.5cm,headheight=15pt]{geometry}
\usepackage{amsmath}
\usepackage{amssymb}
\usepackage{mathtools}
\usepackage{graphicx}
\usepackage{float}
\usepackage{caption}
\usepackage{subcaption}
\usepackage{booktabs}
\usepackage{array}
\usepackage{longtable}
\usepackage{tabularx}
\usepackage{enumitem}
\usepackage{xcolor}
\usepackage[most]{tcolorbox}
\usepackage{fancyhdr}
\usepackage{ragged2e}
\usepackage[colorlinks=true,linkcolor=blue!50!black,citecolor=green!40!black,urlcolor=blue!60!black,bookmarks=true]{hyperref}

\hypersetup{pdftitle={\titulo},pdfauthor={\autor},pdfsubject={\materia}}

\definecolor{azul}{HTML}{1F3864}
\setcounter{tocdepth}{2}
\setcounter{secnumdepth}{2}
\setlength{\parindent}{0pt}
\setlength{\parskip}{0.5em}
\clubpenalty=10000
\widowpenalty=10000
\setlist[itemize]{topsep=0.3em,itemsep=0.2em}
\setlist[enumerate]{topsep=0.3em,itemsep=0.2em}
\captionsetup{font=small,labelfont=bf}

\newcommand{\abs}[1]{\left|#1\right|}
\newcommand{\norm}[1]{\left\lVert#1\right\rVert}
\newcommand{\vect}[1]{\mathbf{#1}}
\newcommand{\deriv}[2]{\frac{d#1}{d#2}}
\newcommand{\pderiv}[2]{\frac{\partial#1}{\partial#2}}

% ---------- Encabezados de capítulo ----------
\makeatletter
\def\@makechapterhead#1{\vspace*{4pt}{\parindent0pt\raggedright\normalfont\color{azul}
  \ifnum\c@secnumdepth>\m@ne\large\scshape\@chapapp\ \thechapter\par\nobreak\vskip4pt\fi
  \interlinepenalty\@M\LARGE\bfseries #1\par\nobreak\vskip 20pt}}
\def\@makeschapterhead#1{\vspace*{4pt}{\parindent0pt\raggedright\normalfont\color{azul}
  \interlinepenalty\@M\LARGE\bfseries #1\par\nobreak\vskip 20pt}}
\makeatother

% ---------- Cajas semánticas ----------
\tcbset{cajabase/.style={enhanced,breakable,fonttitle=\bfseries\small,boxrule=0.6pt,arc=1.5mm,
  left=6pt,right=6pt,top=4pt,bottom=4pt,before skip=8pt,after skip=8pt}}
\newtcolorbox{definicion}[1]{cajabase,title={Definición: #1},colback=blue!3,colframe=blue!50!black}
\newtcolorbox{conceptoclave}[1]{cajabase,title={Concepto clave: #1},colback=teal!4,colframe=teal!50!black}
\newtcolorbox{importante}{cajabase,title={Importante},colback=yellow!8,colframe=orange!70!black}
\newtcolorbox{ejemplo}[1]{cajabase,title={Ejemplo: #1},colback=green!4,colframe=green!40!black}
\newtcolorbox{regla}[1]{cajabase,title={Norma / regla: #1},colback=purple!4,colframe=purple!50!black}
\newtcolorbox{ejercicio}[1]{cajabase,title={Ejercicio: #1},colback=gray!5,colframe=gray!60!black}
\newtcolorbox{nota}{cajabase,title={Nota},colback=white,colframe=gray!50}
\newtcolorbox{advertencia}{cajabase,title={Advertencia},colback=red!4,colframe=red!55!black}
\newtcolorbox{sintesis}{cajabase,title={Síntesis},colback=azul!4,colframe=azul}

% ---------- Cuadro Cornell (tabla real con tabularx) ----------
\newcolumntype{Q}{>{\raggedright\arraybackslash\bfseries\small}p{0.27\textwidth}}
\newcolumntype{R}{>{\raggedright\arraybackslash\small}X}
\newcommand{\cq}[2]{#1 & #2 \\ \addlinespace}
\newcommand{\cresumen}[1]{\midrule \multicolumn{2}{p{\dimexpr\textwidth-2\tabcolsep\relax}}{\small\textbf{Resumen.} #1}\\ \bottomrule}
\newenvironment{cornell}[1]{%
  \par\bigskip\noindent{\color{azul}\bfseries Cuadro Cornell: #1}\par\smallskip\noindent
  \renewcommand{\arraystretch}{1.15}%
  \tabularx{\textwidth}{QR}\toprule
  \textbf{Preguntas / palabras clave} & \textbf{Notas / respuestas}\\ \midrule}{%
  \endtabularx\par\bigskip}

% ---------- Encabezados y pies ----------
\pagestyle{fancy}
\fancyhf{}
\fancyhead[L]{\small\nouppercase{\leftmark}}
\fancyhead[R]{\small\titulocorto}
\fancyfoot[C]{\thepage}
\renewcommand{\headrulewidth}{0.4pt}
\fancypagestyle{plain}{\fancyhf{}\fancyfoot[C]{\thepage}\renewcommand{\headrulewidth}{0pt}}

\begin{document}

\frontmatter
\begin{titlepage}
\thispagestyle{empty}
\begin{center}
\vspace*{1cm}
{\large\textsc{\universidad}}\\[0.3cm]
{\large\textsc{\facultad}}
\vspace{3cm}

{\Huge\bfseries\color{azul}\materia}
\vspace{0.8cm}

{\LARGE Apuntes integrados}
\vspace{1.2cm}

\rule{0.70\textwidth}{0.5pt}\\[0.5cm]
{\large Material de estudio}\\[0.5cm]
\rule{0.70\textwidth}{0.5pt}
\vfill
{\large\autor}\\[0.3cm]
{\large\anio}
\end{center}
\vfill
\begin{flushleft}
\small\textit{Este es un modesto aporte a los estudiantes de Derecho Procesal Civil; no pretende sustituir el material oficial de la materia.}
\end{flushleft}
\end{titlepage}

\tableofcontents
\mainmatter

\chapter*{Introducción general: el proceso como oficio}
\addcontentsline{toc}{chapter}{Introducción general: el proceso como oficio}
\markboth{Introducción general}{}

Este material reúne, en un único recorrido, los contenidos de Derecho Procesal Civil y Comercial desarrollados a lo largo del curso. No sustituye la bibliografía oficial de la materia: pretende servir como herramienta complementaria de estudio y como referencia para organizar los conceptos.

\textbf{El material de trabajo es la palabra.} El abogado trabaja con enunciados, escritos u orales, que debe interpretar, aplicar o invocar. Su especialidad consiste, fundamentalmente, en manejar el proceso judicial: la herramienta que permite interpretar las normas cuando surge un problema y presentarlas ante un juez para obtener un resultado favorable para el cliente. Por eso el proceso se trata aquí como un \emph{oficio}, intelectual pero con reglas operativas que se aprenden practicando.

\textbf{La imagen conductora es la cadena.} El proceso es una sucesión ordenada de actos, como eslabones, que va desde la demanda hasta la sentencia, con salidas anticipadas posibles. Cada acto debe cerrarse para que el siguiente pueda comenzar.

\begin{ejemplo}{un recorrido completo}
Una persona que necesita demandar a un hospital por una mala praxis debe, primero, conseguir la historia clínica que no tiene (prueba documental y diligencias preliminares). Mientras dura el juicio puede necesitar que no se interrumpa su cobertura de obra social (medidas cautelares). Si un juez resuelve mal y viola una garantía constitucional, será necesario saber cómo reclamarlo hasta la última instancia (reserva del caso federal y recursos extraordinarios). Y a lo largo de todo el recorrido, el criterio propio del profesional es lo que, en buena medida, determina el resultado.
\end{ejemplo}

\textbf{Cómo está organizado.} La Parte~I presenta el proceso, el lenguaje jurídico y las clases de proceso. La Parte~II estudia al juez, la jurisdicción y la competencia. La Parte~III trata los escritos, las resoluciones, las notificaciones y los plazos. La Parte~IV sigue el camino de la demanda y de la defensa. La Parte~V se ocupa de la audiencia preliminar y de la prueba. La Parte~VI estudia las nulidades y los recursos. La Parte~VII reúne otros institutos y modos de terminación del proceso, y la Parte~VIII integra casos prácticos y reflexiona sobre la formación profesional. Un caso recurrente, el accidente náutico en el río Tigre (sección~\ref{sec:tigre}), atraviesa varios capítulos.

\textbf{Utilidad más allá de los tribunales.} Distinguir lo que ocurrió de lo que se puede probar, comunicar con claridad y precisión, responder punto por punto en lugar de en bloque, y entender que toda norma de convivencia implica una consecuencia son habilidades que sirven para negociar, argumentar y decidir en cualquier ámbito. Una demanda mal estructurada, una excepción no opuesta a tiempo o una instancia caduca por descuido pueden costarle un derecho al cliente.

\part{El proceso judicial y su lógica}

\chapter{El abogado, el conflicto y el proceso}

Este capítulo presenta el punto de partida de toda la materia: el conflicto humano, la función del abogado frente a él y el proceso judicial como método para resolverlo. Sus conceptos centrales son el debido proceso, la sentencia y la diferencia entre verdad real y verdad procesal.

\section{El conflicto humano y sus vías de solución}

El \textbf{conflicto humano} es la materia prima sobre la que trabaja el abogado. Los conflictos son inherentes a la naturaleza humana, porque las personas tienen necesidades, emociones e intereses distintos que no siempre resultan conciliables por sí mismas; de allí la necesidad de un tercero imparcial. Quienes acuden a un estudio jurídico lo hacen, por regla general, por una situación problemática (un despido, el cierre de una empresa, un conflicto por el régimen de contacto con los hijos), de naturaleza económica, familiar o personal. Detrás de cada expediente hay una persona que atraviesa una situación real, dimensión que no siempre advierten quienes intervienen desde una distancia burocrática.

No todo hecho genera un conflicto, pero cuando esto ocurre las partes cuentan con distintas vías. Algunas lo resuelven por sí mismas mediante un acuerdo amigable; si no pueden, recurren a un mediador (mediación) o a un arbitraje, mecanismo todavía poco extendido en el ámbito local salvo en algunas empresas, aunque en expansión. Cuando ninguna vía resulta posible o suficiente, queda planteado el proceso judicial.

\begin{importante}
Cuando la justicia estatal no ofrece una respuesta suficiente en tiempo o en profundidad, no es desacertado recurrir a vías alternativas como la mediación o el arbitraje.
\end{importante}

\begin{ejemplo}{efecto de la demora judicial}
Una persona que, tras el fallecimiento accidental de su cónyuge, debe atravesar un proceso de aproximadamente diez años para obtener una indemnización cuantiosa (del orden de cincuenta millones de pesos) sufre un perjuicio autónomo, independiente del resultado económico: la resolución del conflicto, aun cuando no sea la más justa posible, permite continuar con la vida; la incertidumbre prolongada lo impide.
\end{ejemplo}

\begin{ejemplo}{la empresa frente a un reclamo}
Una empresa a la que se le reclama una suma elevada evalúa si conviene resistir el reclamo o llegar a un acuerdo por una suma menor para dar por concluido el conflicto y continuar con su actividad. La misma lógica se aplica a un despido o a la compra de un inmueble.
\end{ejemplo}

\section{La función del abogado}

\subsection{Garante del debido proceso}

La actividad del abogado suele describirse como la defensa de los intereses de su cliente, pero no se agota allí: comprende también el resguardo de las \textbf{garantías constitucionales} dentro del \textbf{debido proceso}, cuya finalidad última es asegurar la \textbf{igualdad} de las partes frente a la resolución de los conflictos. Su función social consiste en minimizar los conflictos y ofrecerles una solución, dado que su resolución, aunque difícil y estresante, permite a las personas continuar con su vida.

El rol de los magistrados es relevante, pero el de quienes efectivamente reclaman y ponen en cuestión una situación es tan importante, o a veces más. El fallo es la consecuencia de un \textbf{litigio} (el conflicto y la manera en que las partes lo plantean): se llega a él litigando.

\subsection{La especialidad es el proceso}

La especialidad del abogado es el manejo del proceso. El contenido normativo de otras disciplinas suele ser mejor conocido por quienes las practican a diario que por el abogado.

\begin{ejemplo}{quienes conocen su materia mejor que el abogado}
Quienes se desempeñan en la actividad aeronáutica (un piloto, por ejemplo) conocen las normas de ese ámbito porque las aplican todos los días; lo mismo ocurre con cuestiones marítimas o con el derecho laboral, donde quienes participan de una negociación sindical suelen dominar la materia mejor que el abogado. Lo que el piloto no sabe es, cuando surge un problema, cómo llevarlo adelante: ese es el trabajo técnico del abogado.
\end{ejemplo}

El conocimiento de las instituciones procesales es imprescindible incluso para quien domina otras ramas, porque frente a un conflicto lo primero que se necesita es saber cómo son los mecanismos para reclamar y hacer efectivo un derecho.

\begin{importante}
Un resultado favorable no equivale necesariamente a ganar el cien por ciento del reclamo: puede consistir en una pérdida menor, una ganancia parcial o, en definitiva, «obtener lo que se puede».
\end{importante}

\subsection{Ética profesional}

El abogado debe representar adecuadamente los intereses de su asistido sin fabricar ni manipular pruebas ni cometer actos ilegales. La responsabilidad de actuar conforme a la ética recae exclusivamente en cada profesional: ninguna circunstancia externa obliga a apartarse de la conducta correcta.

\section{El proceso judicial: concepto y estructura}\label{sec:concepto-proceso}

El término «proceso» remite a una secuencia de etapas que se inicia en un punto y culmina en un resultado. Para explicarlo a alguien ajeno al derecho (un estudiante de física, por ejemplo) puede decirse que es un conjunto de actividades coordinadas para producir un resultado, y representarlo como una \textbf{cadena de eslabones}: el primero es la demanda y el último, la sentencia.

\begin{definicion}{Proceso judicial}
Conjunto de actividades coordinadas, ordenadas para producir un resultado: la \textbf{sentencia}. Tiene un extremo inicial (la demanda) y un extremo final (la sentencia), aunque existen modos de terminación anteriores (capítulo~\ref{cap:modos}). Es una especie de proceso caracterizada por su alto grado de formalidad: está regulado por normas que fijan plazos y requisitos cuyo incumplimiento determina la ineficacia del acto procesal.
\end{definicion}

\begin{nota}
Todo el proceso es una construcción intelectual: no hay nada físico ni proveniente de la naturaleza. Es un invento, como cualquier juego: así como en el fútbol vale jugar con el pie y no con la mano porque así se decidió, el sistema procesal funciona conforme a lo que se decidió que funcione.
\end{nota}

\begin{definicion}{Estructura del proceso}
El proceso se estructura sobre tres posiciones: una \textbf{parte actora} que formula un reclamo; una \textbf{parte demandada}, a la que se da traslado para que conteste; y un \textbf{tercero} (juez o árbitro) que resuelve, no por ocupar una posición de superioridad, sino por el lugar que le corresponde en ese esquema.
\end{definicion}

\section{Justicia, debido proceso e igualdad}

El fin del proceso es la \textbf{sentencia}. Idealmente debería ser justa, pero la justicia es un ideal discutible e inalcanzable en sentido absoluto, porque lo que resulta justo para una parte puede no serlo para la otra: opera como un horizonte hacia el cual el proceso tiende. Lo que sí puede garantizarse es que las reglas sean las mismas para ambas partes.

\begin{definicion}{Debido proceso}
Conjunto de normas procesales que garantizan que las dos partes de un conflicto judicial se encuentren en igualdad de condiciones frente al juez. Es la herramienta con la que el ordenamiento procura aproximarse al ideal de justicia.
\end{definicion}

\begin{importante}
El derecho no asegura la justicia del resultado: garantiza un \emph{método} detallado que asegura la igualdad de las partes ante el tercero que resuelve, mecanismos para intentar convencerlo y, ante una resolución desfavorable e incorrecta, mecanismos para recurrir a otra instancia.
\end{importante}

El juez tiene tanto la facultad como la \textbf{obligación} de resolver el conflicto sometido a su conocimiento. Es fundamental, además, sostener un pensamiento estratégico desde el inicio del proceso (al demandar o al contestar) hasta su finalización (véase el capítulo~\ref{cap:construccion}).

\section{Verdad real y verdad procesal}

¿Coincide lo que establece una sentencia con lo que ocurrió en la realidad? En el mejor de los casos se parecen, pero no son idénticas. Existe una verdad de lo que ocurrió (verdad material) a la que ni los abogados ni el juez tienen acceso directo. Lo que existe es un conjunto de pruebas y una contraposición de relatos, de cuya confrontación surge la sentencia, la llamada \textbf{verdad formal o legal}, que se impone sobre hechos que muchas veces ocurrieron años atrás y no pueden reconstruirse con exactitud absoluta.

\begin{table}[H]
\centering
\caption{Verdad material, prueba y verdad legal formal}
\begin{tabularx}{\textwidth}{>{\raggedright\arraybackslash}X >{\raggedright\arraybackslash}X >{\raggedright\arraybackslash}X}
\toprule
\textbf{Verdad material} & \textbf{Transcurso de la prueba} & \textbf{Verdad legal formal} \\
\midrule
Lo que realmente ocurrió, que ni siquiera se conoce con certeza. & Las pruebas se producen; lo narrado se va probando. Es un acercamiento a la verdad. & Lo que el juez decide y establece que ocurrió; puede parecerse poco o mucho a la verdad material. \\
\bottomrule
\end{tabularx}
\end{table}

El derecho acepta esta diferencia porque, como ciencia social, no puede garantizar el acceso a la verdad plena; aun así puede perfeccionarse mediante mejores investigadores, jueces y abogados. Con suerte ambas verdades coinciden, y ello depende de la calidad de la prueba y de la habilidad de los abogados.

\begin{cornell}{El abogado, el conflicto y el proceso}
\cq{¿Cuál es la función del abogado?}{No solo defender intereses del cliente: también garantizar el respeto de las garantías constitucionales y del debido proceso para que los conflictos se resuelvan en igualdad de condiciones.}
\cq{¿Qué es el proceso judicial?}{Conjunto de actividades coordinadas, formales y dispuestas en etapas, que va de la demanda a la sentencia; se estructura con actor, demandado y un tercero que resuelve.}
\cq{¿Puede garantizarse la justicia?}{No en sentido absoluto: es un ideal. Sí puede garantizarse el debido proceso, es decir, que las reglas sean las mismas para ambas partes.}
\cq{¿En qué se diferencia la verdad material de la verdad legal?}{La material es lo que realmente ocurrió y no se conoce con certeza; la legal es lo que el juez declara en la sentencia tras la prueba. Pueden coincidir en mayor o menor medida.}
\cq{¿Qué significa un «resultado favorable»?}{No necesariamente ganar todo: puede ser una pérdida menor, una ganancia parcial u «obtener lo que se puede».}
\cresumen{Todo conflicto que las partes no resuelven por sí mismas puede terminar en un proceso judicial, cuya función social es ofrecer una solución que permita continuar con la vida. El proceso no garantiza la justicia del resultado sino un método que asegura la igualdad de las partes ante quien resuelve; la especialidad del abogado es manejarlo, y la sentencia expresa una verdad legal que puede diferir de la verdad material.}
\end{cornell}

\chapter{Del hecho a la norma: lenguaje, subsunción y persuasión}

Si el capítulo anterior presentó el método, este examina los materiales con los que opera: la palabra, las normas, los hechos y el modo de persuadir a quien decide.

\section{La palabra como herramienta de trabajo}

El abogado trabaja con un único material verdadero: la palabra, hablada o escrita. Las normas son proposiciones lingüísticas que deben ser interpretadas, aplicadas e invocadas; el material del derecho es siempre textual. Describir un hecho, indicar su adecuación normativa y disponer una argumentación adecuada exige, por tanto, entrenarse en expresarse.

\begin{definicion}{Expresarse}
Utilizar la menor cantidad de palabras posible, con la mayor precisión posible, para transmitir un contenido directo y simple que el otro pueda entender sin esfuerzo. Ayuda mucho tener una estructura mental de lo que se quiere decir y de hacia dónde se quiere ir.
\end{definicion}

\begin{ejemplo}{la materia prima del abogado}
Ante el enunciado penal «el que matare a otro tendrá pena de 15 a 25 años», la función del abogado es interpretarlo y luego aplicarlo o invocarlo, siempre sobre la base de la palabra.
\end{ejemplo}

\section{Estructura lógica de las normas: mandato, permiso y sanción}

Juzgar una cuestión jurídica es confrontar un enunciado normativo con la realidad. Toda norma contiene solo dos posibilidades lógicas: \textbf{obligar} o \textbf{permitir}.

\begin{ejemplo}{la prohibición traducida a mandato}
La norma de tránsito «está prohibido estacionar» equivale, en sentido positivo, a «es obligatorio no estacionar». Del mismo modo, la norma que castiga al «que matare» con 15 a 25 años equivale a decir que es obligatorio no matar, bajo pena de sufrir esa consecuencia.
\end{ejemplo}

\begin{importante}
Cuando una norma no prohíbe ni obliga una conducta sino que la autoriza, es facultativa: deja al destinatario la posibilidad de actuar o no, sin consecuencia en ninguno de los dos casos.
\end{importante}

Para ser verdaderamente una norma debe llevar aparejada una \textbf{sanción}. Las hay económicas (como la indemnización de daños), penales, disciplinarias y sociales; estas últimas son sanciones de hecho, sin consecuencia formal, pero cumplen la función estructural de toda sanción.

\begin{ejemplo}{la sanción social}
Un profesor que llega a clase sin saludar, maltratando verbalmente a los presentes o vistiendo un traje rojo con corbata verde y zapatos anaranjados probablemente no recibiría una sanción formal, pero sí una reacción social clara («mal gusto»). Esa reacción silenciosa es una sanción social.
\end{ejemplo}

\begin{importante}
Toda norma implica una sanción; cuando no la implica, ya no cumple función alguna como tal. De esta idea nace una manera de pensar propiamente jurídica.
\end{importante}

\section{El mundo del ser y el del deber ser}

Frente a un hecho, es necesario encuadrarlo en categorías jurídicas para que resulte idóneo a los fines del derecho. La filosofía del derecho distingue dos mundos:

\begin{itemize}
\item el \textbf{mundo del ser}: los hechos, es decir, los acontecimientos de la realidad (por ejemplo, que una persona tome un arma o cause un daño);
\item el \textbf{mundo del deber ser}: las normas, es decir, los enunciados que prescriben cómo deben ser las cosas (por ejemplo, que está prohibido matar y que quien lo haga tendrá consecuencias, salvo una excepción prevista como la legítima defensa).
\end{itemize}

\begin{definicion}{Subsunción}
Adecuación de un hecho a una norma. La calificación jurídica de un hecho (si una muerte es homicidio o legítima defensa) se produce cuando ese hecho logra subsumirse en la norma.
\end{definicion}

Ante un hecho, el jurista debe mostrar y convencer al juez de que existió un daño (un cambio de estado) y de que existe una causa que lo produjo; esa modificación pertenece a alguno de los dos mundos. Si no hay conducta identificable (un rayo), es muy difícil hallar un responsable; si la hay (la maniobra de una lancha), existe relación de causalidad y un responsable. El concepto de daño se desarrolla en el capítulo~\ref{cap:danos}.

\begin{importante}
El razonamiento ante un caso consiste, en general, en verificar si un hecho produjo un daño subsumible en el principio general de no dañar a otro (\emph{alterum non laedere}), consagrado como fundamento de la responsabilidad civil en el Código Civil y Comercial de la Nación.
\end{importante}

El código procesal estructura la demanda con esta misma lógica: primero se ubica el mundo del ser (los hechos efectivamente ocurridos) y después la pretensión, que corresponde al mundo del deber ser. Esa estructura apunta a que exista un «espejo» entre la demanda y la contestación, de modo que ambas partes puedan reclamar y defenderse en igualdad de condiciones, en un discurso emparentado con la retórica clásica griega y romana.

\subsection{Cómo razona cada rama del derecho}

\begin{table}[H]
\centering
\caption{Modo de razonar según la rama del derecho}
\begin{tabularx}{\textwidth}{>{\raggedright\arraybackslash\bfseries}p{0.28\textwidth} X}
\toprule
Rama & Cómo se piensa \\
\midrule
Derecho civil (daños) & Las cosas eran de un modo, se produjo un cambio y ahora son de otro; la diferencia es el daño. Luego corresponde analizar cómo se cuantifica y cómo se prueba. \\
Derecho administrativo & Desde la posición del administrado: la administración puede ser arbitraria al decidir; existen principios que debió respetar. Si no los respetó y ello causó un daño, eso se reclama. \\
Derecho penal & Se analiza si el hecho constituye o no un delito y si la persona a quien se atribuye es o no responsable. \\
\bottomrule
\end{tabularx}
\end{table}

\section{El discurso jurídico como discurso persuasivo}\label{sec:retorica}

Los juicios son una manifestación de lo que la tradición griega llamaba \emph{el discurso}, cuya estructura fija un contenido mínimo sin impedir agregar otras cuestiones. Todo discurso persuasivo busca convencer a un destinatario: en el político, el electorado; en el jurídico, el juez. El deseo de convencer es propio de la vida en sociedad: cada vez que se solicita algo a otra persona se le ofrece una razón, construida a través de una situación retórica.

\begin{nota}
Una anécdota: en una ciudad bilingüe, quienes pedían dinero en la calle repetían su pedido en un idioma y, ante la falta de respuesta, insistían en el otro. La estructura (persuadir para obtener algo) es la misma que emplea cualquier persona que pide algo, dentro o fuera del ámbito jurídico.
\end{nota}

El discurso jurídico procura el convencimiento por dos vías: una \textbf{lógica}, vinculada con la aplicación de las normas, y otra \textbf{emocional}, porque los seres humanos, incluidos los jueces, evalúan también a través de las emociones. Aunque haya juicios obvios, en la mayoría de los casos lo ideológico y lo emocional está presente durante todo el proceso de decisión. Este componente reaparece al estudiar el daño moral y la elección del fuero (secciones~\ref{sec:moral-estrategia} y~\ref{sec:emotivo-fuero}).

\begin{definicion}{Pregunta retórica}
Recurso del discurso persuasivo por el cual quien lo formula se pregunta y se responde a sí mismo para reforzar lo que afirma. Su función es enfatizar y dejar al interlocutor sin posibilidad de disimular la arbitrariedad de una decisión.
\end{definicion}

\begin{ejemplo}{pregunta retórica frente a una resolución arbitraria}
En el incidente de revinculación descripto en la sección~\ref{sec:gesell}, el planteo contra la decisión que negaba todo contacto del padre con sus hijos se construyó así: «Señor juez, ¿dónde está escrito que el interés superior del niño consista en no ver nunca más a su papá, ni a su familia paterna? ¿Dónde está escrito que usted, de esta manera, está ayudando al niño? Esto no es una sentencia: es, poco menos, una falta de respeto hacia el menor. ¿Dónde encontró usted en la ley que el papá no pueda ver una fotografía de su hijo, o comunicarse con él por un medio audiovisual? ¿Cuál es el riesgo que usted encuentra en que el papá pueda ver a la nena a través de una pantalla?»
\end{ejemplo}

\begin{cornell}{Lenguaje, norma y persuasión}
\cq{¿Con qué trabaja el abogado?}{Con la palabra, hablada o escrita: las normas son enunciados que se interpretan, aplican e invocan.}
\cq{¿Qué dos funciones tiene toda norma?}{Obligar (incluida la prohibición, que equivale a obligar a no hacer) o permitir. Toda norma implica una sanción, formal o social.}
\cq{¿Qué distingue el ser del deber ser?}{El ser son los hechos; el deber ser, las normas. La calificación jurídica surge de la subsunción del hecho en la norma.}
\cq{¿Por qué se narran primero los hechos y luego se pide?}{Porque la demanda sigue la lógica ser–deber ser, y porque debe haber un espejo entre demanda y contestación.}
\cq{¿Qué es una pregunta retórica?}{Recurso persuasivo en que quien habla se pregunta y se responde para enfatizar la arbitrariedad de una decisión.}
\cresumen{El derecho se trabaja con la palabra. Las normas obligan o permiten y llevan una sanción; el jurista subsume hechos (ser) en normas (deber ser) y debe convencer al juez por vías lógicas y emocionales, de manera honesta, clara y precisa.}
\end{cornell}

\chapter{Las clases de proceso y la organización del Código}\label{cap:clases}

Una vez presentado el proceso en general, corresponde distinguir sus tipos y ubicarlos dentro del Código Procesal Civil y Comercial de la Nación (CPCCN).

\section{Parte general y parte especial}

El CPCCN se organiza en una \textbf{parte general} y una \textbf{parte especial}. La general es aplicable a todos los tipos de proceso, se divide en libros y comienza con la regulación del órgano judicial, en particular la figura del juez (capítulo~\ref{cap:juez}). La especial contiene la regulación específica de cada tipo procesal. Hay dos grandes tipos de proceso, el de conocimiento (ordinario) y el ejecutivo, distinción fundamental para toda la teoría procesal.

\section{Proceso de conocimiento y proceso ejecutivo}

El \textbf{proceso de conocimiento} se caracteriza por la amplia capacidad probatoria de las partes, las amplias facultades del juez para conocer los hechos y ordenar prueba, y las amplias facultades de las partes para argumentar. Su objeto es conocer los hechos, es decir, la versión de lo ocurrido que cada parte aporta.

\begin{importante}
En el proceso de conocimiento no se conocen los hechos tal como ocurrieron, sino la versión que cada parte relata y cree que sucedió. El objetivo es reconstruirla de la mejor manera posible, aplicar el derecho y arribar a la sentencia.
\end{importante}

El \textbf{proceso ejecutivo} opera sobre otro presupuesto: la existencia de un \textbf{título ejecutivo}. El juez no examina los hechos que dieron origen al título sino su \textbf{validez formal}: la existencia de una fecha, una firma y los demás requisitos que la ley exige para cada instrumento. Los títulos ejecutivos existen para agilizar la circulación comercial y dar seguridad jurídica, evitando que cada operación deba resolverse mediante un proceso de conocimiento, más lento y trabajoso.

\begin{ejemplo}{el cheque sin fondos}
Un vendedor vende una sustancia y recibe un cheque que resulta sin fondos; puede iniciar la ejecución del cheque. El título no expresa la causa de la obligación y el juez solo examinará su validez formal, sin indagar el negocio subyacente.
\end{ejemplo}

\begin{ejemplo}{los quintales de trigo con gorgojo}
Una persona compra cinco quintales de trigo y paga con cheque. Al recibir la mercadería descubre que está infestada de gorgojo y solo se recupera un quintal. No puede negarse a pagar el cheque: el título es válido y puede ser ejecutado. Lo que podrá hacer después es iniciar un proceso de conocimiento contra el vendedor por el incumplimiento de la compraventa, una causa distinta de la validez del título.
\end{ejemplo}

\begin{table}[H]
\centering
\caption{Proceso de conocimiento y proceso ejecutivo}
\begin{tabularx}{\textwidth}{>{\raggedright\arraybackslash\bfseries}p{0.2\textwidth} X X}
\toprule
 & Conocimiento & Ejecutivo \\
\midrule
Presupuesto & Un conflicto sobre hechos. & Un título ejecutivo. \\
Qué se examina & Los hechos que relata cada parte, con amplia prueba y argumentación. & La validez formal del título (fecha, firma, requisitos legales). \\
Cosa juzgada & Material. & Formal respecto del título. \\
Trámite referido & Ordinario, sumario, sumarísimo. & Traslado al demandado por 10 días para oponer excepciones; si no las hay o se rechazan, sentencia de remate. \\
\bottomrule
\end{tabularx}
\end{table}

\section{Cosa juzgada formal y material}

En el proceso de conocimiento la sentencia, una vez firme, produce \textbf{cosa juzgada}: la cuestión no puede volver a discutirse.

\begin{importante}
La sentencia del proceso ejecutivo hace cosa juzgada \emph{formal} respecto del título, pero no impide que el condenado inicie después un proceso de conocimiento para discutir la causa del negocio subyacente (por ejemplo, reclamar por la mercadería defectuosa). La sentencia del proceso de conocimiento produce cosa juzgada \emph{material}: agota definitivamente la controversia y no puede discutirse en ningún proceso posterior.
\end{importante}

\section{Ordinario, sumario y sumarísimo}

Además del ordinario existen otros procesos de conocimiento: el \textbf{sumario} y el \textbf{sumarísimo}. Si bien el sumario está derogado en la práctica en la Capital Federal, el Código conserva referencias a él y continúa vigente en varias provincias, especialmente en la de Buenos Aires. La diferencia radica en la \textbf{celeridad}: tienen plazos más cortos, menos prueba admisible y menos recursos, priorizando la velocidad sobre la extensión del debate. En todos se busca conocer los hechos, a diferencia del ejecutivo, donde se busca conocer la validez del título.

\section{El proceso de ejecución de sentencia}

Tanto el proceso de conocimiento como el ejecutivo concluyen con una sentencia de condena. Si la parte condenada no cumple voluntariamente dentro del plazo, corresponde iniciar el \textbf{proceso de ejecución de sentencia}, que es común a ambos porque lo que se busca es hacer efectivo lo que el juez mandó.

\begin{ejemplo}{los 100 pesos y los 30 días}
Una sentencia hipotética condena al demandado a pagar 100 pesos en 30 días. Si transcurrido el plazo no paga, corresponde iniciar la ejecución de sentencia, igual que si, en un juicio ejecutivo, se ordenara pagar un cheque y transcurrieran los 10 días previstos sin pago.
\end{ejemplo}

La diferencia entre ambos procesos reside solo en el camino previo (estudio de hechos y prueba, o examen de un documento con requisitos formales); una vez dictada la sentencia, la etapa de ejecución es idéntica en su lógica. La capacidad de emplear la fuerza para ejecutar se vincula con la jurisdicción (sección~\ref{sec:jurisdiccion}).

\begin{cornell}{Clases de proceso}
\cq{¿Cómo se organiza el CPCCN?}{En parte general (órgano judicial; aplicable a todo proceso) y parte especial (regulación de cada tipo procesal).}
\cq{¿Qué se conoce en el proceso de conocimiento?}{La versión de los hechos de cada parte, con amplia prueba, para aplicar el derecho y dictar sentencia.}
\cq{¿Qué examina el juez en el proceso ejecutivo?}{La validez formal del título ejecutivo, no el negocio que le dio origen.}
\cq{¿Cosa juzgada formal o material?}{En el ejecutivo es formal (puede discutirse luego la causa en un proceso de conocimiento); en el de conocimiento es material (agota la controversia).}
\cq{¿En qué se distinguen ordinario, sumario y sumarísimo?}{En la celeridad: plazos más cortos, menos prueba y menos recursos en los dos últimos. Todos son de conocimiento.}
\cresumen{Hay dos grandes tipos de proceso: el de conocimiento, que reconstruye hechos, y el ejecutivo, que verifica la validez de un título. Ambos pueden terminar en una sentencia de condena que, si no se cumple voluntariamente, se hace efectiva por un proceso de ejecución común a los dos.}
\end{cornell}

\part{El juez y su ámbito de actuación}

\chapter{Jurisdicción, imperio y estructura del juzgado}\label{cap:juez}

El juez es el actor central del proceso. Este capítulo explica qué lo distingue de cualquier otro funcionario, cómo se resguarda su imparcialidad y cómo se organiza el juzgado en que trabaja.

\section{El juez y la jurisdicción}\label{sec:jurisdiccion}

El juez es un tercero imparcial que adopta una decisión sobre la base de las posiciones sostenidas por cada parte. A diferencia de otros funcionarios públicos, posee \textbf{jurisdicción}. Por eso existe la sentencia: si la parte condenada no cumple espontáneamente, el juez puede hacérsela cumplir, incluso por la fuerza.

\begin{definicion}{Jurisdicción}
Facultad exclusiva del juez de decir el derecho aplicable a un caso (\emph{jus dicere}) y de aplicarlo y ejecutarlo mediante el \textbf{imperio}, esto es, la posibilidad de emplear la fuerza pública para hacer cumplir lo resuelto. Ningún otro funcionario público reúne ambas potestades.
\end{definicion}

\begin{ejemplo}{ejecución forzada de una sentencia civil}
Ante el incumplimiento de una deuda dineraria entre particulares, el acreedor puede iniciar un juicio que concluya en una sentencia condenatoria, notificada por cédula, que fije un plazo (por ejemplo, diez días) para pagar con más sus intereses. Vencido el plazo sin pago, el juez puede disponer la ejecución forzada, incluida la venta de bienes del deudor. Esa capacidad es lo que lo distingue de los demás funcionarios.
\end{ejemplo}

En el derecho administrativo, toda decisión de la administración admite revisión judicial; de lo contrario se reconocería jurisdicción a los funcionarios administrativos. La garantía consiste en que la ejecución final de cualquier decisión requiera la intervención de un juez.

\begin{ejemplo}{juicio de apremio}
Cuando la Administración Federal de Ingresos Públicos (AFIP) pretende cobrar impuestos o aplicar multas, no puede ejecutar directamente: debe recurrir a un juez, único habilitado para disponer la ejecución. Ese proceso se denomina \textbf{juicio de apremio}.
\end{ejemplo}

\section{Imparcialidad, deberes y facultades del juez}

Tras regular la competencia (capítulo~\ref{cap:competencia}), el Código regula las \textbf{recusaciones} y las \textbf{excusaciones}, institutos cuya finalidad es evitar que el juez actúe con parcialidad. Luego regula los \textbf{deberes y facultades del juez}, entre los que figuran el dictado de sentencias y de resoluciones intermedias. Por sobre esas funciones, el juez es el \textbf{director del proceso}: debe velar por el cumplimiento de las normas procesales, cuya finalidad constante es asegurar la igualdad de condiciones entre las partes.

\section{La estructura del juzgado}

\begin{table}[H]
\centering
\caption{Jerarquía del juzgado}
\begin{tabularx}{\textwidth}{>{\bfseries\raggedright\arraybackslash}p{0.28\textwidth} X}
\toprule
Función & Descripción \\
\midrule
Mesa de entradas & Recibe y registra los escritos; es la primera atención al público y el vínculo con los funcionarios. \\
Jefe de despacho / prosecretario & Puede dictar providencias simples, de mero trámite; no puede resolver controversias. \\
Secretario & Tiene atribuciones para cuestiones de trámite de mayor complejidad, pero tampoco resuelve controversias entre las partes. \\
Juez & Máxima autoridad del juzgado; resuelve todas las controversias mediante resoluciones fundadas y dicta sentencias definitivas e interlocutorias. \\
\bottomrule
\end{tabularx}
\end{table}

Los tipos de resoluciones que dicta cada uno se estudian en el capítulo~\ref{cap:escritos}.

\begin{cornell}{Jurisdicción y juzgado}
\cq{¿Qué es la jurisdicción?}{La facultad exclusiva del juez de decir el derecho y ejecutarlo mediante el imperio (la fuerza pública).}
\cq{¿Por qué la administración necesita al juez para ejecutar?}{Porque toda decisión administrativa es revisable judicialmente y su ejecución final requiere un juez; ej.: juicio de apremio.}
\cq{¿Para qué sirven recusación y excusación?}{Para evitar que el juez actúe con parcialidad.}
\cq{¿Quién puede resolver controversias en el juzgado?}{Solo el juez; el jefe de despacho y el secretario dictan providencias de trámite.}
\cresumen{El juez es un tercero imparcial que tiene jurisdicción e imperio, director del proceso que asegura la igualdad de las partes; en el juzgado, solo él resuelve controversias.}
\end{cornell}

\chapter{Competencia y organización de la justicia}\label{cap:competencia}

La jurisdicción es una facultad general; la competencia señala dónde puede ejercerla cada juez. Este capítulo recorre los criterios de competencia, el reparto entre justicia federal y ordinaria, la particular organización de la Ciudad de Buenos Aires y la forma en que la demanda fija la competencia.

\section{Concepto y criterios}

La \textbf{competencia} es un límite a la jurisdicción: el conjunto de asuntos y los límites dentro de los cuales un juez determinado puede resolver un conflicto concreto.

\begin{table}[H]
\centering
\caption{Criterios de competencia}
\begin{tabularx}{\textwidth}{>{\bfseries\raggedright\arraybackslash}p{0.28\textwidth} X}
\toprule
Criterio & Descripción \\
\midrule
Territorio & El lugar donde ocurrió el hecho o donde tiene domicilio alguna de las partes determina el juez territorial competente. \\
Fuero (materia) & La naturaleza del conflicto (civil, penal, comercial, laboral) determina el fuero; un homicidio corresponde a un juez penal, una deuda a uno civil. \\
Improrrogabilidad (regla) & Cada asunto tiene predeterminado un juez; no se puede elegir libremente a cuál acudir. \\
Prórroga (excepción) & Si el conflicto es estrictamente patrimonial, sin orden público, menores ni familia, las partes pueden prorrogar la competencia a otro tribunal. \\
\bottomrule
\end{tabularx}
\end{table}

\begin{nota}
Es habitual que en contratos entre empresas se pacte que los tribunales competentes serán los nacionales de la Capital Federal. Es admisible porque se trata de cuestiones estrictamente patrimoniales, respecto de las cuales la prórroga está permitida.
\end{nota}

\section{Justicia federal y justicia ordinaria}

En todo el país coexisten dos órdenes superpuestos, el federal y el ordinario; no hay zonas con uno sin el otro. La justicia federal es de aplicación excepcional, de modo que la mayoría de las causas son ordinarias.

\begin{importante}
Los artículos 116 y 117 de la Constitución Nacional determinan las cuestiones de competencia federal. Por ejemplo: un conflicto entre dos provincias corresponde a la Corte Suprema de Justicia de la Nación (de manera originaria o por apelación, según el caso); una cuestión que involucre a un cónsul o embajador corresponde a la justicia federal; y lo mismo ocurre, en principio, con telecomunicaciones o ferrocarriles.
\end{importante}

La justicia ordinaria depende de las provincias, por tratarse de una facultad no delegada al gobierno federal. Las provincias fueron en algún momento estados independientes (Buenos Aires funcionó como estado propio antes de 1860) y, al constituirse en Nación mediante un pacto, delegaron ciertos poderes (leyes de fondo, organización de la justicia federal, administración a cargo del poder ejecutivo nacional) y se reservaron otros, entre ellos organizar su propia administración de justicia. Por eso cada provincia tiene su Constitución, sus tres poderes y un proceso propio, aunque con estructura similar: primera instancia, cámara de apelaciones y, eventualmente, una Corte o Superior Tribunal provincial.

Las cuestiones ordinarias surgen del Código Civil y Comercial, del Código Penal y de leyes especiales de aplicación uniforme, como la Ley de Contrato de Trabajo, que funcionan como un código: la ley de fondo es la misma en toda la Nación, pero se aplica mediante procesos diferentes según cada jurisdicción. En supuestos dudosos (un choque en una ruta nacional) la respuesta no depende de contra quién se litiga sino de qué norma es aplicable; en general tramita ante la justicia ordinaria, aunque la cuestión admite discusión.

\subsection{Los distritos judiciales de la Provincia de Buenos Aires}

La Provincia está organizada en distritos judiciales que no coinciden necesariamente con los partidos. El de San Isidro comprende Tigre, Vicente López, San Isidro y San Fernando; el de Lomas de Zamora, ese partido y otros vecinos. En toda la provincia rige el mismo Código Procesal, pero los tribunales tienen asiento en un municipio determinado de cada distrito.

\section{La justicia de la Ciudad de Buenos Aires}

\subsection{Origen histórico}

\begin{nota}
Hacia 1880, mediante la Ley Avellaneda, se federalizó la ciudad de Buenos Aires. Hasta entonces las primeras autoridades nacionales tenían asiento en un territorio que formaba parte de la Provincia de Buenos Aires, con superposición de poderes. La provincia cedió el territorio, que pasó a ser la Capital Federal, y se creó La Plata como nueva capital provincial.
\end{nota}

Desde entonces existieron en la Ciudad jueces «nacionales», distintos de los federales, que resolvían las cuestiones ordinarias de sus habitantes, pues la Ciudad, al no ser provincia, carecía de gobierno propio que organizara su justicia. Un efecto curioso: el juez civil de la Capital termina siendo pagado con impuestos de todo el país y no solo de los habitantes de la Ciudad.

\subsection{La autonomía porteña}

Antes de 1994 la Ciudad se gobernaba por una intendencia cuyo titular designaba el presidente. Con la reforma constitucional de 1994 adquirió el carácter de Ciudad Autónoma, equiparable en varios aspectos a una provincia, con Constitución y poder legislativo propios y tres senadores. Correspondía que tuviera su propia justicia, pero los jueces nacionales ya en ejercicio se resistieron a integrarla, por lo que se crearon tribunales nuevos: los de lo Contencioso Administrativo y Tributario de la Ciudad, con sus cámaras y, en la cúspide, el Tribunal Superior de Justicia porteño.

\begin{ejemplo}{fuero ordinario o contencioso administrativo porteño}
Quien sufre una mala praxis en un sanatorio privado (el Hospital Alemán, por ejemplo) debe demandar ante la justicia ordinaria. Quien la sufre en un hospital público de la Ciudad (el Hospital Fernández) debe demandar ante el fuero contencioso administrativo y tributario porteño, porque la propia Ciudad es parte demandada.
\end{ejemplo}

\subsection{El fallo «Vinas»}

En las provincias el proceso avanza de primera instancia a la Cámara y luego, eventualmente, al Tribunal Superior o Corte provincial antes de llegar a la Corte Suprema de la Nación. En la Capital, en cambio, los procesos de los tribunales nacionales iban directamente a la Corte.

\begin{importante}
El fallo «Vinas» de la Corte Suprema dispuso que, antes de llegar a ella, los procesos originados en los tribunales nacionales de la Ciudad deben pasar por el Tribunal Superior de Justicia de la Ciudad, generando una instancia intermedia adicional.
\end{importante}

\subsection{La reforma laboral y el traspaso de competencias desde 2027}

A partir de 2027, con la nueva ley, los juicios laborales que se inicien tramitarán ante los nuevos tribunales de la justicia de la Ciudad. El fuero laboral nacional seguirá entendiendo solo en las causas iniciadas hasta el 31 de diciembre de 2026; los iniciados desde el 1 de enero de 2027 tramitarán ante la justicia porteña.
% TODO: verificar referencia presente en las fuentes originales (el apunte no precisa el número o la denominación de la ley ni su régimen transitorio completo).

\section{La demanda fija la competencia}\label{sec:fija-competencia}

La interposición de la demanda fija la competencia (y los hechos controvertidos; véase el capítulo~\ref{cap:inicio}). Si bien el actor no dispone en general de la competencia, el propio ordenamiento procesal le permite, en ciertos supuestos, elegir entre distintas opciones, y en esa elección el abogado define la alternativa más conveniente para su cliente. Interponer la demanda ante un fuero determinado implica proponer esa jurisdicción y fijarla, aunque después la calificación jurídica del hecho no se confirme.

\begin{ejemplo}{el accidente entre San Isidro y Santa Fe}
Un vehículo cuyo titular vive en la Ciudad Autónoma choca contra otro cuya aseguradora tiene sede en Santa Fe; el hecho ocurrió en San Isidro. El ordenamiento concede al actor elegir entre opciones (domicilio del demandado, lugar del hecho o domicilio de la aseguradora, según el caso). Esa elección queda definitivamente fijada una vez presentada la demanda.
\end{ejemplo}

\subsection{Competencia civil y competencia penal}

La competencia penal no puede prorrogarse, modificarse ni elegirse: se determina exclusivamente por el lugar en que se produjo el hecho. En la Provincia de Buenos Aires, a diferencia de la Ciudad, la investigación penal preparatoria está a cargo del fiscal, que asienta los hechos y reúne las pruebas para determinar si hubo un hecho delictivo y si puede atribuirse a una persona; existe siempre un juez de garantías que vela por que la instrucción respete las garantías procesales, y ante quien pueden recurrir el querellante o la defensa. La relación con el proceso civil se estudia en la sección~\ref{sec:civil-penal}.

\subsection{La ley aplicable y la calificación del hecho}

La calificación jurídica puede incidir en la competencia. En el caso del río Tigre (sección~\ref{sec:tigre}), además del Código Civil existe otra norma en discusión: la Ley de Navegación, en tanto el hecho pueda calificarse como \textbf{abordaje}. Existe discusión doctrinaria y jurisprudencial acerca de si corresponde calificar como abordaje un siniestro entre una lancha comercial de gran porte y una embarcación pequeña; esa calificación puede atraer la intervención de la justicia federal. No se trata de un choque de esquina: si la lancha hubiese pasado por encima de la plataforma sin llegar a haber contacto entre los cascos, no se plantearía el mismo problema.

\subsection{El componente emotivo del juez}\label{sec:emotivo-fuero}

La elección del foro puede analizarse también desde el componente emotivo. Una persona ajena a las actividades náuticas y residente en la Ciudad de Buenos Aires no tiene el mismo vínculo con el río que quien reside en la zona del Tigre, de modo que en San Isidro el hecho resulta más cercano y habitual para el juzgador, no porque sea más sensible sino porque el tema le resulta más próximo. El profesional debe saber identificar y manejar ese componente al definir su estrategia.

\begin{cornell}{Competencia y organización judicial}
\cq{¿Qué es la competencia?}{Un límite a la jurisdicción: el conjunto de asuntos que un juez determinado puede resolver. Se determina por territorio y fuero.}
\cq{¿Cuándo se prorroga?}{Solo en cuestiones estrictamente patrimoniales, sin orden público, menores ni familia. La regla es la improrrogabilidad.}
\cq{¿Cómo se reparten justicia federal y ordinaria?}{Coexisten siempre. La federal es excepcional (arts. 116 y 117 CN); la ordinaria depende de las provincias, facultad no delegada.}
\cq{¿Qué rol tuvo el fallo «Vinas»?}{Dispuso que los procesos de los tribunales nacionales de la Ciudad pasen por el Tribunal Superior de la Ciudad antes de llegar a la Corte.}
\cq{¿Quién fija la competencia?}{La demanda la fija; el actor elige solo entre las opciones que el ordenamiento le da. La competencia penal se determina solo por el lugar del hecho.}
\cresumen{La competencia delimita la jurisdicción por territorio y fuero y es, en principio, improrrogable. En el país coexisten justicia federal (excepcional) y ordinaria (provincial); la Ciudad de Buenos Aires, tras su autonomía, suma tribunales propios junto a los nacionales. La demanda fija la competencia, y su elección es una decisión estratégica.}
\end{cornell}

\part{Escritos, resoluciones, notificaciones y plazos}

\chapter{Escritos y resoluciones judiciales}\label{cap:escritos}

Las partes y el juez se comunican de modo distinto dentro del proceso. Este capítulo estudia cómo se redactan los escritos, qué resoluciones dicta el juez, qué diferencia hay entre cargas y obligaciones y en qué consiste la sustanciación.

\section{El lenguaje de las partes y del juez}

Las partes se expresan mediante \textbf{escritos}, cuyo contenido consiste en peticiones; son una suerte de peticionantes que siempre solicitan algo. El juez se expresa mediante \textbf{resoluciones}, en modo imperativo (\emph{ordénase}, \emph{dígase}, \emph{hágase}), porque su función es ordenar.

Cuando una demanda llega al juzgado, habitualmente no la lee primero el juez sino personal con múltiples expedientes pendientes; por eso los escritos deben ser concretos, precisos y ordenados, lo que requiere un genuino trabajo de síntesis.

\section{Estructura formal de un escrito}

\begin{regla}{elementos del escrito judicial, en orden}
\begin{enumerate}
\item \textbf{Título o sumario}: resumen mínimo de lo que se solicita (por ejemplo, «Solicita apertura a prueba»).
\item \textbf{Encabezamiento}: a quién se dirige («Señor Juez») y quién se presenta (el abogado con patrocinio o apoderamiento, o la parte por su propio derecho).
\item \textbf{Individualización de la causa}: autos caratulados y número de expediente.
\item \textbf{Exposición y petitorio}: desarrollo de lo solicitado y pedido concreto al juez.
\end{enumerate}
\end{regla}

El domicilio real y el DNI del presentante se consignan solo en el \textbf{primer escrito} del expediente (usualmente la demanda o la contestación), momento en que se constituye también el domicilio procesal para las notificaciones sucesivas; después no se repiten.

\subsection{Escrito de mero trámite}

Es un escrito simple, en el que la parte informa el estado de los autos y formula un pedido puntual sin argumentación extensa.

\begin{ejemplo}{escrito modelo}
\textit{«Solicita apertura a prueba. Señor Juez: Fulano de Tal, por derecho propio, con domicilio constituido en los autos caratulados `[carátula]', expediente n.\textsuperscript{o} [número], a Vuestra Señoría respetuosamente digo: Que vengo por el presente, conforme al estado de autos, a solicitar se abra la causa a prueba. Proveer de conformidad. Será justicia.»}
\end{ejemplo}

\subsection{Escrito fundado}

Requiere mayor desarrollo argumentativo. Si un testigo clave debe viajar en la fecha de la audiencia testimonial, por razones laborales o médicas, el escrito debe justificar la ausencia y pedir el cambio de fecha. Un primer intento (manifestar que el testigo no puede presentarse por cuestiones laborales hasta fin de enero y pedir que la nueva audiencia no se fije antes) es válido en su estructura, pero puede perfeccionarse con dos elementos:

\begin{enumerate}
\item calificar la situación como \textbf{fuerza mayor} (por ejemplo, la imposibilidad de cambiar pasajes ya adquiridos);
\item sobre todo, explicar \textbf{por qué el testigo es determinante} para la causa; si no, el juez puede denegar el cambio por considerarlo poco fundado.
\end{enumerate}

\begin{importante}
No alcanza con solicitar: hay que fundamentar por qué. Cuanto mayor es la carga argumentativa que exige el pedido (como reprogramar una audiencia ya fijada), mayor debe ser el esfuerzo de justificación, explicando tanto la causa de fuerza mayor como la relevancia de la prueba que se busca preservar.
\end{importante}

\section{Los tipos de resoluciones judiciales}

\begin{regla}{arts. 160 a 163 CPCCN}
Las resoluciones judiciales se clasifican en providencias simples, sentencias interlocutorias y sentencias definitivas. Toda sentencia debe estar razonablemente fundada.
\end{regla}

\begin{table}[H]
\centering
\caption{Tipos de resoluciones judiciales}
\begin{tabularx}{\textwidth}{>{\bfseries\raggedright\arraybackslash}p{0.22\textwidth} X}
\toprule
Tipo & Función, autor y fundamentación \\
\midrule
Providencia simple & Mero trámite: impulsa el proceso sin resolver controversias (por ejemplo, librar un nuevo oficio). La dicta el jefe de despacho o el secretario. No requiere fundamentación extensa. \\
Sentencia interlocutoria & Resuelve una controversia surgida durante el proceso (por ejemplo, un incidente de competencia). Siempre la dicta el juez; requiere sustanciación y fundamentación. \\
Sentencia definitiva & Resuelve el fondo del asunto; es la decisión final. Siempre la dicta el juez, con fundamentación completa y detallada. \\
\bottomrule
\end{tabularx}
\end{table}

Si una petición no genera controversia (pedir un nuevo oficio) corresponde una providencia simple; si surge controversia entre las partes (cuál es el juez competente), debe resolverse mediante sentencia interlocutoria. La fundamentación exigida a las interlocutorias puede ser menos extensa que la de las definitivas, sin que ello implique una diferencia sustancial: en todos los casos debe haber un fundamento (véase el capítulo~\ref{cap:sentencia}).

\section{Cargas procesales y obligaciones}

\begin{definicion}{Carga procesal}
Conducta cuya omisión no genera sanción del juez, pero coloca a quien no la realiza en una posición procesal desfavorable: lo que se pierde es la oportunidad. En el proceso civil, todos los actos que las partes deben realizar son cargas; en la medida en que se ejercen, se incrementan las probabilidades de un resultado favorable.
\end{definicion}

\begin{definicion}{Obligación}
Deber de hacer o no hacer, exigible judicialmente y susceptible de ser compelido por vía coactiva; su incumplimiento genera una consecuencia jurídica directa, como la mora o los daños y perjuicios.
\end{definicion}

\begin{ejemplo}{carga frente a obligación}
No ofrecer prueba testimonial no genera sanción, pero priva a la parte de la posibilidad de acreditar sus afirmaciones con esa prueba. En cambio, el incumplimiento en el pago del alquiler constituye en mora al inquilino y puede derivar en desalojo.
\end{ejemplo}

Las partes no están obligadas a realizar ningún acto procesal en particular, pero cumplir las cargas mejora su posición frente al juez: ofrecer prueba testimonial idónea, aunque no sea obligatorio, coloca a la parte en mejores condiciones de convencerlo.

\section{La sustanciación y la vista}

La \textbf{sustanciación} es el traslado de una petición de una parte a la contraria, para que exprese su posición antes de que el juez resuelva. Así el juez cuenta con ambas posiciones, en resguardo del principio de contradicción y de la igualdad procesal. Debe distinguirse el \textbf{traslado}, que se corre a las partes, de la \textbf{vista}, que se corre a otros funcionarios que intervienen sin ser parte, como el Ministerio Público.

\begin{cornell}{Escritos y resoluciones}
\cq{¿Cómo se expresan partes y juez?}{Las partes con escritos (peticiones); el juez con resoluciones, en modo imperativo.}
\cq{¿Qué lleva un escrito, en orden?}{Sumario, encabezamiento, individualización de la causa y exposición con petitorio. Domicilio real y DNI solo en el primer escrito.}
\cq{¿Qué agrega un escrito fundado?}{Fundamento del pedido: calificación (p. ej., fuerza mayor) y relevancia de lo que se pretende preservar.}
\cq{¿Qué resoluciones existen?}{Providencia simple (trámite), interlocutoria (controversia en el curso del proceso) y definitiva (fondo). Todas las sentencias deben estar fundadas.}
\cq{¿Carga u obligación?}{La carga no se sanciona pero su omisión debilita la posición; la obligación es exigible y su incumplimiento tiene consecuencia jurídica directa.}
\cq{¿Qué es la sustanciación?}{Traslado de una petición a la contraria antes de resolver; garantiza contradicción e igualdad.}
\cresumen{El proceso se comunica mediante escritos con estructura fija y resoluciones de tres tipos. Las partes actúan por cargas, no obligaciones; antes de resolver controversias, el juez sustancia las peticiones y debe fundar sus sentencias.}
\end{cornell}

\chapter{Notificaciones, domicilios, plazos y preclusión}\label{cap:plazos}

Cada paso del expediente llega a las partes según reglas de tiempo que, si se ignoran, hacen perder el derecho a reclamar. Este capítulo estudia cómo se notifica, a qué domicilios, cómo se cuentan los plazos y cómo opera la preclusión.

\section{Formas de notificación}

Notificar es comunicar una resolución judicial. Toda resolución (traslado de demanda, providencia simple, sentencia) debe notificarse de alguna de las formas previstas por el Código:

\begin{enumerate}
\item \textbf{Notificación personal}: típicamente, la parte o su letrado concurren a tomar vista del expediente y quedan notificados en ese acto. Es una estrategia habitual de los actores concurrir a notificarse cuando el demandado no se muestra interesado en hacerlo, para no depender de su diligencia.
\item \textbf{Notificación por cédula}: formulario específico, en papel o digital, hoy mayoritariamente digital. Subsisten excepciones que requieren papel, como el traslado de la demanda cuando el demandado aún no constituyó domicilio procesal.
\item \textbf{Notificación por nota o ministerio de la ley}: toda resolución queda notificada automáticamente los \textbf{martes y viernes}, salvo las que según el artículo 133 CPCCN deban notificarse por cédula o personalmente. Si el martes o viernes fuera feriado, se traslada al primer día hábil siguiente.
\end{enumerate}

Las dos primeras formas cumplen, en la práctica, una función equivalente. La audiencia preliminar se notifica por cédula (art. 363 CPCCN, según el apunte), en su versión electrónica dirigida al domicilio procesal de cada parte.

\section{Los domicilios en el proceso}

\begin{table}[H]
\centering
\caption{Clases de domicilio}
\begin{tabularx}{\textwidth}{>{\bfseries\raggedright\arraybackslash}p{0.24\textwidth} X}
\toprule
Tipo & Descripción \\
\midrule
Real & Aquel en que la persona efectivamente vive, o donde la sociedad está legalmente constituida. \\
Denunciado & El que el actor indica como domicilio real de la contraria al demandar, aún sin verificar. \\
Constituido & El que cada parte fija a fines procesales en un expediente; debe estar dentro del radio del juzgado. \\
Electrónico & El domicilio procesal constituido en el sistema informático del Poder Judicial, que se suma al constituido tradicional. \\
Fiscal & El registrado ante el organismo de recaudación (AFIP/ARCA) o la provincia; no se confunde con el real ni tiene, en principio, efectos procesales sobre la notificación. \\
Legal & El que la ley asigna a ciertas personas por su función; como ejemplo excepcional, el del nuncio apostólico. \\
\bottomrule
\end{tabularx}
\end{table}

\begin{importante}
En el caso de una sociedad demandada, el domicilio al que se notifica debe coincidir con el domicilio real registrado (por ejemplo, el que figura ante la Inspección General de Justicia en el contrato social), no con cualquier otro que figure en otros registros. Verificar el domicilio de cada notificación es el primer paso ante cualquier acto procesal sorpresivo (véase la sección~\ref{sec:nulidad-cadena}).
\end{importante}

\section{Cómputo de plazos}

El día de salida de la resolución nunca corre a efectos del cómputo. Si se dicta un martes, no queda notificada ese día: queda notificada el próximo día de nota disponible (el viernes).

\begin{ejemplo}{el cómputo del plazo}
Una resolución dictada el martes 18 de agosto no se notifica ese día. El siguiente día de nota es el viernes 21: la notificación se produce ese viernes y el plazo para responder comienza a correr el primer día hábil siguiente, es decir, el lunes posterior; nunca desde el día en que se dictó ni desde el de la notificación.
\end{ejemplo}

\section{La preclusión}\label{sec:preclusion}

\begin{definicion}{Preclusión}
Principio por el cual, vencida una etapa o el plazo de un acto procesal, el acto ya no puede realizarse; lo que ha precluido no puede reabrirse salvo excepciones expresas de la ley. Es el cierre de una etapa: cada acto procesal es un eslabón que debe cerrarse para poder unirse al siguiente.
\end{definicion}

Un error en el cómputo puede hacer que opere la preclusión. Cada acto se cierra de dos maneras: por \textbf{cumplimiento} (la parte realiza el acto dentro del plazo) o por \textbf{vencimiento} (deja pasar el plazo). En ambos casos, cerrado el acto, no puede volver a ejercerse.

\begin{ejemplo}{el caso trampa de la preclusión}
Quien cuenta con quince días para contestar la demanda y la contesta el día tres ya no puede completarla ni modificarla en los días restantes, aunque el plazo no haya vencido. Si omitió ofrecer prueba, no podrá subsanarlo: al ejercer el acto anticipadamente, perdió los días restantes.
\end{ejemplo}

\begin{importante}
La \textbf{ampliación de la demanda} es el único acto que no precluye por su sola realización: puede ampliarse después de presentada, pero únicamente \textbf{antes de notificarla} al demandado, porque hasta entonces la relación procesal no involucra a la contraparte. Una vez notificada, se traba la litis y precluye la posibilidad de ampliarla.
\end{importante}

\begin{cornell}{Notificaciones, domicilios y preclusión}
\cq{¿Qué formas de notificación hay?}{Personal, por cédula y por nota (martes y viernes), salvo resoluciones del art. 133 CPCCN que requieren cédula o notificación personal.}
\cq{¿Cómo se cuenta el plazo?}{No corre el día de salida de la resolución; se notifica en el próximo día de nota y el plazo corre desde el primer día hábil siguiente.}
\cq{¿Qué domicilio importa en una sociedad?}{El domicilio real registrado (p. ej., ante la IGJ); el fiscal no tiene, en principio, efectos de notificación.}
\cq{¿Cómo se cierra un acto procesal?}{Por cumplimiento o por vencimiento del plazo; cerrado, no se reabre.}
\cq{¿Cuál es la excepción a la preclusión?}{La ampliación de la demanda, posible solo antes de su notificación al demandado.}
\cresumen{Toda resolución debe notificarse por alguna de las tres formas legales y los plazos se cuentan desde el día hábil siguiente a la notificación. La preclusión cierra cada acto por cumplimiento o vencimiento; solo la demanda puede ampliarse, y únicamente antes de notificada.}
\end{cornell}

\chapter{El expediente electrónico, las audiencias virtuales y los costos de litigar}

La práctica actual se desarrolla en gran medida por medios electrónicos. Este capítulo reúne las cuestiones prácticas de acceso al expediente, desarrollo de audiencias a distancia y costos asociados a litigar.

\section{El expediente electrónico}

El acceso a una causa para su consulta o seguimiento se hace a través del sistema de consulta pública de expedientes. En un expediente civil se selecciona la Cámara correspondiente (en el caso analizado, la Cámara Nacional Civil) y se indica el año de inicio; el ejemplo trabajado fue el expediente 6923/25. Una audiencia no es un acto aislado ni espontáneo: es consecuencia de actos previos (demanda, contestación) que desembocan en la fijación de fecha por el juzgado. La fecha se notifica por \textbf{cédula electrónica} al domicilio procesal de cada parte, y la audiencia preliminar se conoce en la práctica del fuero como «audiencia de 360» (capítulo~\ref{cap:audiencia}).

\section{La audiencia virtual}

Las audiencias virtuales son, en principio, públicas, pero el ingreso a la sala requiere un enlace que otorga el juzgado. Dado que la audiencia puede ser revisada por la Cámara, la imagen y el sonido deben permitir visualizar correctamente a cada interviniente, en un entorno adecuado (por ejemplo, sentado frente a un escritorio).

El lugar y la forma del acto judicial no son un detalle: los espacios cumplen una función simbólica además de práctica, y cuando el acto se vuelve excesivamente doméstico se debilita la solemnidad que históricamente caracterizó a la administración de justicia.

\begin{ejemplo}{la \emph{traditio} romana}
En el derecho romano, ante la falta de alfabetización y de instrumentos escritos, se recurría a ceremonias simbólicas para que un acto trascendente fuera recordado. En la compraventa de un terreno, el vendedor entregaba al comprador un pedacito de tierra (la \emph{traditio}) y luego se la quitaba simbólicamente, ante dos o tres testigos, para generar un recuerdo duradero. Del mismo modo, cuando el acto judicial pierde su marco formal, se debilita su capacidad de ser recordado, comprendido y respetado por quienes participan.
\end{ejemplo}

\section{Tasa de justicia, bono y cajas previsionales}

Todo proceso exige el pago de la \textbf{tasa de justicia} y del \textbf{bono de derecho fijo}, cuya forma de pago varía según la jurisdicción. En la Ciudad Autónoma el bono se abona en una plataforma del Banco Ciudad: se elige el pago del bono, la Cámara o fuero, el número de expediente y el CUIT del letrado, y se paga por transferencia, con acreditación directa en el expediente sin comprobante adicional. En la provincia de Buenos Aires el pago se gestiona a través del Colegio de Abogados (CRA): el comprobante se descarga en PDF y se acompaña con la documentación de la demanda.

En la provincia de Buenos Aires y en la mayoría de las provincias existe una \textbf{caja de previsión} específica para abogados, de afiliación obligatoria: entre el 10\% y el 15\% de los honorarios se destina a ella, con aportes a cargo del profesional y, en parte, de la contraparte. Su incumplimiento impide obtener regulaciones de honorarios o testimonios, entre otras consecuencias.

\begin{importante}
La obligación no es uniforme. En la Capital Federal los abogados pueden inscribirse como autónomos ante la AFIP y no hay caja especial más allá del bono. En la provincia de Buenos Aires, Tierra del Fuego y otras (como Córdoba) con caja independiente, el aporte es obligatorio; si al cierre del año los depósitos de honorarios no alcanzan el mínimo, el profesional debe aportar la diferencia para mantener la matrícula y la afiliación. El costo de litigar varía según la jurisdicción y debe calcularse antes de iniciar cualquier proceso.
\end{importante}

\begin{cornell}{Práctica electrónica y costos}
\cq{¿Cómo se notifica la audiencia?}{Por cédula electrónica al domicilio procesal; la preliminar se llama coloquialmente «audiencia de 360».}
\cq{¿Quién puede entrar a una audiencia virtual pública?}{Quien cuente con el enlace otorgado por el juzgado.}
\cq{¿Por qué importa la formalidad del acto?}{Porque el marco formal da solemnidad y capacidad de ser recordado y respetado (ejemplo de la \emph{traditio}).}
\cq{¿Cómo difieren los aportes previsionales?}{Obligatorios con caja propia en PBA y otras provincias; en Capital, autónomos ante AFIP.}
\cresumen{El litigio hoy se gestiona por expediente electrónico, cédulas electrónicas y audiencias virtuales, que exigen calidad técnica y formalidad. Tasa de justicia, bono y aportes previsionales varían según la jurisdicción y deben calcularse antes de litigar.}
\end{cornell}

\part{La demanda y la defensa}

\chapter{La construcción de la demanda: cliente, relato y estrategia}\label{cap:construccion}

Un conflicto entre dos personas llega a convertirse en proceso judicial a través de la demanda. Este capítulo explica la lógica que la sustenta: cómo se recibe el relato del cliente, qué hechos se seleccionan y cómo se piensa estratégicamente frente a una contraparte que también juega.

\section{La demanda como pretensión y como estrategia}

Reducida al mínimo, la demanda es un \textbf{reclamo}: una pretensión apoyada en tres pilares. Constituye además la estrategia procesal del actor.

\begin{definicion}{Pretensión (demanda)}
Reclamo o pedido que se basa en tres pilares: (1) los \textbf{hechos}, como narración; (2) las \textbf{pruebas}; y (3) el \textbf{derecho}. En cuanto acto, la demanda fija los hechos controvertidos mediante un relato acompañado de las pruebas que lo sostienen o que se ofrecen para producirse durante el proceso.
\end{definicion}

La demanda integra el relato aportado por el cliente, la prueba que lo respalda y el aporte propio del abogado: la \textbf{fundamentación jurídica}, el encuadre de los hechos en el derecho aplicable. El objetivo es que el juez tenga por establecidos los hechos tal como los presenta el cliente: convencerlo es el camino hacia una sentencia favorable. El actor supone que los testigos declararán en el sentido esperado y que los peritos se expedirán medianamente a su favor, pero debe recordar que la contraparte desplegará su propia estrategia.

\section{La entrevista con el cliente}

El conflicto llega al abogado a través del cliente, porque el profesional no presenció los hechos. De allí la importancia de investigar y formular la mayor cantidad posible de preguntas. Una vez reconstruido el hecho, corresponde evaluar el derecho aplicable y decidir si se recomienda iniciar el juicio: es obligación del abogado comunicar al cliente una estimación de las probabilidades de éxito.

Todo relato es necesariamente incompleto.

\begin{ejemplo}{los 200 millones de sensaciones}
Una persona cuenta cómo llegó a la facultad: se preparó, tomó un colectivo unos 40 minutos y combinó dos líneas de subte. Durante ese trayecto percibió, consciente o no, enormes cantidades de estímulos (gente que pasó, un vendedor ambulante, alguien que la empujó) que ni siquiera registró. Por mejor que sea el esfuerzo de síntesis, el relato es incompleto.
\end{ejemplo}

Lo mismo ocurre con el relato de un cliente, que además relata desde su propio punto de vista y con su carga emocional. El trabajo del abogado no es transcribirlo sino \textbf{reconstruirlo con criterio profesional}: elaborar una estrategia y organizar la narración en función de ella. La entrevista es instancia central: requiere preguntas bien diseñadas y capacidad de interpretar relatos de personas de muy distintos niveles socioculturales.

Las personas se comunican fundamentalmente mediante relatos: el pensamiento necesita estructura narrativa (una historia que empieza, se desarrolla y termina), porque somos seres temporales que no perciben la realidad de modo instantáneo. Por eso la demanda se construye como un relato.

\section{Selección de los hechos}

Tres criterios determinan si un hecho se incluye en el relato estratégico:

\begin{enumerate}
\item su \textbf{relevancia jurídica}: si favorece o no la posición del cliente;
\item la \textbf{calificación jurídica} adecuada del hecho;
\item la \textbf{posibilidad de probarlo}, criterio decisivo.
\end{enumerate}

\begin{importante}
Todo aquello que se afirma y no se puede probar juega en contra. Si un hecho es relevante pero no hay elementos para acreditarlo, lo estratégicamente correcto es omitirlo o dejarlo en segundo plano. El abogado debe anticipar la posición de la contraparte para no dejar flancos débiles en el propio relato.
\end{importante}

\section{Estrategia personal y estrategia procesal}

\begin{table}[H]
\centering
\caption{Estrategia personal frente a estrategia procesal}
\begin{tabularx}{\textwidth}{>{\raggedright\arraybackslash\bfseries}p{0.30\textwidth} X X}
\toprule
 & Estrategia personal (gimnasio, régimen) & Estrategia competitiva o procesal \\
\midrule
¿Con quién se enfrenta? & Consigo mismo; no hay rival que juegue. & Con otro que también juega, y a veces muy bien. \\
¿Qué debe considerar? & Lo que uno quiere y planifica. & Lo que uno quiere y lo que se supone que el otro puede sostener y afectará. \\
¿Qué flexibilidad exige? & Baja. & Alta: lo que ocurra no será necesariamente lo que uno quiere. \\
\bottomrule
\end{tabularx}
\end{table}

Debe sostenerse la estructura estratégica durante todo el proceso; la ausencia de estrategia equivale a litigar sin rumbo.

\begin{cornell}{Construcción de la demanda}
\cq{¿Qué es la demanda?}{Un reclamo basado en hechos (narración), pruebas y derecho; es la estrategia del actor.}
\cq{¿Por qué el abogado no transcribe el relato del cliente?}{Porque es incompleto y subjetivo; debe reconstruirlo con criterio y organizarlo según una estrategia.}
\cq{¿Qué criterios deciden qué hecho incluir?}{Relevancia jurídica, calificación jurídica y posibilidad de probarlo (el decisivo).}
\cq{¿Qué lo diferencia de una estrategia personal?}{Hay un oponente activo, por lo que se exige tener en cuenta lo que el otro puede sostener y gran flexibilidad.}
\cresumen{La demanda es relato, prueba y derecho. El abogado reconstruye con criterio el relato del cliente, selecciona hechos relevantes, bien calificados y probables, evalúa e informa las chances del caso y piensa estratégicamente de principio a fin, anticipando a la contraparte.}
\end{cornell}

\chapter{Estructura y contenido de la demanda}\label{cap:estructura-demanda}

El Código fija el contenido mínimo de la demanda. Este capítulo estudia ese contenido, el objeto, los sujetos legitimados y demandados, la citación en garantía del asegurador y la narración de los hechos.

\section{Requisitos formales: el artículo 330}

\begin{regla}{art. 330 CPCCN --- contenido de la demanda}
La demanda será deducida por escrito y contendrá:
\begin{enumerate}[label=\arabic*.\textsuperscript{o})]
\item El nombre y domicilio del demandante.
\item El nombre y domicilio del demandado.
\item La cosa demandada, designándola con toda exactitud.
\item Los hechos en que se funde, explicados claramente.
\item El derecho expuesto sucintamente, evitando repeticiones innecesarias.
\item La petición en términos claros y positivos. Deberá precisarse el monto reclamado, salvo cuando al actor no le fuera posible determinarlo por las circunstancias del caso o porque la estimación dependiera de elementos aún no definitivamente fijados y la promoción de la demanda fuese imprescindible para evitar la prescripción.
\end{enumerate}
\end{regla}

La demanda comienza por el sumario y el encabezamiento, que admiten fórmulas equivalentes («promueve demanda», «inicia demanda», «demanda»); el único dato que varía entre casos es el domicilio real del actor. Aplicado a un modelo de demanda por daños, se verificó cada inciso: el actor aparece como cónyuge supérstite (inc.~1); se identifica a patrón, marinero, empresa fluvial y aseguradora (inc.~2); se reclama por daños y perjuicios con detalle (inc.~3); se relata el accidente (inc.~4); el derecho debía desarrollarse con las normas aplicables (inc.~5); y el monto puede diferirse si hay imposibilidad fundada (inc.~6).
% TODO: verificar consistencia entre fuentes originales (el modelo presenta a la actora como «la señora Pérez»; otros apuntes nombran distintas personas para el caso del río Tigre).

\section{El objeto de la demanda}

El \textbf{objeto} es lo que se pide; tiene importancia trascendente porque fija el tema decidendum desde la perspectiva del actor (sección~\ref{sec:tema}): el juez no podrá conceder más de lo pedido. Debe redactarse con claridad, precisión y concisión.

\begin{regla}{fórmula del objeto}
«[Nombre del actor] viene a promover demanda de daños y perjuicios contra [demandado/s], por la suma de [monto], en razón de los hechos y el derecho que se expondrán, como consecuencia del accidente relatado.»
\end{regla}

\begin{ejemplo}{objeto en el caso guía}
«Vengo a promover demanda por daños y perjuicios contra los señores patrón y marinero, y también contra la empresa de colectivos, por ser responsables de ocasionar lesiones graves que causaron la muerte de quien fuera mi mujer y compañía de trabajo, la señora [nombre], y cito en garantía a la compañía de seguros en virtud del artículo 118 de la Ley de Seguros.»
\end{ejemplo}
% TODO: el apunte original no consigna el nombre de la víctima, reemplazado por el expositor con «[nombre]».

\begin{importante}
No es indispensable indicar el monto en el objeto; puede precisarse en otro momento. Pero el monto, o al menos una pauta para estimarlo, no puede faltar en el escrito: su ausencia total habilitaría a la contraparte a oponer una excepción de defecto legal. Precisar qué se reclama permite además que el demandado evalúe su defensa, incluso allanarse.
\end{importante}

\section{Legitimación activa y pasiva}

\begin{definicion}{Legitimación}
Condición de ser titular del derecho que se reclama (\textbf{legitimación activa}: quién puede reclamar; requiere titularidad del derecho, perjuicio actual e interés legítimo) o de ser la persona a quien puede exigírsele su cumplimiento (\textbf{legitimación pasiva}: a quién puede dirigirse el reclamo). Si se demanda a quien no puede responder por la pretensión, la acción es rechazable por falta de legitimación pasiva.
\end{definicion}

La legitimación activa corresponde, en el caso guía, al cónyuge supérstite. Plantearla expresamente desde el inicio es una estrategia que anticipa y neutraliza una eventual excepción previa de falta de legitimación (capítulo~\ref{cap:defensa}).

Todos los demandados responden por el mismo hecho, aunque no por el mismo título:

\begin{table}[H]
\centering
\caption{Títulos de responsabilidad en el caso guía}
\begin{tabularx}{\textwidth}{>{\raggedright\arraybackslash\bfseries}p{0.28\textwidth} X}
\toprule
Demandado & Título de responsabilidad \\
\midrule
Marinero & Responsable directo por negligencia: no estaba en la popa verificando las condiciones para la maniobra. \\
Patrón & Responsable directo por negligencia: dio la orden o permitió la maniobra en condiciones inseguras. \\
Empresa de transporte fluvial & Responsabilidad solidaria como dueña de la cosa riesgosa (la lancha; cosa riesgosa es un objeto como un automóvil o una embarcación) y empleadora de los responsables directos. \\
Aseguradora & Legitimada pasiva por el contrato de seguro con la empresa (art. 118 Ley de Seguros). \\
\bottomrule
\end{tabularx}
\end{table}

% TODO: verificar consistencia entre fuentes originales (otro apunte caracteriza la responsabilidad del patrón como subjetiva y la de la empresa como objetiva, y la menciona como «empresa de catamaranes»).

\section{La elección de los demandados}

En el caso guía los demandados identificados son los dueños de la lancha, la empresa naviera y la aseguradora; se descarta demandar a la municipalidad y se analiza sumar al Estado (por ejemplo, la Prefectura Naval) como codemandado. La cantidad no es indiferente: cada demandado implica otra cédula, más carga procesal y más tiempo. Conviene evaluar la solvencia de cada uno, la carga que implica cada codemandado, los plazos aplicables (los que rigen contra el Estado son más largos) y la probabilidad de que termine siendo responsable y efectivo pagador. Puede convenir incluir a organismos estatales aun con plazos más largos. Es una decisión estratégica en cada caso.

\section{La citación en garantía del asegurador}

\begin{definicion}{Citación en garantía}
Acto procesal por el cual se incorpora al asegurador como parte del proceso, para que responda por la condena en la medida del seguro contratado. El asegurador no es responsable directo del hecho dañoso, pero si los demandados son condenados y no pagan, el seguro es la garantía de pago.
\end{definicion}

La obligación del asegurador nace del contrato de seguro: algunas actividades imponen legalmente contar con seguro (el transporte público de pasajeros) y la circulación de vehículos tampoco puede hacerse sin cobertura. El contrato cubre un riesgo determinado; cuando se concreta, hay siniestro.

\begin{importante}
Para que la condena alcance al asegurador, este debe haber sido parte. Si no fue citado, la sentencia no podrá alcanzarlo. Por eso es determinante la primera resolución sobre la demanda, «téngase por presentado y por parte», que fija a quiénes alcanzará la sentencia. Así como los demandados, de saber que tienen cobertura, denunciarían el hecho a su aseguradora para que se defienda, el actor tiene interés en que el asegurador sea citado desde el inicio.
\end{importante}

\begin{regla}{art. 118, Ley de Seguros n.\textsuperscript{o} 17.418}
\begin{itemize}
\item El damnificado puede citar en garantía al asegurador hasta que se reciba la causa a prueba.
\item Debe interponer la demanda ante el juez del lugar del hecho o del domicilio del asegurador.
\item El asegurador citado puede oponer a la víctima las defensas nacidas con anterioridad al siniestro.
\item Si no comparece, es representado por el asegurado.
\item La sentencia dictada contra el asegurado alcanza al asegurador en la medida del seguro.
\end{itemize}
\end{regla}

\section{La narración de los hechos}

La demanda debe narrar dos categorías de hechos: los \textbf{hechos del accidente} (mecánica del siniestro, desarrollo y sujetos) y los \textbf{hechos que fundan los rubros indemnizatorios} (profesión, ingresos, familia y actividades de la víctima, que justifican cada rubro). Existencia de hijos, tiempo de convivencia, sufrimiento, deportes o actividades también son hechos; ubicarlos en el capítulo del accidente o distribuirlos en cada rubro es una decisión estratégica.

\begin{ejemplo}{relato modelo}
«El hecho que dio origen a la presente demanda ocurrió el 9 de [mes] de 2020, a las 11:30 horas aproximadamente, mientras las señoras Lorena Montes y Carla González circulaban a bordo de una pequeña embarcación por el río Tigre, perteneciente al Club Delta, cuando, cerca de la estación fluvial de Tigre, fueron embestidas por la lancha comercial de pasajeros de la línea Tigre Argentina S.R.L. La lancha colectiva realizaba maniobras de marcha atrás para salir de la terminal y, al retroceder, succionó y pasó por encima de la pequeña embarcación. Como consecuencia, la señora Lorena Montes, luego de permanecer hospitalizada cinco días en un centro de salud de Pacheco, falleció; la señora Carla González sufrió lesiones leves. Los pormenores de la mecánica surgirán de las pruebas a producirse, en especial de la pericia accidentológica. Lo sucedido dio origen a la causa penal «N.N. s/ homicidio culposo», ante el juzgado de instrucción en lo criminal y correccional.»
\end{ejemplo}
% TODO: verificar consistencia entre fuentes originales (fecha y datos de personas, días de internación y composición de la tripulación difieren entre apuntes; las propias notas califican el relato de incompleto).

Otro relato, reproducido de las actuaciones del expediente penal correccional de San Isidro, está en primera persona:

\begin{ejemplo}{relato según el expediente penal}
«Navegábamos en un kayak por el río Tigre a la altura de la estación fluvial, habiendo partido de las instalaciones del club. La lancha, amarrada como último de varias embarcaciones, inició una maniobra de retroceso para zarpar, invirtiendo nuestra embarcación. El kayak dio vuelta campana y la señora fue succionada por la turbulencia de las hélices, con lesiones gravísimas que provocaron su muerte días después. Según testigos, recogido por la Prefectura Naval Argentina, y un video de celular, la maniobra violó el procedimiento usual: el marinero charlaba con el patrón en lugar de estar en la popa verificando condiciones e indicando la zarpada más segura. Contábamos con todos los elementos reglamentarios de seguridad y navegábamos en la zona habilitada para embarcaciones a remo, observando las reglas de tránsito fluvial, lo que acreditan los testigos y el personal del club que autorizó la salida.»
\end{ejemplo}
% TODO: el apunte original no consigna la fecha ni la hora exactas ni los nombres, reemplazados por marcadores.

El último párrafo tiene especial importancia: anticipa la defensa previsible de los demandados (que el hecho no les fue imputable y que las víctimas se colocaron en situación de riesgo) y acredita que estaban en zona habilitada, eran socios con habilitación y cumplían las normas. Se señala además que el puerto fluvial es de tránsito compartido por embarcaciones de distintos clubes y que, en los hechos, no hay un representante de la Prefectura que ordene el tránsito.

\section{Reservas, representación y otras cuestiones de la demanda}

La demanda debe contener el relato de los hechos, el derecho invocado, los daños, los fundamentos de la responsabilidad, el objeto y el ofrecimiento de prueba, con toda la documental disponible. Al notificarse, el demandado recibe la cédula con copia de la demanda y de la documental, salvo que esta sea muy voluminosa y se pida no trasladarla físicamente por estar en el expediente digital. La reserva del caso federal se estudia en la sección~\ref{sec:reserva-federal}.

\begin{nota}
Cuando un padre o madre representa a un hijo menor, la demanda debe consignarlo expresamente («Fulano de tal, en representación de su hijo menor\ldots»), y se acredita con la partida de nacimiento, que es una de las primeras piezas de prueba documental. Quedan enunciados, sin desarrollo, otros temas de la parte general: la intervención de terceros (arts.~90 y 91 CPCCN), el régimen de notificaciones (en particular, de la demanda) y las excepciones procesales.
\end{nota}

\begin{ejercicio}{trabajo práctico integrador}
A partir del caso guía del accidente en el río Tigre: (1) plantear la competencia y fundamentar el fuero; (2) definir a quiénes se demanda; (3) redactar la narración completa de los hechos, clara, ordenada y cronológica; (4) elaborar el objeto conforme a la fórmula estudiada.
\end{ejercicio}

\begin{cornell}{Estructura de la demanda}
\cq{¿Qué exige el art. 330 CPCCN?}{Escrito con nombre y domicilio de actor y demandado, cosa demandada, hechos, derecho sucinto y petición clara, con monto salvo imposibilidad fundada.}
\cq{¿Para qué sirve el objeto?}{Fija lo que se pide y limita lo que el juez puede conceder; sin monto ni pauta, procede excepción de defecto legal.}
\cq{¿Qué es la legitimación activa y la pasiva?}{Titularidad del derecho reclamado y condición de persona a quien puede exigirse; se explicita desde el inicio para neutralizar excepciones.}
\cq{¿Por qué citar al asegurador desde el inicio?}{Porque solo la sentencia dictada con él como parte lo alcanza, en la medida del seguro (art. 118 Ley de Seguros).}
\cq{¿Qué hechos narra la demanda?}{Los del accidente y los que fundan cada rubro indemnizatorio.}
\cresumen{La demanda sigue el art. 330 CPCCN: identifica partes, objeto, hechos, derecho y petición. El objeto limita la decisión del juez; la legitimación y los títulos de responsabilidad de cada demandado se explicitan; el asegurador debe citarse en garantía para ser alcanzado por la condena, y los hechos se narran distinguiendo el accidente de lo que funda cada rubro.}
\end{cornell}

\chapter{Los daños y los rubros indemnizatorios}\label{cap:danos}

Los rubros indemnizatorios son el eje económico de todo juicio civil. Identificarlos y acreditarlos con la prueba adecuada define en gran medida el resultado económico para el cliente.

\section{El daño como cambio}

\begin{definicion}{Daño}
Afectación de un derecho: modificación perjudicial de un estado de cosas, producida por un hecho con el que debe existir relación de causalidad. Consiste en la diferencia entre la situación de la persona \textbf{antes} del hecho y \textbf{después} de él; esa brecha (física, psicológica, patrimonial o social) es lo que debe acreditarse con precisión.
\end{definicion}

\section{Clasificación de los daños}

\begin{table}[H]
\centering
\caption{Daño patrimonial y extrapatrimonial}
\begin{tabularx}{\textwidth}{X X}
\toprule
\textbf{Daño patrimonial} & \textbf{Daño extrapatrimonial} \\
\midrule
Daño emergente (por ejemplo, el kayak destruido). & Daño moral: angustia, pena, tristeza permanente. \\
Lucro cesante: ingresos que la víctima dejó de percibir. & Daño psicológico: perturbación de la psiquis, tratable y cuantificable. \\
Pérdida de chance: expectativas futuras no concretadas. & Daño estético (si correspondiera). \\
Gastos médicos, de internación, velatorio y sepelio. & \\
\bottomrule
\end{tabularx}
\end{table}

\section{Daño moral y daño psicológico}

\begin{definicion}{Daño moral}
Daño inmaterial, espiritual y afectivo: una tristeza, angustia o pena que no desaparece. En algunos supuestos requiere mayor comprobación; en la pérdida de un ser querido no, pues se entiende que toda persona sufre la pérdida de un familiar. Como la reparación integral (volver las cosas al estado anterior) es imposible, el derecho encuentra en el dinero el único medio de reparación.
\end{definicion}

\begin{definicion}{Daño psicológico}
Perturbación producida en la psiquis, que es tratable. Un perito puede establecer que requiere una terapia de cierta frecuencia semanal durante un período, y a partir del valor de las sesiones se obtiene el monto. El daño moral no admite esa reparación específica. Puede manifestarse como depresión, estrés postraumático u otras afectaciones, según la evolución de cada persona.
\end{definicion}

\begin{table}[H]
\centering
\caption{Daño moral frente a daño psicológico}
\begin{tabularx}{\textwidth}{>{\raggedright\arraybackslash\bfseries}p{0.2\textwidth} X X}
\toprule
 & Daño moral & Daño psicológico \\
\midrule
Naturaleza & Inmaterial, espiritual, afectivo. & Perturbación de la psiquis. \\
Prueba & En la muerte de un familiar se presume; no requiere pericia específica. & Requiere pericia: el perito dictamina tratamiento y duración. \\
Cuantificación & Difícil de precisar con exactitud. & Sesiones $\times$ valor $\times$ tiempo de tratamiento. \\
\bottomrule
\end{tabularx}
\end{table}

\section{Lucro cesante y pérdida de chance}

El \textbf{lucro cesante} es lo que la víctima deja de percibir de manera concreta y mensurable (el ingreso mensual efectivo). La \textbf{pérdida de chance} es lo que eventualmente podría haber ganado, menos concreto pero estimable.

\begin{ejemplo}{estimación según la actividad de la víctima}
Si trabajaba en relación de dependencia, la expectativa se proyecta hasta la edad jubilatoria; si desarrollaba una labor profesional en crecimiento, debe guardar relación con ese crecimiento; si tenía una actividad comercial en auge o era deportista, las circunstancias son distintas. En todos los casos deben enunciarse y probarse.
\end{ejemplo}

\section{El componente emotivo en la prueba del daño moral}\label{sec:moral-estrategia}

La determinación del daño moral tiene además una dimensión estratégica: existe un componente emocional, ideológico y de prejuicios sociales con el que el profesional trabaja. Mostrar la magnitud real de la pérdida para una persona determinada, con sus circunstancias particulares, incide en la valuación.

\begin{ejemplo}{deporte amateur y movilidad}
Si el actor practicaba un deporte de modo amateur y perdió la movilidad de un miembro, es pertinente ofrecer testigos, acreditar su condición de socio del club desde hace años y exhibir constancias de participación, para mostrar que la pérdida representa para esa persona un daño moral mayor que para una persona sedentaria.
\end{ejemplo}

\begin{cornell}{Rubros indemnizatorios}
\cq{¿Qué es el daño?}{Diferencia entre la situación antes y después del hecho, con relación de causalidad.}
\cq{¿Cómo se clasifica?}{Patrimonial (emergente, lucro cesante, pérdida de chance, gastos) y extrapatrimonial (moral, psicológico, estético).}
\cq{¿Moral o psicológico?}{El moral es inmaterial y, en la pérdida de un familiar, se presume; el psicológico es una perturbación tratable que exige pericia y se cuantifica por el tratamiento.}
\cq{¿Lucro cesante o pérdida de chance?}{Lo concreto que se dejó de percibir frente a lo que eventualmente se habría ganado.}
\cresumen{El daño es el cambio entre el antes y el después. Los rubros se dividen en patrimoniales y extrapatrimoniales; deben enunciarse y probarse según la actividad y circunstancias de la víctima, distinguiendo daño moral (presumible en la muerte de un familiar) de daño psicológico (pericial y cuantificable).}
\end{cornell}

\chapter{El inicio del proceso y la traba de la litis}\label{cap:inicio}

La demanda abre el proceso y produce efectos inmediatos. Este capítulo sigue el camino desde la mediación previa hasta que, con demanda y contestación, queda trabada la litis y fijado el tema decidendum.

\section{La mediación prejudicial}

La mediación prejudicial es, cuando se exige, un requisito previo para ciertos juicios (la mayoría de las provincias la exigen con carácter obligatorio). Pero es \textbf{previa} al proceso, no parte de él.

\begin{importante}
El proceso judicial en sentido técnico se inicia con la interposición de la demanda. La mediación, aunque sea etapa ineludible en muchas jurisdicciones, es anterior al proceso.
\end{importante}

\section{La demanda como acto procesal y sus efectos}

\begin{definicion}{Acto procesal}
Todo acto realizado dentro del proceso. Tienen por objeto ir de la demanda a la sentencia y pueden provenir de las partes, del juez (resoluciones) o de terceros auxiliares (peritos, oficiales de justicia).
\end{definicion}

La interposición de la demanda produce efectos fundamentales:

\begin{itemize}
\item la fijación de los \textbf{hechos controvertidos};
\item la fijación de la \textbf{competencia} (sección~\ref{sec:fija-competencia});
\item la \textbf{interrupción de la prescripción} del derecho reclamado, lo que la conecta con la caducidad de instancia (sección~\ref{sec:caducidad-prescripcion}).
\end{itemize}

\section{Traslado y contestación}

Presentada la demanda, el juez \textbf{da traslado} mediante \textbf{cédula} (en el caso analizado, en papel). En el juicio ordinario el traslado es por \textbf{quince días}, dentro de los cuales el demandado \textbf{contesta la demanda} (capítulo~\ref{cap:defensa}). Si hay reconvención, se corre traslado de ella al actor y recién entonces queda trabada la litis.

\section{El tema decidendum}\label{sec:tema}

\begin{definicion}{Tema decidendum}
Límite dentro del cual puede resolver el juez, fijado por la pretensión del actor, la defensa del demandado y la prueba producida. Es el conjunto de cuestiones sometidas a decisión a partir de la demanda y la contestación: el juez no puede dar más de lo pedido ni reconocer mayores defensas que las introducidas (principio de congruencia).
\end{definicion}

Para las partes también es una limitación: lo que no se plantea en su oportunidad no puede introducirse después, por preclusión y por el derecho de defensa. Por eso importa estudiar el planteo de los abogados y no solo la sentencia, que sería «el final de la película sin haber visto el desarrollo».

\section{El proceso civil y el proceso penal}\label{sec:civil-penal}

No es necesario aguardar la resolución del proceso penal para demandar civilmente, pues el daño se produjo con independencia de aquel. Lo que aporta la resolución penal es que los hechos allí acreditados quedan, de hecho, prácticamente probados para el proceso civil; si el hecho no se califica como delito, incide en otro sentido, pero el daño no cambia. La condena penal de los demandados favorece al actor civil; en el caso del río Tigre, patrón y marinero fueron condenados en sede penal, no así la empresa.

\begin{cornell}{Inicio y traba de la litis}
\cq{¿Cuándo empieza el proceso?}{Con la demanda; la mediación prejudicial es previa y no forma parte de él.}
\cq{¿Qué efectos produce la demanda?}{Fija hechos controvertidos y competencia, e interrumpe la prescripción.}
\cq{¿Cómo se traslada la demanda y en qué plazo se contesta?}{Por cédula; en el juicio ordinario, quince días.}
\cq{¿Qué es el tema decidendum?}{El límite de la decisión del juez, fijado por demanda, contestación y prueba; consagra la congruencia.}
\cq{¿Hay que esperar la sentencia penal?}{No; pero los hechos acreditados en lo penal quedan prácticamente probados en lo civil.}
\cresumen{El proceso comienza con la demanda, que fija competencia y hechos controvertidos e interrumpe la prescripción. Se traslada por cédula y se contesta en quince días; con demanda y contestación queda trabada la litis y fijado el tema decidendum, que limita al juez y a las partes.}
\end{cornell}

\chapter{La defensa del demandado}\label{cap:defensa}

El demandado no se limita a negar: \emph{defiende}. Este capítulo sistematiza cómo responde a una demanda, desde las excepciones previas hasta la reconvención.

\section{Un orden de trabajo}

Se propone un mapa mental: primero se revisa si hay defectos de forma que permitan oponer excepciones; luego, sobre el fondo, se decide si se niegan los hechos o se aceptan y se discute su consecuencia jurídica. El orden de trabajo para cualquier demanda es: \textbf{forma $\rightarrow$ hechos $\rightarrow$ derecho $\rightarrow$ prueba}.

\begin{ejemplo}{un reclamo de dinero entre amigos}
Un amigo reclama una deuda que, según él, nunca fue devuelta: es la \textbf{demanda} (hechos: se prestó una suma en una fecha; pruebas: mensajes, un comprobante; regla: los préstamos se devuelven). Antes de discutir si se debe, podría plantearse una cuestión previa («esto no se arregla ante este árbitro», «ya lo discutimos ante otro», «vos no sos quien me prestó»): las \textbf{excepciones}. Si se sigue adelante, o se niegan los hechos («nunca me prestaste nada») o se aceptan y se discute su consecuencia («me lo prestaste, pero fue un regalo»): la \textbf{defensa de fondo}. Se necesita un escrito ordenado, que reconozca o desconozca cada papel y niegue cada afirmación; y si fue el amigo quien quedó debiendo por la misma operación, se lo plantea como contrarreclamo: la \textbf{reconvención}.
\end{ejemplo}

\begin{table}[H]
\centering
\caption{Defensas de forma y defensas de fondo}
\begin{tabularx}{\textwidth}{>{\raggedright\arraybackslash\bfseries}p{0.17\textwidth} X X}
\toprule
 & Defensas de forma (excepciones) & Defensas de fondo \\
\midrule
Qué discuten & Una cuestión técnica: el proceso no está acabado o correctamente planteado. & El centro del reclamo: hechos o derecho. \\
Qué se dice al juez & «Hasta que no se resuelva esta cuestión técnica no podemos seguir.» & «Los hechos no fueron esos» o «son esos, pero tienen otra consecuencia jurídica.» \\
Si prosperan & El planteo del actor puede frustrarse o trasladarse a otro lugar. & Se rechaza total o parcialmente la pretensión. \\
Regla práctica & Se plantean primero, pero nunca se deja de defender el fondo. & Se defienden siempre, subsidiariamente si la excepción no prospera. \\
\bottomrule
\end{tabularx}
\end{table}

\section{Las excepciones previas}

\begin{definicion}{Excepción previa}
Planteo por el cual se dice al juez que existe una cuestión técnica no del todo acabada o correctamente planteada y que, hasta resolverla, no se puede seguir. No afirma que el otro no tenga razón: señala un defecto técnico. Deben oponerse en la primera oportunidad procesal.
\end{definicion}

\begin{itemize}
\item \textbf{Incompetencia}: el juez ante quien se demanda no es el competente.
\item \textbf{Litispendencia}: ya existe otro juicio igual por esta cuestión.
\item \textbf{Falta de legitimación}: quien reclama no está legitimado, o el demandado no es la persona a quien debe reclamarse.
\item \textbf{Falta de personería}: no se acreditó poder suficiente.
\item \textbf{Arraigo}.
\item \textbf{Prescripción}: lo reclamado está prescripto.
\item \textbf{Defecto legal} en el planteo.
\item \textbf{Transacción}, como causal de extinción.
\end{itemize}

\begin{nota}
\textbf{Referencia complementaria, art.~347 CPCCN.}
% TODO: verificar referencia presente en las fuentes originales (texto abreviado; cotejar con el Código vigente).
«Sólo se admitirán como previas las siguientes excepciones: 1) Incompetencia. 2) Falta de personería en el demandante, en el demandado o en sus representantes, por carecer de capacidad civil para estar en juicio o de representación suficiente. 3) Falta de legitimación para obrar en el actor o en el demandado, cuando fuere manifiesta, [\ldots] 4) Litispendencia. 5) Defecto legal en el modo de proponer la demanda. 6) Cosa juzgada. 7) Transacción, conciliación y desistimiento del derecho. 8) Prescripción. 9) Arraigo [\ldots]». Las fuentes enumeran las defensas sin citar el artículo.
\end{nota}

Lo primero ante una demanda es verificar si hay un defecto de forma que permita excepcionar. Si el actor no supera la excepción, la cuestión puede terminar allí; si prospera (por ejemplo, la incompetencia y el actor debe litigar en otra jurisdicción), puede ser una complicación muy grande para él. La \textbf{prescripción} debe oponerse en la primera oportunidad: si el demandado contesta y omite plantearla, pierde la posibilidad de invocarla después.

\begin{ejemplo}{los trabajadores del buque y la empresa de Tierra del Fuego}
En juicios en trámite contra una empresa de Tierra del Fuego, los actores, casi todos correntinos que viajan a la temporada de pesca, trabajaban en un barco y demandaron en Capital Federal. Para ello, en un juicio laboral, el actor debe acreditar alguna de tres cuestiones (véase el cuadro siguiente). Los actores alegaron que la empresa tenía sede en Capital; la demandada acreditó estar inscripta desde hace años en la IGJ de Tierra del Fuego, y no se acreditó que el contrato se celebrara o cumpliera en Capital. Varios juicios tuvieron sentencia: se decretó la \textbf{incompetencia territorial}. Se planteó además una segunda incompetencia: el barco había salido del puerto y se sostuvo que, saliendo aunque sea por un instante, corresponde la justicia federal; fiscales y jueces concluyeron que hay incompetencia. Los actores y sus abogados deben ahora contactar abogados de la otra jurisdicción.
\end{ejemplo}

\begin{table}[H]
\centering
\caption{Criterios de competencia laboral del caso}
\begin{tabularx}{\textwidth}{>{\raggedright\arraybackslash\bfseries}p{0.34\textwidth} X}
\toprule
Criterio & Pregunta que debe responder el actor \\
\midrule
1. Lugar de trabajo & ¿Se trabajó en la Ciudad Autónoma de Buenos Aires? \\
2. Lugar de celebración del contrato & ¿El contrato fue firmado en CABA? \\
3. Domicilio del empleador & ¿El empleador tiene domicilio en CABA? \\
\bottomrule
\end{tabularx}
\end{table}

\begin{importante}
Las excepciones previas no sirven solo para «ganar tiempo»: permiten ordenar la estrategia y, en ciertos casos, trasladar a otra jurisdicción una serie de juicios (a la empresa del ejemplo le conviene litigar en Tierra del Fuego). Aun así, quien excepciona no debe dejar de defender el fondo, porque si la excepción no prospera tendrá que hacerlo.
\end{importante}
% TODO: un comentario final aparece trunco en las notas originales.

\section{Las defensas de fondo}

Al defenderse sobre el fondo hay dos posibilidades: \textbf{negar los hechos} (los hechos fueron otros y por lo tanto tienen otra consecuencia) o \textbf{aceptarlos y defenderse por el derecho} (se reconoce que pasó lo que dice el actor, pero se sostiene que corresponde otra norma o una consecuencia jurídica distinta).

\begin{ejemplo}{la empresa constructora en Puerto Madero}
Se es abogado de una empresa que levanta un edificio en Puerto Madero. El actor dice que debe cumplir cierta obligación del contrato. Se puede negar la existencia misma del contrato o de los hechos (con lo cual se niega también el derecho que se pretende fundar en ellos), o reconocer que el contrato existe y que los hechos ocurrieron, pero sostener que la cláusula debe interpretarse de otra manera.
\end{ejemplo}

\begin{importante}
\textbf{La gran estructura para contestar.} Primero se despeja si hay excepciones y se las plantea; nunca se deja de lado el fondo. Al defender el fondo se miran primero los hechos: si no son ciertos, se discute el hecho; si lo son, se discute su consecuencia jurídica. Cuando no hay discusión sobre los hechos, puede declararse que no son controvertidos porque se concuerda con ellos y plantear que es una cuestión de derecho en que se aplica otra norma.
\end{importante}

\section{El escrito de contestación}

El objeto es \textbf{contestar la demanda y ofrecer prueba}: ahora o nunca. El escrito es, en buena medida, un \textbf{espejo de la demanda}.

\begin{table}[H]
\centering
\caption{Orden del escrito de contestación}
\begin{tabularx}{\textwidth}{>{\centering\arraybackslash}p{0.1\textwidth} X}
\toprule
Orden & Contenido \\
\midrule
1 & Excepciones previas (si corresponden). \\
2 & Objeto: contestar la demanda y ofrecer prueba. \\
3 & Desconocimiento de los documentos presentados por el actor. \\
4 & Negativa puntual de cada hecho (y de cada rubro de la liquidación, si la hay). \\
5 & Relato de la realidad de los hechos: estrategia y argumentación de la defensa de fondo. \\
6 & Ofrecimiento de prueba. \\
7 & Reconvención, al final y por separado (si corresponde). \\
\bottomrule
\end{tabularx}
\end{table}

\subsection{Desconocer los documentos y negar los hechos}

De la prueba del actor, lo único que puede y debe hacerse en ese momento es desconocer lo \textbf{documental}; una pericia, un testigo o un oficio no pueden negarse ni desconocerse. El documento no desconocido se considera reconocido. Luego vienen las negativas.

\begin{nota}
\textbf{Referencia complementaria, art.~356, inc.~1 CPCCN.}
% TODO: verificar referencia presente en las fuentes originales (texto abreviado; cotejar con el Código vigente).
«En la contestación opondrá el demandado todas las excepciones o defensas [de fondo] que tuviere [\ldots] y deberá: 1) Reconocer o negar categóricamente cada uno de los hechos expuestos en la demanda, la autenticidad de los documentos acompañados que se le atribuyan y la recepción de las cartas y telegramas a él dirigidos, cuya copia se acompañe. Su silencio, sus respuestas evasivas o la negativa meramente general podrán estimarse como reconocimiento de la verdad de los hechos pertinentes y lícitos a que se refieran. En cuanto a los documentos, se los tendrá por reconocidos o recibidos, según el caso [\ldots]». Las fuentes remiten a este contenido sin citar el número.
\end{nota}

No conviene negar lo que no tiene sentido negar (hay quienes niegan hasta el nombre del actor); no se gana nada. La negativa se hace \textbf{por cada hecho}, de manera pormenorizada: se niega que en cierta fecha sucedió tal cosa, que el actor sufrió un accidente, que fue en tal lugar, que tuvo secuelas psicológicas, que fue pasajero del barco. Cada negativa debe ser \textbf{puntual, clara y precisa}, con una sola negación por negativa. Si hay liquidación, se niega cada partida. «Desconozco» y «niego» son equivalentes, aunque normalmente se usa «niego».

\begin{importante}
La negativa general («niego todo aquello que no haya sido expresamente reconocido») no tiene consecuencia alguna; puede incluirse, pero no sirve. Al negar puntualmente, cada hecho negado se vincula luego con los hechos controvertidos que se discutirán en la audiencia. Sobre los documentos, véase la sección~\ref{sec:documental}.
\end{importante}

\subsection{Pruebas comunes con la parte actora}

Con frecuencia las partes ofrecen las mismas pruebas.

\begin{ejemplo}{testigos ofrecidos también por el actor}
Para defenderse el demandado debe llamar a cinco testigos; al leer la demanda advierte que el actor ofreció a tres de ellos (Julián, Oscar y Pedro), que también le son útiles. Debe \textbf{ofrecerlos igual}: aunque sean las mismas personas, las preguntas serán distintas según los hechos que cada parte quiera probar; y, además, el actor puede desistir después de un testigo que ya no le convenga, y si el demandado no lo había ofrecido se queda sin él.
\end{ejemplo}

\begin{importante}
\textbf{La prueba es de la causa, no de quien la ofrece.} Una vez incorporada, no importa quién la presentó. Aun cuando la otra parte la haya ofrecido, conserva la facultad de desistir; por eso, si algo importa, debe ofrecerse igual. No es superabundante: cada parte manejará la prueba como propia. Vale para pericias y cualquier otra prueba.
\end{importante}

\section{La reconvención}

\begin{definicion}{Reconvención}
Contrademanda del demandado, planteada por los mismos hechos y con las mismas personas, con relación con la pretensión de la demanda; con ella el demandado se convierte en una suerte de actor. Requiere tres elementos: las \textbf{mismas partes}, los \textbf{mismos hechos} y la \textbf{misma relación jurídica}.
\end{definicion}

\begin{ejemplo}{el usufructo vitalicio de la pareja}
Una pareja que convivió muchos años sin casarse compró varias propiedades; una la compró el señor y constituyó un usufructo vitalicio a favor de la señora. Ella reclama que el usufructo es un acto simulado (no tiene sentido tras comprar tantas propiedades juntos; ella no podía acreditar ingresos y él sí, de modo que se constituyó para resguardarla) y pide que se declare que nunca reflejó una realidad. Él responde que la propiedad es suya y reconviene: pasaron diez años y ella nunca ejerció el usufructo (que se ejerce habitando o percibiendo los frutos). Los sujetos son los mismos, los hechos son los mismos (la compra y el usufructo) y se discute la misma relación jurídica (si existe o no el usufructo, que cada parte interpreta de modo distinto); los demás elementos se probarán luego.
\end{ejemplo}

\begin{nota}
\textbf{Referencia complementaria, art.~357 CPCCN.}
% TODO: verificar referencia presente en las fuentes originales (texto abreviado; cotejar con el Código vigente).
«Reconvención. Si el demandado reconviniere, deberá hacerlo en el escrito de contestación [\ldots]. La reconvención [será admisible] si las pretensiones deducidas derivaren de la misma relación jurídica o fueren conexas por el objeto o por el título [\ldots]». Las fuentes no citan el artículo.
\end{nota}

La reconvención debe plantearse \textbf{en la contestación}, dentro del mismo plazo de quince días; si se olvida, habrá que iniciar un juicio nuevo (que irá al mismo juzgado pero será otro). Puede ubicarse en medio de la contestación, lo cual no se sugiere (en un caso se corrió traslado de la contestación pero nunca de la reconvención), o, preferentemente, \textbf{al final}, de manera separada y clara. Primero se acreditan los tres requisitos y luego lo que se reconviene. Debe ofrecerse prueba para ambas cosas aunque coincidan (se la menciona dos veces). En el objeto puede decirse: «vengo a contestar demanda, a ofrecer prueba y a plantear reconvención», y al final pedir que se tenga por contestada la demanda, ofrecida la prueba, movida la reconvención y ofrecida su prueba. De la reconvención se corre traslado al actor, que se defiende como en una contestación; recién entonces queda trabada la litis.

\begin{table}[H]
\centering
\caption{Contestación de demanda y reconvención}
\begin{tabularx}{\textwidth}{>{\raggedright\arraybackslash\bfseries}p{0.2\textwidth} X X}
\toprule
 & Contestación & Reconvención \\
\midrule
Posición & Defensa. & Contraataque: el demandado es una suerte de actor. \\
Ubicación & Primera parte del escrito. & Al final, separada y clara. \\
Plazo & 15 días. & El mismo; si se olvida, juicio nuevo. \\
Prueba & Ofrece la suya. & Ofrece la suya (aunque coincida, se menciona dos veces). \\
Pedido final & Tener por contestada la demanda y ofrecida la prueba. & Tener por movida la reconvención y ofrecida su prueba. \\
Qué sigue & --- & Traslado al actor; luego queda trabada la litis. \\
\bottomrule
\end{tabularx}
\end{table}

\begin{cornell}{La defensa del demandado}
\cq{¿Cuál es el orden de trabajo ante una demanda?}{Forma, hechos, derecho, prueba: excepciones primero; luego, negar hechos o aceptarlos y discutir su consecuencia jurídica.}
\cq{¿Qué es una excepción previa?}{Un defecto técnico que debe resolverse antes de seguir (incompetencia, litispendencia, falta de legitimación o personería, arraigo, prescripción, defecto legal, transacción). No reemplaza la defensa de fondo.}
\cq{¿Cómo se desconoce y se niega?}{Se desconoce documento por documento (el no desconocido se tiene por reconocido) y se niega hecho por hecho, con negativas puntuales; la negativa general no sirve.}
\cq{¿Por qué ofrecer la prueba que ya ofreció el actor?}{Porque la prueba es de la causa y el actor puede desistir de ella.}
\cq{¿Qué requisitos tiene la reconvención y cuándo se plantea?}{Mismas partes, mismos hechos y misma relación jurídica; solo en la contestación, preferentemente al final y con prueba propia.}
\cresumen{El demandado trabaja en capas: excepciones (que no excusan defender el fondo), desconocimiento de documentos, negativa puntual de hechos, relato de la realidad como estrategia y ofrecimiento de prueba («ahora o nunca»). Debe ofrecer aun la prueba que el actor ya ofreció y, si tiene un contrarreclamo con las mismas partes, hechos y relación jurídica, reconvenir en la contestación; de lo contrario deberá iniciar un juicio nuevo.}
\end{cornell}

\part{La audiencia preliminar y la prueba}

\chapter{La audiencia preliminar y la fijación del objeto del juicio}\label{cap:audiencia}

Trabada la litis, el proceso llega a un momento decisivo: el juez debe decidir de qué se discutirá y qué se probará. Este capítulo estudia la cuestión de puro derecho, la audiencia preliminar, el protocolo de oralidad y una audiencia observada en la práctica.

\section{Hechos controvertidos o puro derecho}

Trabada la litis, el juez analiza la cuestión planteada. Si puede resolverse aplicando únicamente el derecho, sin hechos controvertidos, la declara de \textbf{puro derecho}; si hay hechos en disputa y conducentes, abre a prueba y cita a audiencia preliminar.
% TODO: verificar consistencia entre fuentes originales (un apunte atribuye esta regla al «artículo 139» del Código y otro, al artículo 359 CPCCN).

\begin{definicion}{Cuestión de puro derecho}
Simplificación procesal que opera cuando los hechos no son controvertidos o no hay prueba que producir; el debate se concentra en la interpretación y aplicación de las normas, sin etapa probatoria. No se discute qué ocurrió sino quién tiene mayor apoyo en la norma. Si se discute, por ejemplo, la constitucionalidad de una ley, no hace falta probar cómo ocurrió algo sino interpretar la Constitución.
\end{definicion}

\begin{regla}{art.~359 CPCCN}
Contestada la demanda o la reconvención, y no quedando cuestiones previas que resolver, el juez tiene la facultad de resolver la cuestión de puro derecho.
\end{regla}

Procede, por ejemplo, cuando el demandado admite los hechos como los describe el actor, cuando las partes no tienen pruebas que ofrecer y aceptan continuar sin etapa probatoria, o cuando en la audiencia preliminar el juez declara que los hechos están claros. El procedimiento es: el juez concede 5 días a quien inició el proceso para que se expida; se corre traslado a la otra parte por 5 días; las partes presentan \textbf{alegatos}; la resolución es apelable si causa agravio; de no haber apelación, el juez dicta sentencia. En el \textbf{juicio ejecutivo}, el título otorga presunción de legitimidad: se da traslado por 10 días para oponer excepciones y, si no las hay o se rechazan, se dicta sentencia de remate; también configura cuestión de puro derecho, pues se discute la validez del título y no los hechos subyacentes.

\section{La audiencia del artículo 360}

La audiencia preliminar es obligatoria y es el corazón del proceso civil moderno. En ella ocurren tres cosas: la conciliación, la determinación del objeto de la litis y la fijación de la prueba.

\begin{table}[H]
\centering
\caption{Funciones de la audiencia preliminar}
\begin{tabularx}{\textwidth}{>{\raggedright\arraybackslash\bfseries}p{0.28\textwidth} X}
\toprule
Etapa & Descripción \\
\midrule
Conciliación & El juez llama a conciliar. Generalmente no hay acuerdo, pero debe formular la pregunta. \\
Objeto de la litis & Identifica los hechos controvertidos y conducentes y pregunta la versión de cada parte. \\
Fijación de la prueba & Se ofrecen los medios; el juez va depurando la prueba, admitiéndola o desestimándola según su conducencia. \\
\bottomrule
\end{tabularx}
\end{table}

\begin{regla}{art.~360 CPCCN}
El juez, en la audiencia preliminar, debe fijar los hechos articulados que sean conducentes a la decisión del juicio sobre los que versará la prueba, recibir la prueba confesional si fue ofrecida y proveer las demás pruebas que considere admisibles. Su primer paso, antes de fijar el objeto, es promover la conciliación; las partes pueden pedir una audiencia de conciliación en cualquier momento anterior a la sentencia.
\end{regla}
% TODO: verificar referencia presente en las fuentes originales (no queda claro si el texto entrecomillado de las notas es cita literal del artículo).

\begin{definicion}{Hechos controvertidos y conducentes}
Un hecho es \textbf{controvertido} cuando una parte lo afirma y la otra lo niega; es \textbf{conducente} cuando es relevante para resolver el pleito. Solo los hechos controvertidos y, a la vez, conducentes son los que el juez fija en la audiencia y sobre los que se produce prueba.
\end{definicion}

A veces las partes afirman muchas cosas y la mitad, aunque controvertida, no sirve para resolver. El objetivo de la audiencia es que el juez tenga mayor contacto directo con el caso y las partes, conozca los hechos de primera mano y, si se hace bien, reduzca significativamente la duración del proceso, descartando pruebas impertinentes y redundantes.

\begin{ejemplo}{una audiencia mal llevada}
En una audiencia complicada, ya virtual, el juez no fijó los hechos controvertidos y conducentes: dijo directamente «bueno, abrimos a prueba». Se abrió la prueba y se pidió, por ejemplo, un oficio para acreditar un dato irrelevante: se pierde tiempo porque ese dato no cambia la resolución. Frente a cada hecho el juez debe preguntarse si sirve o no para resolver.
\end{ejemplo}

\section{Oralidad}

\begin{definicion}{Protocolo de oralidad}
Política impulsada desde el Ministerio de Justicia antes de la pandemia, orientada a aplicar en el proceso civil, en la mayor medida posible, la producción oral de la prueba (especialmente la pericial) en la audiencia preliminar, en lugar de resolver todo por escrito.
\end{definicion}

No es lo mismo que el perito presente su informe por escrito que el juez pueda preguntarle directamente cómo se produjo un hecho técnico.

\begin{ejemplo}{el chalet con la filtración}
En un juicio por filtración en un chalet no era claro si el defecto obedecía a un error de construcción o a las tejas; la demanda se había dirigido contra todos los posibles responsables. En la audiencia el juez interrogó al perito, que explicó que el problema se originó en la colocación de las tejas y no en su fabricación, lo que permitió identificar rápidamente al responsable. Por escrito habría sido más extenso, discutible y difícil. Tras la explicación, las partes dejaron de tener motivos para discutir, porque el resultado ya era previsible.
\end{ejemplo}

De aplicarse plenamente, permitiría reducir el tiempo de los procesos escritos a una fracción estimada del 20\% o menos, acercar al juez a los hechos y limitar maniobras dilatorias. La «audiencia de 360» equivale a la «audiencia de vista de causa» de varios códigos provinciales.

\begin{importante}
La oralidad exige que partes y letrados lleguen preparados para responder preguntas en vivo, sin ampararse en un escrito diferido. Quien no tiene claro qué sostiene, qué pretende y qué prueba ofrece queda expuesto ante el tribunal.
\end{importante}

\section{Una audiencia preliminar observada: el caso de la cochera}

La audiencia observada correspondió al expediente «Chester, Claudio c/ Construcciones Mirage S.A. s/ prescripción adquisitiva» (véase la sección~\ref{sec:usucapion}). Al abrirse, se dejó constancia de que no compareció representante del Gobierno de la Ciudad (citado por el interés fiscal) ni de la demandada.

\textbf{Interrogatorio al actor.} Relató que en 1989 compró por boleto la cochera, unidad funcional 28 del edificio de Independencia 1452.
% TODO: contenido incompleto o ambiguo en las notas originales (el piso indicado en la transcripción es «64», dato que no resulta verosímil y que se conserva tal como figura en el apunte).
Pagó todo el precio pero nunca se escrituró; la usó personalmente un tiempo y luego la alquiló por contratos verbales, pagando siempre los gastos. No tenía certeza de si otras unidades atravesaban una situación similar; no participaba en las asambleas pero siempre se sintió propietario; insistió sin éxito en la escrituración y aclaró que no quiere vender sino regularizar su situación.

\textbf{Testigos.} Antes de admitir a los propuestos por el actor (Alberto Damián Cabuli y María Teresa Gamboa), la jueza examinó su utilidad: uno lo acompañó en trámites y le alquiló la cochera; la segunda trabajaba con el primero. Ambos fueron admitidos por su vinculación con los hechos.

\begin{importante}
La admisión de un testigo no depende de la mera oferta: la prueba debe estar dirigida a acreditar los hechos controvertidos (aquí, la posesión por más de veinte años). No corresponde citar a quien no pueda aportar algo relevante. Es un error frecuente ofrecer un testigo sin explicar para qué hecho concreto.
\end{importante}

\textbf{Prueba desistida.} El actor mantuvo la testimonial y la informativa (Gobierno de la Ciudad, Aguas Argentinas y consorcio), pero desistió de la pericial caligráfica y de la inspección ocular por redundantes: el desinterés del Gobierno y la informativa cumplían la misma función. Cuando una prueba puede ser suplida por otra ya ordenada, es razonable desistir; en ciertos casos la inspección ocular puede sustituirse por medios audiovisuales. Se dejó además constancia de que quien firmó el boleto por Construcciones Mirage S.A. representaba a una sociedad disuelta desde 2001.

\begin{regla}{acta de la audiencia (síntesis)}
«Buenos Aires, a los 18 días de septiembre de 2026 [\ldots] se resuelve abrir la causa a prueba. Se ordena prueba informativa a Aguas Argentinas, al Gobierno de la Ciudad y al consorcio propietario [\ldots]. La testimonial se llevará a cabo el día de la audiencia de vista de causa [\ldots], con comparecencia presencial en la sede del juzgado.» (Reconstruida a partir de las notas originales.)
\end{regla}

La jueza mantuvo el control del acto y preguntó directamente al actor todo lo necesario; de tratarse de un caso más complejo (aquí había boleto, pago de expensas y un único punto débil, la falta de escrituración) habría sido probablemente más incisiva. Queda planteado un desafío: cómo continuar frente a una demandada que ya no existe, cuando el Gobierno de la Ciudad manifestó su desinterés, y cómo obtener información institucional cuando no hay contraparte activa: de allí la figura del oficio (sección~\ref{sec:oficios}).

\begin{cornell}{Audiencia preliminar}
\cq{¿Cuándo hay puro derecho?}{Cuando no hay hechos controvertidos o prueba que producir; se debate solo la norma (art.~359).}
\cq{¿Qué ocurre en la audiencia del art.~360?}{Conciliación, fijación del objeto (hechos controvertidos y conducentes) y fijación de la prueba.}
\cq{¿Qué es un hecho conducente?}{El relevante para resolver el pleito; solo los controvertidos y conducentes se prueban.}
\cq{¿Qué exige la oralidad?}{Preparación para responder en vivo; a cambio, abrevia el proceso y acerca al juez a los hechos.}
\cq{¿Cuándo se admite un testigo?}{Cuando aporta algo relevante sobre un hecho controvertido; hay que explicar para qué se lo ofrece.}
\cresumen{Trabada la litis, el juez decide si hay puro derecho o si abre a prueba. En la audiencia preliminar promueve la conciliación, fija los hechos controvertidos y conducentes y depura la prueba; la oralidad la hace más eficaz pero exige preparación. La audiencia observada muestra interrogatorio a la parte, admisión de testigos por utilidad y desistimiento de prueba redundante.}
\end{cornell}

\chapter{Teoría general de la prueba}\label{cap:prueba}

Probar es convencer al juez. Este capítulo reúne los conceptos comunes a todos los medios: qué es la prueba, cuándo se ofrece, quién debe probar, qué debe probarse y cómo se valora.

\section{Concepto, fuentes y medios}

\begin{definicion}{Prueba}
Conjunto de medios, instrumentos y actuaciones mediante los cuales las partes llevan convicción de sus pretensiones al juez.
\end{definicion}

\begin{definicion}{Fuente de prueba y medio de prueba}
Las fuentes existen en la realidad (un documento, un testigo, un perito); los medios son los mecanismos legales para introducirlas en el proceso. Puede existir una fuente sin medio legal para ingresarla.
\end{definicion}

\begin{table}[H]
\centering
\caption{Clasificación de los medios de prueba}
\begin{tabularx}{\textwidth}{>{\raggedright\arraybackslash\bfseries}p{0.32\textwidth} X}
\toprule
Medio & Descripción \\
\midrule
Documental & Documentos públicos o privados presentados en el expediente. \\
Testimonial & Declaración de testigos que presenciaron los hechos. \\
Informes & Datos técnicos o administrativos provenientes de organismos. \\
Absolución de posiciones & Declaración de la parte en juicio sobre los hechos. \\
Reconocimiento judicial & Inspección directa del lugar por el juez. \\
Pericial & Dictamen de peritos expertos. \\
Medios judiciales & Otros medios que disponga el juez. \\
\bottomrule
\end{tabularx}
\end{table}

\section{Momento de ofrecimiento}

\begin{importante}
En el proceso civil y comercial la prueba debe ofrecerse con la demanda y con la contestación; si no, no hay posibilidad de hacerlo después. Este régimen difiere de otros fueros, como el laboral, donde ciertas pruebas (salvo la documental) pueden ofrecerse en otro momento. La excepción es el hecho nuevo (capítulo~\ref{cap:hecho-nuevo}).
\end{importante}

\section{Probar el hecho y probar el daño}

\begin{importante}
Para reclamar daños deben probarse \textbf{dos cosas distintas}: el \textbf{hecho causante} y los \textbf{daños concretos} sufridos. Probar el hecho no basta para acreditar los daños: sin demostrar el «antes» y el «después» no hay daño acreditable. La estrategia también debe decidir qué hechos se probarán por qué medios: primero se definen los hechos, luego los medios (por ejemplo, testigos) y finalmente el cuestionario de cada uno.
\end{importante}

\section{Carga de la prueba y carga dinámica}

\begin{definicion}{Derecho dispositivo}
El proceso civil es dispositivo cuando las partes disponen y son ellas quienes proponen las actuaciones.
\end{definicion}

\begin{definicion}{Carga de la prueba}
Corresponde a quien alega un hecho: el actor prueba sus pretensiones y el demandado sus excepciones. Originalmente era el actor quien probaba; hoy se interpreta menos literalmente y \textbf{ambas partes} prueban lo que alegan, aunque el actor tenga en principio una carga mayor.
\end{definicion}

\begin{nota}
\textbf{Referencia complementaria, art.~377 CPCCN.}
% TODO: verificar referencia presente en las fuentes originales (las fuentes no consignan el número; texto abreviado).
«Incumbirá la carga de la prueba a la parte que afirme la existencia de un hecho controvertido o de una norma jurídica que no sea de derecho nacional. Cada una de las partes deberá probar el presupuesto de hecho de la norma o normas que invocare como fundamento de su pretensión, defensa o excepción. [\ldots]»
\end{nota}

\begin{definicion}{Carga dinámica de la prueba}
Desplazamiento de la carga probatoria hacia quien está en mejor posición para probar, aun cuando no sea quien alega el hecho; el juez puede invertir la carga cuando la otra parte está en mejor posición.
\end{definicion}

\begin{ejemplo}{responsabilidad médica}
Un paciente operado queda con secuelas y demanda por mala praxis. Debe probar la lesión (fácil) y que la causó una negligencia (difícil); el médico está en mejor posición de probar que actuó correctamente, por lo que la carga se desplaza hacia él. En otro caso, una paciente operada en una clínica grande en una intervención sencilla presenta, al recibir anestesia, un cuadro de presión; tras otra medicación queda casi en coma, y tres años después recién empieza a salir sin saberse qué secuelas quedarán. La clínica deberá demostrar que, siendo la medicina una ciencia no exacta, siguió los protocolos, tomó previsiones, hizo los análisis y el prequirúrgico correspondientes y que, aun así, ocurrió algo inesperado; para la actora es mucho más difícil.
\end{ejemplo}

\begin{importante}
La carga dinámica no es automática: la decisión es del juez, pero la parte debe invocarla activamente, explicar por qué la otra está en mejor situación y pedir su aplicación. Es una herramienta poderosa que muchos abogados dejan pasar.
\end{importante}
% TODO: contenido incompleto o ambiguo en las notas originales (las fuentes no consignan el artículo aplicable).

\section{Sana crítica y valoración de la prueba}

\begin{definicion}{Sana crítica}
Sistema por el cual el juez valora las pruebas: no es un criterio único ni objetivo sino un estándar de racionalidad; las conclusiones deben ser lógicas, coherentes y fundadas.
\end{definicion}

Una sentencia correcta se estructura paso a paso: se identifican los hechos controvertidos probados (A, B, C); de la testimonial surge X; de la pericial surge Y; por lo tanto, se concluye Z. Entre lo que ocurrió, lo que se prueba y lo que sentencia el juez hay tres cosas distintas: la verdad real, la prueba como acercamiento a ella y la verdad legal (véase la sección~\ref{sec:concepto-proceso} y el capítulo~\ref{cap:sentencia}).

\begin{cornell}{Teoría de la prueba}
\cq{¿Fuente o medio de prueba?}{La fuente existe en la realidad; el medio es el mecanismo legal para introducirla al proceso.}
\cq{¿Cuándo se ofrece la prueba?}{Con la demanda y con la contestación; después, solo por hecho nuevo.}
\cq{¿Qué hay que probar para reclamar daños?}{El hecho causante y los daños concretos (antes y después).}
\cq{¿Quién prueba?}{Cada parte lo que alega; la carga dinámica se desplaza a quien está en mejor posición, pero debe invocarse.}
\cq{¿Qué es la sana crítica?}{Estándar de racionalidad con que el juez valora la prueba: conclusiones lógicas, coherentes y fundadas.}
\cresumen{La prueba lleva convicción al juez a través de medios legales que introducen fuentes. Debe ofrecerse con demanda y contestación; hay que probar hecho y daño; ambas partes prueban lo que alegan y puede invocarse la carga dinámica; el juez valora según la sana crítica.}
\end{cornell}

\chapter{Prueba documental, informativa y oficios}

Estos medios acreditan hechos mediante documentos y la información que poseen entidades. Se estudian juntos porque en la práctica se combinan: lo que no consta en un documento propio suele requerirse a un tercero.

\section{La prueba documental}\label{sec:documental}

\begin{definicion}{Prueba documental}
Todo instrumento o soporte que acredita la existencia de un hecho, de una relación jurídica o de una circunstancia; no se limita al papel sino que comprende cualquier soporte con información relevante.
\end{definicion}

Ejemplos: el acta de matrimonio, el acta constitutiva de una empresa, el permiso de navegación, filmaciones de celular, el acta de cierre de mediación, comprobantes de gastos médicos, el expediente penal.

\begin{importante}
La documental debe ofrecerse con la demanda. Si el documento no se posee, se cumple la carga indicando cuál es, dónde se encuentra y todos los elementos disponibles (por ejemplo, una historia clínica que obra en un hospital). El instrumento no acompañado con la demanda no puede presentarse después, salvo hecho nuevo, y no pueden presentarse en la audiencia del art.~360 documentos no ofrecidos: corresponde entonces iniciar el incidente respectivo mediante escrito específico.
\end{importante}

\subsection{La carga del demandado ante los documentos}

Al contestar, el demandado debe pronunciarse sobre cada documento: reconocerlo o desconocerlo; la omisión implica reconocimiento. El desconocimiento debe ser puntual, indicando el motivo (por inexacto o falso), no genérico.

\begin{table}[H]
\centering
\caption{Instrumento público e instrumento privado}
\begin{tabularx}{\textwidth}{>{\raggedright\arraybackslash\bfseries}p{0.22\textwidth} X X}
\toprule
 & Instrumento público & Instrumento privado \\
\midrule
Cómo se lo cuestiona & Mediante redargución de falsedad, procedimiento específico. & Se lo desconoce. \\
Qué ocurre después & Se presume que hace fe pública; cuestionarlo es dificultoso. & Quien lo presentó debe probar que la firma es auténtica, o acreditarlo por otro medio. \\
\bottomrule
\end{tabularx}
\end{table}

\begin{importante}
Desconocer sistemáticamente todos los documentos, incluso los evidentes, no es una estrategia inteligente: resta credibilidad ante el juez. Conviene reconocer lo que no puede negarse y controvertir con precisión lo restante.
\end{importante}

\section{Exhibición de documentos en poder de la contraparte o de terceros}

\begin{regla}{art.~388 CPCCN --- exhibición de documentos}
Las partes y los terceros en cuyo poder se encuentren documentos esenciales para la solución del litigio estarán obligados a exhibirlos o a designar el protocolo o archivo en que se hallan los originales.
\end{regla}

\begin{regla}{art.~389 CPCCN --- documentos en poder de un tercero}
Si el documento que deba reconocerse se encontrare en poder de un tercero, se lo intimará para que lo presente. Si lo acompañase, podrá solicitar su oportuna devolución, dejando testimonio en el expediente. El requerido podrá oponerse a su presentación si el documento fuere de su exclusiva propiedad y la exhibición pudiera ocasionarle perjuicio. Ante la oposición formal del tenedor, no se insistirá en el requerimiento.
\end{regla}

Cuando una parte no cuenta con una documental esencial pero sabe que un tercero la posee, puede invocar el art.~389 sin que ese tercero sea demandado; con frecuencia el poseedor es la propia contraparte.

\begin{ejemplo}{la póliza de seguro}
La póliza habitualmente no está en poder del actor sino de la aseguradora demandada. Intimar su presentación permite probar la existencia del contrato, que este obliga a la aseguradora en ciertos supuestos y que el riesgo reclamado está cubierto. Puede intimarse en la demanda o mediante una diligencia preliminar previa (capítulo~\ref{cap:incidentes}).
\end{ejemplo}

\section{Prueba informativa}

Consiste en solicitar, mediante un oficio, que una entidad pública o privada informe sobre una cuestión (por ejemplo, a un banco sobre una cuenta).

\begin{ejemplo}{acreditar ingresos mediante oficios}
Para acreditar que una persona era docente en ciertas universidades, puede librarse oficio a esas instituciones para que informen desde cuándo ejercía el cargo y cuánto percibía; así se acredita indirectamente el ingreso dejado de percibir por el hecho dañoso.
\end{ejemplo}

\section{Los oficios judiciales}\label{sec:oficios}

\begin{definicion}{Oficio judicial}
Pedido de comunicación o de información dirigido siempre a una \textbf{entidad}, pública o privada, nunca a una persona física, para que los datos de sus registros se pongan a disposición del proceso.
\end{definicion}

\begin{nota}
Si se busca información de una entidad, corresponde un oficio; si depende del conocimiento de una persona física, corresponde ofrecerla como testigo.
\end{nota}

Hay dos categorías según su finalidad. El \textbf{oficio 3003} se emplea en procesos sucesorios o concursales: antes de promover una sucesión, una quiebra o un proceso universal, se dirige al Registro de Procesos Universales para verificar si ya existe uno a nombre de la persona o empresa. El otro, de uso general, es el del artículo 400 CPCCN.

\begin{regla}{art.~400 CPCCN --- atribuciones de los letrados patrocinantes}
«Los pedidos de informes, testimonios y certificados, así como los de remisión de expedientes ordenados en el juicio, serán requeridos por medio de oficios firmados, sellados y diligenciados por el letrado patrocinante, con transcripción de la resolución que los ordena y que fija el plazo en que deberán remitirse.»
\end{regla}

Con el expediente digital, buena parte de esas exigencias formales se simplificó; solo cobran relevancia cuando la entidad no figura en el sistema.

\subsection{Diligenciamiento electrónico}

El módulo de oficios requiere completar, en orden: jurisdicción y número de expediente; entidad destinataria; tipo de oficio (art.~400); texto del despacho que ordena la medida; y oficina de recepción. Completado, el oficio queda cargado como diligenciado, sin escrito adicional pidiendo que se diligencie, como se exigía en papel.

\subsection{El oficio en papel}

Si el destinatario no figura en el sistema (por ejemplo, ciertas empresas de comercio electrónico), se vuelve al procedimiento en papel.

\begin{regla}{estructura del oficio en papel (modelo)}
«Tengo el agrado de dirigirme a usted en mi calidad de letrado patrocinante [o apoderado] de la parte actora [o demandada] en los autos caratulados «\ldots», que tramitan por ante el Juzgado [\ldots], a cargo de [\ldots], Secretaría a cargo de [\ldots], con domicilio en [\ldots], a fin de solicitarle quiera tener a bien informar [el objeto del pedido] [\ldots]. [Se transcribe el auto que ordena la medida]. Lo saludo atentamente.»
\end{regla}

Al transcribir la resolución no se copia íntegra sino solo el fragmento operativo (por ejemplo, «líbrese el oficio solicitado»). Para acreditar el diligenciamiento cuando la respuesta se recibe en papel y sin recibo digital, se confeccionan dos copias: una queda con el letrado y otra se presenta ante la institución para obtener el sello de recibido; se digitaliza y se presenta un escrito en el expediente electrónico acreditando el diligenciamiento. Al cargarlo se selecciona el expediente principal (no un incidente con la misma carátula), el tipo de escrito (habitualmente «mero trámite») y se adjunta en un único PDF el escrito y la copia del oficio.

\begin{cornell}{Documental, informativa y oficios}
\cq{¿Cuándo se ofrece la documental y qué pasa si no se posee?}{Con la demanda; si no se posee se indica cuál es y dónde está. Lo no ofrecido no puede presentarse después, salvo hecho nuevo.}
\cq{¿Cómo se desconoce un documento?}{Puntualmente y con motivo; el no desconocido se reconoce. El instrumento público se cuestiona por redargución de falsedad.}
\cq{¿Qué prevén los arts.~388 y 389?}{La obligación de exhibir documentos esenciales y la intimación a terceros que los poseen, con posibilidad de oposición.}
\cq{¿A quién se dirige un oficio?}{A una entidad, nunca a una persona física; si depende del conocimiento de una persona, se la ofrece como testigo.}
\cq{¿Cómo se acredita un oficio en papel?}{Con dos copias, una sellada por el destinatario; se digitaliza y se presenta escrito de mero trámite en el expediente principal.}
\cresumen{La documental se ofrece con la demanda y se reconoce o desconoce documento por documento. Los documentos en poder de la contraparte o de terceros pueden exigirse (arts.~388 y 389). La informativa se canaliza por oficios a entidades (art.~400), hoy diligenciados electrónicamente; el oficio en papel se acredita con copia sellada.}
\end{cornell}

\chapter{La prueba testimonial}

El testimonio permite reconstruir hechos a través de quien los percibió. Este capítulo estudia el concepto de testigo, las cargas de quien lo ofrece, su preparación y la técnica de interrogatorio.

\section{El testigo}

\begin{definicion}{Testigo}
Quien conoce por sus propios sentidos (lo vio, lo escuchó o lo percibió) un hecho relevante para el proceso y lo relata de manera directa, sin que se trate de algo que le contaron (lo que lo distinguiría de un testigo de oídas o de referencia).
\end{definicion}

Se destaca la \textbf{imparcialidad}: no debe tener interés personal en el resultado. Cualquiera de las partes puede ofrecer testigos que reúnan esas condiciones. No todo testigo declara sobre el hecho principal.

\begin{table}[H]
\centering
\caption{Testigos del hecho y de las consecuencias}
\begin{tabularx}{\textwidth}{>{\raggedright\arraybackslash\bfseries}p{0.22\textwidth} X X}
\toprule
 & Testigo del hecho & Testigo de las consecuencias \\
\midrule
Quién & Una persona presente en el momento del hecho. & Amigos o conocidos del actor. \\
Qué acredita & El hecho mismo. & La situación anterior y posterior de la persona, usada para el daño moral. \\
\bottomrule
\end{tabularx}
\end{table}

\begin{ejemplo}{testimonial para daño moral y para pérdida de chance}
Testigos que frecuentaban el club con la pareja pueden declarar: «Venían al club todos los días como pareja; ahora este señor ya no viene más». Socios o compañeros de trabajo pueden declarar: «Eran socios hace diez años; trabajaban juntos y participaban de la vida social. Ahora, a raíz del hecho, el sobreviviente ya no rema, no navega, le tiene miedo al agua». Los ingresos dejados de percibir se acreditan con prueba informativa.
\end{ejemplo}

\section{Cargas de la parte que ofrece un testigo}

\begin{table}[H]
\centering
\caption{Cargas del oferente}
\begin{tabularx}{\textwidth}{>{\raggedright\arraybackslash\bfseries}p{0.30\textwidth} X}
\toprule
Momento & Carga \\
\midrule
Al ofrecerlo & Identificarlo (nombre y, en la práctica, DNI) e indicar sobre qué hechos declarará. \\
Al redactar la demanda & Precisar respecto de qué hechos se cita a cada testigo (no es obligatorio incluir ya el cuestionario). \\
Después de ofrecido & Asegurar su comparecencia: el testigo es citado por el juzgado, salvo que la parte asuma hacerlo comparecer. \\
\bottomrule
\end{tabularx}
\end{table}

Aunque los jueces, por la amplitud probatoria, rara vez desestiman un testigo por no precisar su objeto, la contraparte tiene derecho a saber a qué viene cada uno (defensa en juicio y bilateralidad). Respecto de la comparecencia hay una carga que debe cumplirse sí o sí: \textbf{notificar} al testigo. Si no se le envía cédula y el día de la audiencia avisa que no puede presentarse, la parte se queda sin ese testigo y con las preguntas preparadas. Si es conocido y se lo avisa y concurre, no hay problema aunque falte la cédula. Otra alternativa, a veces más violenta pero válida, es hacerle firmar una constancia de que tomó conocimiento de que determinado día y hora fue citado como testigo en tal causa, juzgado y lugar, y presentarla al juzgado en lugar de la cédula; con todo, lo más seguro es la cédula real. Si los testigos no comparecen por falta de notificación (en el ejemplo, por carta documento), la contraparte debe pedir en el acto la \textbf{caducidad} de la prueba testimonial (art.~432; véase la sección~\ref{sec:caducidad-prueba}); no corresponde «impugnar», porque no hay acto que impugnar.

\section{Preparación del testigo y generales de la ley}

Antes de ofrecer un testigo hay que tomar contacto con él y confirmar que tiene conocimiento directo; suele ser el cliente quien identifica a esas personas. Preparar al testigo, sin aleccionarlo, es explicarle el procedimiento y para qué se lo llama («nosotros lo llamamos para que acredite tal cosa, simplemente eso»): el testigo suele comparecer con temor. Cuando declara deficientemente, la responsabilidad suele ser del abogado que lo ofreció por no haberlo preparado.

\begin{definicion}{Generales de la ley}
Procedimiento inicial por el cual se interroga al testigo sobre su situación personal frente a las partes, antes de declarar sobre los hechos.
\end{definicion}

\begin{importante}
Debe explicarse al testigo que declara bajo juramento y que, si miente, hay consecuencias penales (aunque es parte normal del procedimiento y no debe generarle temor); que se le preguntarán datos personales; que se le consultará si conoce a la parte o mantiene vínculo con ella, debiendo responder con precisión y sinceridad sobre el grado de esa relación, sin que la amistad le impida ser objetivo; y si hay deudas entre él y la parte o si es acreedor.
\end{importante}

\section{Técnica de interrogatorio}

La prueba se produce según dos modalidades: el tribunal interroga primero (juez o secretario formulan las preguntas y luego consultan al abogado si quiere ampliar) o las partes formulan las preguntas directamente.

\begin{importante}
La pregunta nunca debe poder contestarse con sí o no ni presuponer el dato que se pretende obtener. Se comienza con preguntas amplias y abiertas y se avanza sobre las respuestas. Frente a una pregunta \textbf{indicativa} (de la contraparte o del tribunal) corresponde oponerse.
\end{importante}

\begin{table}[H]
\centering
\caption{Preguntas cerradas o indicativas y preguntas abiertas}
\begin{tabularx}{\textwidth}{X X}
\toprule
\textbf{A evitar} & \textbf{Pregunta abierta (correcta)} \\
\midrule
«¿Hace cuánto formaban parte del club?»: presupone que el testigo conocía al actor y que eran socios. & «¿De dónde lo conoce?» \\
«¿Frecuentaba tal club?»: incorpora el dato en la pregunta. & «¿Qué actividades realizaba el actor?» \\
«¿El semáforo estaba en rojo cuando se produjo el choque?» & «¿Qué sucedió?», «¿qué vio?» \\
Cualquier pregunta que se conteste por sí o por no. & «¿Con qué frecuencia se veían?», «¿dónde?», «¿desde cuándo?»: cada una se apoya en la respuesta anterior. \\
\bottomrule
\end{tabularx}
\end{table}

La forma de objetar es pedir la palabra y decir, por ejemplo, «me opongo, la pregunta es indicativa», pidiendo su reformulación. En lo penal el interrogatorio suele ser aún más estrictamente abierto («¿qué hizo?, ¿dónde estaba?, ¿tenía el arma?»).

\begin{ejemplo}{el juez que preguntaba de manera indicativa}
En un juicio, el juez interrogó con preguntas indicativas; al advertírsele y pedírsele que preguntara qué actividad realizaba el testigo, respondió que siempre había preguntado así, lo que derivó en un cuestionamiento posterior ante la Cámara.
% TODO: las notas originales presenta fragmentos poco claros en esta anécdota; se conserva el sentido general.
\end{ejemplo}

\begin{table}[H]
\centering
\caption{Orden y derechos en el interrogatorio}
\begin{tabularx}{\textwidth}{>{\raggedright\arraybackslash\bfseries}p{0.2\textwidth} >{\raggedright\arraybackslash}p{0.27\textwidth} X}
\toprule
Momento & Quién pregunta & Qué ocurre \\
\midrule
Primero & La parte oferente (o el tribunal, según la modalidad). & Formula sus preguntas; el tribunal puede preguntar lo que considere que falta. \\
Luego & El tribunal consulta a la oferente. & Puede ampliar o considerar suficiente lo preguntado. \\
Después & La contraparte. & Tiene derecho a preguntar aunque no sea su testigo. \\
Si hay discusión & --- & Se retira al testigo de la sala, se discute y luego se lo hace volver. \\
\bottomrule
\end{tabularx}
\end{table}

Ante una discusión sobre si una pregunta es indicativa, primero se retira al testigo; quien preside invita a las partes a reformularla de común acuerdo; si no lo logran, el tribunal puede reformularla por su decisión o disponer que no se formule. La parte disconforme con una decisión sobre la prueba dispone de la reposición (sección~\ref{sec:reposicion}). Frente a la jueza que interroga directamente a la parte, ello es una facultad del tribunal, no una irregularidad: el cliente debe estar preparado porque puede suceder.

\begin{cornell}{Prueba testimonial}
\cq{¿Quién es testigo?}{Quien percibió por sus sentidos un hecho relevante y lo relata directamente; se exige imparcialidad. No es testigo el de oídas.}
\cq{¿Qué cargas tiene el oferente?}{Identificarlo, indicar sobre qué hechos declarará y asegurar su comparecencia, en especial notificándolo.}
\cq{¿Qué ocurre si el testigo no fue notificado y no comparece?}{La contraparte pide la caducidad de la prueba (art.~432); la parte pierde el testigo.}
\cq{¿Qué son las generales de la ley?}{La indagación inicial sobre la relación del testigo con las partes, previa a declarar sobre los hechos.}
\cq{¿Cómo se pregunta?}{Con preguntas abiertas, que no sugieran la respuesta; las indicativas se objetan.}
\cresumen{El testimonio depende de la preparación del abogado: elegir testigos con conocimiento directo, explicarles el procedimiento, notificarlos y diseñar el cuestionario. En la audiencia se pregunta con preguntas abiertas, se objetan las indicativas y se resuelve en el acto cualquier discusión.}
\end{cornell}

\chapter{La prueba pericial}

Cuando los hechos exceden el conocimiento del juez, interviene un auxiliar de la justicia: el perito. Buena parte de la labor profesional consiste en trabajar adecuadamente con la pericia.

\section{Procedencia, perito único y puntos de pericia}

\begin{regla}{art.~457 CPCCN --- procedencia}
Será admisible la prueba pericial cuando la apreciación de los hechos controvertidos requiera conocimientos especiales en alguna ciencia, arte, industria o actividad técnica especializada.
\end{regla}

\begin{regla}{art.~458 CPCCN --- perito único; consultores técnicos}
La prueba pericial estará a cargo de un perito único designado de oficio por el juez, salvo cuando una ley especial establezca un régimen distinto. En los procesos de incapacidad e inhabilitación se estará al art.~626, inc.~3. En el juicio por nulidad de testamento, el juez podrá nombrar de oficio tres peritos cuando por la importancia y complejidad del asunto lo considere conveniente, impartiéndoles directivas sobre el modo de proceder. Cada parte tiene la facultad de designar un consultor técnico.
\end{regla}

\begin{regla}{art.~459 CPCCN --- designación; puntos de pericia (primer párrafo)}
Al ofrecer la prueba pericial se indicará la especialización que ha de tener el perito y se propondrán los puntos de pericia; si la parte ejerciera la facultad de designar consultor técnico, deberá indicar, en el mismo escrito, su nombre, profesión y domicilio.
\end{regla}

\begin{definicion}{Puntos de pericia}
Temas concretos sobre los que el perito debe trabajar e informar. Se indican al ofrecer la prueba, junto con la especialización requerida. Corrido el traslado, la contraparte puede impugnarlos, añadir o observar los propuestos. Si no hay acuerdo sobre perito y puntos, el juez los fija de oficio en la audiencia del art.~360, pudiendo agregar o eliminar puntos. Si nada se dice, el plazo para presentar el dictamen es, por principio general, de quince días.
\end{definicion}

\begin{importante}
Al ofrecer prueba pericial se debe ser específico sobre qué debe determinar el perito mediante preguntas concretas. Nunca debe escribirse al final «y cualquier otro punto de interés»: es un error grave que viola el debido proceso y hace que la contraparte pierda el derecho a cuestionar los puntos. Como el juez suele otorgar a la pericia un valor casi determinante, el diseño de los puntos es decisivo.
\end{importante}

\begin{ejemplo}{puntos de una pericia psicológica}
Se designa un perito psicólogo para determinar: grado o porcentaje de incapacidad psíquica; si es permanente o transitoria; si requiere tratamiento; en qué actividades la persona está impedida; tiempo estimado de tratamiento; y costo estimado. Cada punto es concreto y mensurable; el perito no decide qué responder y la contraparte puede impugnar los puntos si le parecen erróneos. Debe vincularse el daño con el hecho por relación causal.
\end{ejemplo}

\section{El perito: título, recusación, aceptación y remoción}

\begin{table}[H]
\centering
\caption{Reglas sobre la actuación del perito}
\begin{tabularx}{\textwidth}{>{\raggedright\arraybackslash\bfseries}p{0.25\textwidth} X}
\toprule
Tema & Regla \\
\midrule
Título habilitante & Si la actividad está reglamentada, debe presentarse título. Si no hay perito con título en la materia, puede designarse a otra persona con conocimientos suficientes (idóneo). \\
Recusación & Pueden recusarlo las partes, invocando causa justa, o el juez de oficio, por falta de título o incompetencia. Si reconoce la causa o calla, se lo reemplaza; si la niega, se tramita por separado, sin interrumpir el proceso. \\
Aceptación & Dentro del tercer día; si carece de título, debe jurar desempeñar fielmente el cargo. Hoy se formaliza por escrito electrónico. \\
Remoción & El juez puede removerlo de oficio si renuncia, no presenta el informe en plazo o se rehúsa a entregarlo. \\
\bottomrule
\end{tabularx}
\end{table}

Los peritos se inscriben anualmente en las listas de la cámara correspondiente (civil, penal, etcétera); excepcionalmente, si no hubiera perito matriculado, se llama a un idóneo. La remoción puede pedirse si no presenta el dictamen, se niega a hacerlo o renuncia sin motivo atendible (una enfermedad sobreviniente, por ejemplo, puede justificar una prórroga).

\section{Adelanto de gastos}

El perito puede solicitar adelanto de gastos y el juez resuelve; las partes pueden oponerse u observar el monto.

\begin{importante}
Si el juez concede el adelanto y la parte obligada no cumple, la prueba pericial no se produce, y la parte queda sin una prueba que puede ser fundamental.
\end{importante}

El beneficio de litigar sin gastos permite, al menos temporariamente, no afrontar los gastos, entre ellos los de la pericia (sección~\ref{sec:litigar-sin-gastos}).

\section{El consultor técnico}

\begin{table}[H]
\centering
\caption{Perito y consultor técnico}
\begin{tabularx}{\textwidth}{>{\raggedright\arraybackslash\bfseries}p{0.22\textwidth} X X}
\toprule
 & Perito & Consultor técnico \\
\midrule
Designación & El juez, de oficio (perito único). & Cada parte tiene la facultad de designar uno. \\
Dictamen & Presenta el informe pericial. & No presenta dictamen; no puede actuar como perito oficial. \\
Función & Acepta el cargo, realiza las operaciones y presenta el informe en plazo. & Asiste a las operaciones, acompaña al perito, puede sugerirle aspectos y asesora a la parte para consentir u objetar la pericia. \\
Requisitos & Título si la actividad está reglamentada; inscripción anual; recusable y removible. & Profesional o idóneo de confianza de la parte. \\
\bottomrule
\end{tabularx}
\end{table}

\begin{ejemplo}{las filtraciones en un departamento}
Ante una demanda por filtraciones en un edificio, el perito oficial debe determinar el origen, cómo resolverlo y su costo. El consultor de parte puede asistir a las operaciones, acompañar al perito cuando se toman fotografías y sugerirle que registre otro sector o analice otro aspecto.
\end{ejemplo}

\section{Impugnación, aclaraciones y valoración}

La objeción técnica de la pericia se denomina \textbf{impugnación}; es tan importante como poder pedir aclaraciones. Cuando el perito concluye que el trabajador no presenta incapacidad y la actora sostiene lo contrario, se impugna señalando con precisión las deficiencias: falta de aplicación de un método, omisión de un estudio necesario o de algún elemento relevante; según el caso puede ordenarse una nueva pericia o rechazarse la impugnación.

\begin{ejemplo}{recorrido de la pericia médica en un accidente laboral}
(1) Se sortea el perito médico, que debe aceptar el cargo. (2) Cita a la persona, la revisa y solicita estudios complementarios (por ejemplo, radiografías). (3) Elabora el informe en el plazo; si no lo presenta, puede ser removido. (4) Se da conocimiento a ambas partes, que pueden impugnar o pedir aclaraciones. (5) El juez lo tiene en cuenta al sentenciar.
\end{ejemplo}

\begin{importante}
El juez tiene en cuenta el informe pero no está obligado a basarse en él: puede acogerlo o apartarse de sus conclusiones, siempre con fundamento. De allí la importancia de impugnar y pedir aclaraciones en tiempo oportuno. En un caso con prueba pericial mecánica se observó que puede resultar sobreabundante si otros medios ya cubren el aspecto técnico.
\end{importante}

\begin{cornell}{Prueba pericial}
\cq{¿Cuándo procede?}{Cuando apreciar los hechos controvertidos requiere conocimientos especiales (art.~457).}
\cq{¿Quién es el perito y qué es el consultor?}{Perito único designado de oficio, que dictamina; el consultor técnico lo designa cada parte, asiste y asesora, pero no dictamina.}
\cq{¿Qué se propone al ofrecerla?}{Especialización del perito y puntos de pericia concretos; jamás «cualquier otro punto de interés».}
\cq{¿Qué pasa si no se paga el adelanto de gastos?}{La pericia no se produce y la parte pierde la prueba.}
\cq{¿Cómo se controla el dictamen?}{Con impugnación y pedido de aclaraciones; el juez lo valora sin quedar atado, con fundamento.}
\cresumen{La pericia suple el saber técnico que el juez no tiene. Requiere puntos de pericia concretos, perito único con título (o idóneo), atención al adelanto de gastos y, si es posible, un consultor técnico. Su resultado puede impugnarse y el juez lo valora libremente, con fundamento.}
\end{cornell}

\chapter{Hecho nuevo, caducidad y negligencia en la prueba}\label{cap:hecho-nuevo}

La regla de que toda la prueba se ofrece al inicio tiene una excepción, el hecho nuevo, y la producción de la prueba está sujeta a dos institutos que sancionan la inactividad: la caducidad y la negligencia.

\section{El hecho nuevo}

\begin{definicion}{Hecho nuevo}
Hecho jurídicamente relevante (simple hecho o acto jurídico) que ocurre con posterioridad al intercambio de demanda y contestación y guarda relación directa con el objeto del litigio. Es la excepción a la regla de ofrecer toda la prueba con la demanda: permite incorporar después una prueba antes desconocida.
\end{definicion}

Requisitos: que no se conociera al presentar la demanda o se haya producido con posterioridad, y que sea relevante para resolver. Plazos: cinco días desde que se tomó conocimiento; límite máximo, la audiencia del art.~360 (o cinco días antes de su celebración). La audiencia se notifica por cédula (art.~363) y los hechos nuevos pueden alegarse hasta el quinto día de notificada. Si aparece después de la audiencia, se canaliza por el sistema de recursos.
% TODO: verificar consistencia entre fuentes originales (los apuntes formulan de manera no del todo coincidente el límite temporal: «hasta la audiencia», «hasta cinco días antes» y «hasta el quinto día de notificada la audiencia»).

\begin{table}[H]
\centering
\caption{Ubicación temporal del hecho nuevo}
\begin{tabularx}{\textwidth}{>{\centering\arraybackslash}X >{\centering\arraybackslash}X >{\centering\arraybackslash\bfseries}X >{\centering\arraybackslash}X}
\toprule
Demanda & Contestación & Hechos nuevos (hasta 5 días antes de la audiencia art.~360) & Audiencia art.~360 \\
\bottomrule
\end{tabularx}
\end{table}

Se corre traslado a la otra parte «cuando se lo considere pertinente»; es ambiguo, pero por la bilateralidad y la defensa en juicio, si no se corre traslado la contraria no podrá refutarlos, lo que equivale en la práctica a rechazarlos. El rechazo se apela con efecto diferido (sección~\ref{sec:apelacion}).

\begin{ejemplo}{hecho nuevo en una causa de alimentos}
En una causa de alimentos, el demandado es despedido después de la demanda: el despido modifica su capacidad económica, base de la cuota, y se vincula con el objeto. A la inversa, si mejorara sustancialmente su situación patrimonial, la actora podría denunciarlo para fundar una cuota mayor.
\end{ejemplo}

\section{Caducidad y negligencia probatorias}\label{sec:caducidad-prueba}

Ambos institutos operan en la etapa probatoria y se vinculan con el principio dispositivo: quien no cumple una carga puede sufrir consecuencias adversas. No deben confundirse con la caducidad de instancia (sección~\ref{sec:caducidad-instancia}).

\begin{definicion}{Caducidad de la prueba}
Extinción de una carga procesal por el mero transcurso de un plazo fijado por la ley. Su elemento esencial es objetivo. Opera sobre prueba pericial, testimonial, confesional y documental.
\end{definicion}

\begin{definicion}{Negligencia en la producción de la prueba}
Comportamiento omisivo del litigante: falta de cuidado, diligencia o aplicación en producirla. Su elemento definitorio es subjetivo.
\end{definicion}

\begin{table}[H]
\centering
\caption{Caducidad y negligencia}
\begin{tabularx}{\textwidth}{>{\raggedright\arraybackslash\bfseries}p{0.2\textwidth} X X}
\toprule
Criterio & Caducidad & Negligencia \\
\midrule
Elemento & Objetivo (transcurso del plazo). & Subjetivo (conducta omisiva). \\
Declaración & De oficio o a pedido de parte. & Solo a pedido de parte. \\
Operación & Automática: el juez verifica el plazo. & El juez aprecia la conducta del litigante. \\
Efecto & Hace caer la prueba de la parte contra quien se declara. & El mismo. \\
\bottomrule
\end{tabularx}
\end{table}

En prueba informativa, si una entidad no contesta dentro del plazo legal (por ejemplo, 20 días para entidades públicas), el abogado debe requerir el \textbf{oficio reiteratorio}; si no lo hace, la contraria puede acusar negligencia, situación difícil de sostener. Las resoluciones que decretan negligencia o caducidad son inapelables de inmediato y la responsabilidad recae en el abogado; al apelar la sentencia puede pedirse a la Cámara el \textbf{replanteo de la prueba} mal denegada o mal declarada en negligencia, aunque las posibilidades son entonces reducidas. Si la prueba afectada es la central, la causa puede quedar comprometida; por eso el abogado debe controlar los plazos probatorios, verificar el avance de la prueba propia y monitorear la de la contraria para solicitar, en su caso, caducidad o negligencia.

\begin{cornell}{Hecho nuevo, caducidad y negligencia}
\cq{¿Qué es un hecho nuevo?}{Hecho relevante, posterior a demanda y contestación (o desconocido), conectado con el objeto; excepción a ofrecer toda la prueba al inicio.}
\cq{¿En qué plazo se denuncia?}{Cinco días desde que se conoce; límite, la audiencia del art.~360 (o cinco días antes).}
\cq{¿Caducidad o negligencia?}{La caducidad es objetiva (plazo) y puede declararse de oficio; la negligencia es subjetiva (conducta) y solo a pedido de parte.}
\cq{¿Qué evita el oficio reiteratorio?}{Que se acuse negligencia cuando una entidad no contesta en plazo.}
\cresumen{Fuera del hecho nuevo, la prueba no ofrecida a tiempo se pierde. En la producción de la prueba, la inactividad se sanciona con caducidad (objetiva) o negligencia (subjetiva), con resoluciones inapelables de inmediato; el abogado debe vigilar plazos propios y ajenos.}
\end{cornell}

\part{Nulidades, recursos y vías extraordinarias}

\chapter{La sentencia, su fundamentación y la nulidad}\label{cap:sentencia}

El proceso culmina en una decisión que debe ser razonada. Cuando no lo es, o cuando se omite un requisito esencial, el ordenamiento ofrece la nulidad. Este capítulo estudia ambos conceptos y los ilustra con casos reales.

\section{La fundamentación de las sentencias}

\begin{regla}{art.~163 CPCCN --- fundamentación}
Toda sentencia debe estar razonablemente fundada. La exigencia rige para la sentencia definitiva y para las interlocutorias; solo las providencias simples, al no resolver controversias, no requieren fundamentación extensa.
\end{regla}

La fundamentación garantiza que la decisión no sea arbitraria, permite a las partes conocer las razones del fallo y ejercer su derecho de defensa. Los fundamentos de una interlocutoria pueden ser menos extensos que los de una definitiva, pero siempre debe existir fundamento: su ausencia, o un fundamento meramente aparente, determina la nulidad. Una sentencia correcta valora la prueba paso a paso según la sana crítica (capítulo~\ref{cap:prueba}).

\begin{ejemplo}{la fundamentación en el caso de la nueva cámara Gesell}
El tribunal hizo lugar a una nueva cámara Gesell remitiéndose a la conformidad de la fiscalía y del Ministerio Público, sin exponer fundamentos propios ni remitirse expresamente a los argumentos del dictamen fiscal. La sola referencia a la conformidad de otros no es un fundamento jurídico válido; para considerar propios esos argumentos era necesario manifestarlo expresamente (caso completo en la sección~\ref{sec:gesell}).
\end{ejemplo}

\section{La nulidad procesal}

\begin{definicion}{Nulidad procesal}
Consecuencia que se produce cuando en el curso del proceso no se respeta una norma procesal o no se cumple un requisito esencial. Cuando la decisión se adopta de manera arbitraria, sin fundamento o con un fundamento meramente aparente, el acto es pasible de ser tachado de nulo.
\end{definicion}

Se distinguen tres vías ante una resolución desfavorable sobre prueba: la \textbf{apelación} se descarta, pues las cuestiones de producción o denegación de prueba son, como regla, insusceptibles de ese recurso (de admitirse, los procesos se prolongarían); la \textbf{revocatoria} es un recurso de resultado incierto; y la \textbf{nulidad} es la vía adecuada cuando se compromete el derecho de defensa. En ese planteo se articulan como agravios la arbitrariedad de la decisión y el menoscabo del derecho de defensa, por no estar las partes en igualdad frente a una resolución sin fundamento. El fundamento general invocado en el caso penal es que toda resolución debe estar fundada bajo pena de nulidad (art.~125 del Código Procesal Penal, según se citó en clase).

\section{La nulidad en cadena: notificación defectuosa, rebeldía y embargo}\label{sec:nulidad-cadena}

\begin{ejemplo}{embargo originado en una notificación defectuosa}
El cliente es una sociedad anónima demandada en un juicio laboral. La cédula se dirigió a una calle Rivadavia, domicilio que, según el contrato social presentado ante la IGJ, no es el real de la sociedad. La recibió una persona que se identificó como empleada y firmó el acuse, pero la sociedad no tiene relación con ese domicilio (probablemente el particular del dueño). El proceso siguió sin que la demandada se enterara: hubo rebeldía y finalmente un embargo de unos cien millones de pesos sobre su cuenta. Los plazos para apelar ya habían vencido, por lo que esa vía no estaba disponible.
\end{ejemplo}

\begin{importante}
\textbf{Sobre qué se pide la nulidad.} No corresponde pedir la «nulidad de la sentencia» en abstracto, sino primero la \textbf{nulidad de la notificación defectuosa}, acto que originó el vicio. Como consecuencia caen los actos posteriores que dependían de ella: la declaración de rebeldía y, en cadena, la sentencia dictada en rebeldía. El embargo en sí no es nulo, porque se trabó conforme a derecho a partir de la sentencia vigente en ese momento: lo que corresponde pedir, declarada la nulidad de lo anterior, es su \textbf{levantamiento}, porque ya no subsiste la sentencia que le dio origen.
\end{importante}

La vía es un escrito con un sumario que identifique con precisión lo pedido (nulidad de la notificación y, como consecuencia, de los actos posteriores, con levantamiento del embargo), acompañado de la prueba del vicio (aquí, el contrato social inscripto ante la IGJ, que demuestra el verdadero domicilio real). Un segundo caso: una clienta de un pueblo de Córdoba, cliente del mismo banco hacía casi cincuenta años, fue informada de que su cuenta estaba embargada por un juicio ejecutivo en el que se la notificó de manera ficta en un domicilio en el que no vivía; la solución fue sustancialmente la misma. Muchas veces el vicio no está en el fondo sino en la forma en que se notificó (véase la sección sobre domicilios del capítulo~\ref{cap:plazos}).

\section{Los errores judiciales que fundan recursos}

\begin{table}[H]
\centering
\caption{Tipos de error judicial}
\begin{tabularx}{\textwidth}{>{\bfseries\raggedright\arraybackslash}p{0.30\textwidth} X}
\toprule
Tipo de error & Descripción \\
\midrule
\emph{In procedendo} & Violación de una norma procesal que afecta el derecho de defensa; provoca la nulidad de la sentencia. \\
\emph{In iudicando} (hechos) & El juez apreció incorrectamente los hechos: lo probado difiere de lo afirmado en la sentencia. \\
\emph{In iudicando} (derecho) & El juez aplicó una norma equivocada para resolver el caso. \\
\bottomrule
\end{tabularx}
\end{table}

\begin{cornell}{Sentencia y nulidad}
\cq{¿Qué exige el art.~163 CPCCN?}{Que toda sentencia esté razonablemente fundada; sin fundamento (o con uno aparente), la resolución es nula.}
\cq{¿Para qué sirve fundar?}{Evita la arbitrariedad, permite conocer las razones y ejercer la defensa.}
\cq{¿Qué es la nulidad procesal?}{Consecuencia de omitir una norma o requisito esencial del proceso, en especial la fundamentación o la notificación válida.}
\cq{¿Sobre qué se pide la nulidad ante una notificación defectuosa?}{Sobre la notificación; caen rebeldía y sentencia; el embargo se levanta por falta de sentencia que lo sustente.}
\cq{¿Qué tipos de error hay?}{\emph{In procedendo}, \emph{in iudicando} de hecho y \emph{in iudicando} de derecho.}
\cresumen{Toda sentencia debe estar razonablemente fundada. La ausencia o apariencia de fundamento y la notificación defectuosa son causas de nulidad, que se pide con precisión técnica sobre el acto viciado y sus consecuencias. Los errores del juez pueden ser de procedimiento o de juzgamiento de hechos o derecho.}
\end{cornell}

\chapter{El sistema recursivo y las vías extraordinarias}

Una resolución que causa agravio puede revisarse. Este capítulo ordena las vías: la apelación, la queja, el efecto diferido, la reposición, el replanteo de prueba, y las vías extraordinarias (reserva del caso federal, casación y queja ante la Corte).

\section{El agravio y la apelación}\label{sec:apelacion}

Para recurrir, la resolución debe causar un \textbf{agravio}: un resultado diferente a lo pretendido. El primer paso es identificar el tipo de resolución, porque de ello depende qué recurso procede.

\begin{table}[H]
\centering
\caption{Pasos del recurso de apelación}
\begin{tabularx}{\textwidth}{>{\centering\arraybackslash}p{0.08\textwidth} X >{\raggedright\arraybackslash}p{0.30\textwidth}}
\toprule
Paso & Acción & Plazo / nota \\
\midrule
1 & Interponer el recurso (escrito simple: «vengo a apelar porque me causa agravio»). & 5 días (interlocutorias y definitivas). \\
2 & El mismo juez de primera instancia concede o deniega. & Análisis formal: tipo de resolución y plazo. \\
3A & Si concede: se fundan los agravios ante la cámara. & Trámite inmediato o diferido. \\
3B & Si deniega: recurso de queja ante la cámara para que ordene conceder. & Si no se va en queja, queda firme (cosa juzgada). \\
\bottomrule
\end{tabularx}
\end{table}

\section{La queja}

\begin{definicion}{Recurso de queja}
No discute el fondo sino exclusivamente la denegatoria de un recurso anterior. Se interpone directamente ante el tribunal superior al que denegó, que revisa la denegatoria y, si corresponde, concede el recurso rechazado.
\end{definicion}

\begin{importante}
Si se hace lugar a la queja, solo se obtiene que se conceda el recurso denegado; no se resuelve aún el fondo. Si el superior no hace lugar, la denegatoria queda firme, salvo que esté comprometido un derecho de raigambre constitucional, caso en que podría intentarse la queja ante la Corte Suprema, previo depósito que exige la reglamentación. El esquema se repite: si la cámara no concede el recurso extraordinario, se acude en queja directamente ante la Corte, con el depósito correspondiente.
\end{importante}

\section{El efecto diferido}

\begin{definicion}{Efecto diferido}
Modalidad de tramitación de la apelación concedida: no se eleva de inmediato a la cámara, sino que queda en suspenso hasta la sentencia definitiva, y solo se tramita si el apelante pierde el juicio; si gana, carece de interés práctico.
\end{definicion}

\begin{nota}
Lo que se difiere no es el efecto del recurso sino su \emph{trámite} ante la cámara, para evitar dilaciones. Se aplica, entre otros supuestos, cuando el juez rechaza un hecho nuevo y el litigante apela.
\end{nota}

\section{Las resoluciones sobre prueba: irrecurribilidad y reposición}\label{sec:reposicion}

Las cuestiones relativas a la producción y denegación de prueba son irrecurribles, inapelables de manera directa; la vía que queda es la reposición («por efecto diferido, no; reposición sí»). La vía que queda es la \textbf{reposición}: debe interponerse \emph{en el mismo acto}, con fundamento inmediato («pido reposición porque esta pregunta es determinante y el criterio para dejarla sin efecto es arbitrario»). El tribunal corre traslado a la contraparte, que puede sostener su posición u oponerse, y recién entonces resuelve «hágase lugar» o «no hágase lugar», con fundamento breve. Si no escuchara a una de las partes, podría plantearse la nulidad de lo actuado. No hacerlo en el acto significa perder la posibilidad de revertir la decisión. Al apelar la sentencia, es posible el replanteo ante la cámara de la prueba mal denegada (capítulo~\ref{cap:hecho-nuevo}).
% TODO: verificar consistencia entre fuentes originales (se afirma la irrecurribilidad de las resoluciones sobre prueba, pero también la apelación con efecto diferido del rechazo de un hecho nuevo).

\section{Reserva del caso federal y recursos extraordinarios}\label{sec:reserva-federal}

Se advierte contra el uso automático de la reserva del caso federal en los escritos: carece de sentido cuando no se sustenta en fundamentación concreta. Aproximadamente el 99\% de las demandas la incluyen, pero la mayoría no sirve porque no reserva nada específico; los jueces no objetan nada al presentarla, pero resulta inútil al recurrir ante la Corte si no se formuló correctamente.

\begin{definicion}{Reserva del caso federal (art.~14, ley 48)}
Norma que fundamenta la reserva de la vía de acceso a la Corte Suprema por recurso extraordinario federal cuando está comprometida una garantía de la Constitución Nacional o de un tratado internacional con jerarquía constitucional.
\end{definicion}

\begin{importante}
La reserva tiene sentido solo cuando puede identificarse, de manera concreta y fundada, qué artículo constitucional o tratado está comprometido y por qué. Sin esa fundamentación es una fórmula vacía.
\end{importante}

\begin{ejemplo}{reserva bien fundada y reserva sin sentido}
En la cautelar contra el decreto de aumento de las prepagas (sección~\ref{sec:cautelares}) correspondería una reserva fundada: podrían estar comprometidos el art.~14 de la Constitución Nacional (salud), el art.~17 (propiedad), el derecho a la vida, tratados con jerarquía constitucional y la propia constitucionalidad del decreto, si el Ejecutivo excedió sus atribuciones invadiendo las del Legislativo. En cambio, en la demanda de daños del curso no había violación concreta que justificara reserva, salvo un error muy grosero, como aplicar una ley manifiestamente distinta. Tampoco correspondía reservar en el planteo de nulidad del caso Gesell: lo en juego era la arbitrariedad y el derecho de defensa.
\end{ejemplo}

\subsection{El recurso de casación}

Rechazado el planteo de nulidad en el caso penal (el tribunal consideró que se garantizaba la participación de la defensa en la nueva cámara Gesell), se interpuso casación.

\begin{regla}{requisitos para la procedencia de la casación}
\begin{enumerate}
\item Acreditar la existencia de un \textbf{agravio}.
\item Que la resolución recurrida sea \textbf{arbitraria}.
\item Que sea \textbf{definitiva} o \textbf{equiparable a definitiva}.
\end{enumerate}
\end{regla}

\begin{definicion}{Resolución equiparable a definitiva}
La sentencia definitiva pone fin al proceso o es un modo anormal de terminarlo. Una resolución que no lo es puede equipararse cuando produce un efecto que no podrá repararse después. En el caso, ese efecto irreparable era someter al adolescente a una nueva pericia que lo revictimizaría, sin posibilidad de deshacer esa exposición si luego se revocara.
\end{definicion}

\begin{importante}
«El procedimiento es tanto o más importante, a veces, que el conocimiento del derecho de fondo»: este suele responder a una lógica sencilla, pero sin saber incorporarlo al proceso pierde utilidad. Para determinar si una resolución es equiparable a definitiva hay que anticipar sus efectos concretos y no limitarse a su forma.
\end{importante}

\begin{cornell}{Recursos}
\cq{¿Qué requiere recurrir?}{Un agravio; y, según el tipo de resolución, un recurso determinado.}
\cq{¿Cómo es el trámite de la apelación?}{Se interpone (5 días), el juez concede o deniega; si concede, se fundan agravios ante la cámara; si deniega, procede queja.}
\cq{¿Qué resuelve la queja?}{Solo la denegatoria; si prospera, se concede el recurso, sin resolver el fondo.}
\cq{¿Qué es el efecto diferido?}{El recurso concedido queda en suspenso hasta la sentencia y se tramita solo si el apelante pierde.}
\cq{¿Qué se hace contra una resolución sobre prueba?}{Reposición en el acto, fundada y con traslado; luego, replanteo en la cámara al apelar la sentencia.}
\cq{¿Cuándo corresponde reservar el caso federal?}{Cuando se identifica concretamente la garantía constitucional comprometida (art.~14 ley~48).}
\cresumen{Todo recurso exige agravio. La apelación puede concederse con trámite inmediato o diferido y, si se deniega, se acude en queja. Las resoluciones sobre prueba se atacan con reposición en el acto. La reserva del caso federal solo vale si está fundada; la casación exige agravio, arbitrariedad y resolución definitiva o equiparable.}
\end{cornell}

\part{Modos de terminación y otros institutos procesales}

\chapter{Modos anormales de terminación del proceso}\label{cap:modos}

Si el proceso es una cadena que va de la demanda a la sentencia, existen salidas anticipadas. Se las llama «anormales» no porque lo sean en sentido estricto, sino porque normalmente se supone que se llegará a la sentencia; los legisladores no hallaron mejor calificativo. Son cinco: allanamiento, transacción, desistimiento, conciliación y caducidad de instancia. El denominador común de los cuatro primeros es que suponen un acuerdo o una decisión voluntaria de las partes; la caducidad, en cambio, es una sanción a la inactividad.

\section{Allanamiento}

\begin{definicion}{Allanamiento}
El demandado se allana cuando reconoce el reclamo del actor. Para que sea válido, el reconocimiento debe ser total: una \textbf{aceptación lisa y llana} de la pretensión.
\end{definicion}

\begin{importante}
Si el demandado introduce un condicionamiento («sí, pero\ldots»), ese «pero» significa que no hay allanamiento.
\end{importante}

\begin{ejemplo}{por qué ocurre}
A veces las partes hablan, van, vienen, se cansan y una inicia la demanda; el otro consulta a su abogado, que le dice que no tiene cómo discutir y que le resultará más caro, y se allana.
\end{ejemplo}

\section{Transacción}

\begin{definicion}{Transacción}
Acuerdo entre las partes basado en concesiones recíprocas. Los abogados y las partes lo llevan al juez para que lo \textbf{homologue}; si no hay una cuestión de orden público que lo impida, lo homologa.
\end{definicion}

\begin{importante}
Homologado, el acuerdo adquiere efecto de \textbf{cosa juzgada} (la cuestión ya no puede discutirse) y se convierte en \textbf{título ejecutivo}.
\end{importante}

\section{Desistimiento}

\textbf{Desistimiento de la acción.} Quien desiste de la acción desiste del proceso, pero se reserva volver a demandar siempre que no opere la prescripción; por eso requiere el \textbf{consentimiento} del demandado. \textbf{Desistimiento del derecho.} El actor renuncia al derecho mismo; como no puede perjudicar a la contraparte, esta no puede impedirlo, queda liberada para siempre y, homologado, hace cosa juzgada.

\begin{ejemplo}{por qué se requiere la conformidad}
Se inicia un juicio y, unos días después, el actor desiste de la acción. El demandado puede oponerse: no quiere quedarse con esa «espada de Damocles» por un tiempo indeterminado, por si dentro de seis meses el actor se arrepiente y vuelve; ya en juicio, prefiere seguirlo y terminarlo.
\end{ejemplo}

\begin{table}[H]
\centering
\caption{Desistimiento de la acción y del derecho}
\begin{tabularx}{\textwidth}{>{\raggedright\arraybackslash\bfseries}p{0.22\textwidth} X X}
\toprule
 & De la acción & Del derecho \\
\midrule
Qué se abandona & El proceso. & El derecho mismo. \\
Nueva demanda & Posible, si no operó la prescripción. & El demandado queda liberado para siempre. \\
Conformidad del demandado & Se requiere. & No se requiere: no puede perjudicarlo. \\
Homologación & % TODO: verificar consistencia entre fuentes originales (el efecto de la homologación en el desistimiento de la acción no se precisa en las notas).
No se precisa en las notas. & Se convierte en cosa juzgada. \\
\bottomrule
\end{tabularx}
\end{table}

\section{Conciliación}

En la conciliación, como en la transacción, las partes llegan a un acuerdo; la diferencia es que en la conciliación \textbf{participa un tercero} y se promueve, muchas veces, en la audiencia preliminar. En la transacción, el acuerdo lo hacen los abogados con las partes y lo llevan al juez para su homologación. El primer paso del juez en la audiencia del art.~360, antes de fijar el objeto, es promover la conciliación (capítulo~\ref{cap:audiencia}); las partes pueden pedir una audiencia de conciliación en cualquier momento anterior a la sentencia. Cuando la controversia versa sobre la validez constitucional de una norma, no hay margen de acuerdo y no es posible conciliar.

\section{Caducidad de instancia}\label{sec:caducidad-instancia}

\begin{advertencia}
No puede operar la caducidad de instancia si no hay responsabilidad: cuando se produce, alguno de los abogados, se lo reclamen o no después, tiene una responsabilidad involucrada.
\end{advertencia}

\subsection{Fundamento y plazo}

El proceso civil y comercial, a diferencia del penal, es \textbf{impulsado por las partes}, que lo llevan etapa por etapa mediante peticiones al juez, quien siempre ordena. Cuando las partes dejan de pedir, el legislador supone que ya no tienen interés en continuar.

\begin{definicion}{Acto impulsorio}
Acto que hace avanzar el proceso hacia su etapa siguiente. No todo acto procesal lo es: una petición que no cumple esa función (por ejemplo, solicitar una copia de un expediente) no interrumpe el plazo de caducidad.
\end{definicion}

\begin{regla}{art.~310 CPCCN --- plazos de caducidad}
«Se producirá la caducidad de la instancia cuando no se instare su curso dentro de los siguientes plazos: 1) De seis (6) meses, en primera o única instancia.» Es el plazo aplicable al proceso de conocimiento ordinario.
\end{regla}

\subsection{Efectos, incidente y escrito}

La caducidad extingue el proceso, \textbf{pero no caduca el derecho}: solo se pierde la instancia (tiempo, costo económico, angustia), salvo que haya operado la prescripción. El pedido se tramita como incidente: el juez corre traslado. La parte frente a la cual se pide procurará oponerse: buscará algún acto que haya impulsado el proceso, mostrar que fue el tribunal quien lo detuvo, o discutir desde cuándo corre el plazo y si hubo días inhábiles para sostener que no se cumplieron los seis meses.
% TODO: verificar consistencia entre fuentes originales (las notas dicen «como demandados» al referirse a quien se opone a la caducidad).
Si se hace lugar, el proceso se extingue y la actora paga las costas; puede volver a iniciarse el juicio porque no se perdió el derecho.

\begin{ejercicio}{escrito de caducidad de instancia}
Pasar el conocimiento al papel es distinto de saberlo. Tres propuestas: (1) «Presento para la extinción del proceso en virtud de haber pasado los seis meses requeridos por la ley, artículo 310, sin impulso procesal, para configurarse la caducidad de instancia.» (2) «Vengo a solicitar la caducidad de instancia en los autos caratulados, por haberse cumplido el plazo perentorio de seis meses desde el último acto procesal del actor, a fojas\ldots» (3) «Que tengo presente el plazo transcurrido, sin que la contraparte impulsara acción alguna, de acuerdo al plazo del artículo 310 del Código Procesal, que es de seis meses. Vengo a solicitar la caducidad de instancia.» Se consideró perfecta la segunda; podría agregarse el sistema de cargas y la imposición de costas.
\end{ejercicio}

\subsection{Caducidad y prescripción}\label{sec:caducidad-prescripcion}

El plazo de prescripción se \textbf{interrumpe al iniciar la demanda}. Al declararse la caducidad, ese efecto desaparece y se vuelve al primer momento: si desde el hecho hasta la declaración no transcurrió el plazo de prescripción, se puede demandar de nuevo; si transcurrió, en principio se puede igualmente iniciar el juicio, pero la contraparte puede oponer la excepción de prescripción, que debe plantearse en la primera oportunidad (capítulo~\ref{cap:defensa}).

\subsection{Responsabilidad profesional}

\begin{ejemplo}{la casa con filtraciones}
Una familia tenía una casa antigua de 1920, de más de doscientos metros, con una casita al fondo que alquilaron; el ocupante cubrió cada vez más la terraza y comenzaron las filtraciones en cada lluvia. Iniciaron juicio, pero el abogado, por razón desconocida, dejó caducar la instancia. La casa, que valía casi medio millón de dólares, terminó valiendo mucho menos, al punto de tener que venderse el terreno para demolición y reconstrucción. Los clientes demandaron a su propio abogado por daños y perjuicios y lo ganaron con seguridad: no hay defensa frente a una caducidad.
\end{ejemplo}

\begin{importante}
Cuando opera la caducidad y, además, ya se cumplió el plazo de prescripción de la acción de fondo, el perjuicio puede ser irreversible: el cliente pierde el derecho. La responsabilidad recae en el abogado que dejó de impulsar y puede dar lugar a un reclamo de daños en su contra.
\end{importante}

\begin{cornell}{Modos anormales de terminación}
\cq{¿Cuáles son los cinco modos?}{Allanamiento, transacción, desistimiento, conciliación y caducidad de instancia.}
\cq{¿Qué requiere el allanamiento?}{Aceptación lisa y llana, total; un «sí, pero» no es allanamiento.}
\cq{¿Efecto de la transacción homologada?}{Cosa juzgada y título ejecutivo.}
\cq{¿Acción o derecho?}{Desistir de la acción abandona el proceso y requiere conformidad del demandado; desistir del derecho lo libera para siempre, sin conformidad y con cosa juzgada.}
\cq{¿En qué difieren conciliación y transacción?}{En la conciliación participa un tercero (p. ej., la audiencia preliminar); la transacción la acuerdan las partes y la homologa el juez.}
\cq{¿Qué efecto tiene la caducidad de instancia?}{Extingue el proceso, no el derecho; borra el efecto interruptivo de la prescripción; plazo de seis meses (art.~310) sin acto impulsorio.}
\cresumen{El proceso puede terminar antes de la sentencia por voluntad de las partes (allanamiento, transacción, desistimiento, conciliación) o por inactividad (caducidad de instancia). La caducidad extingue la instancia pero no el derecho, salvo que haya prescripción, y compromete la responsabilidad profesional.}
\end{cornell}

\chapter{Incidentes, diligencias preliminares y medidas cautelares}\label{cap:incidentes}

Junto al proceso principal existen cuestiones accesorias y medidas que lo preparan o aseguran. Este capítulo reúne los incidentes, las diligencias preliminares y las medidas cautelares.

\section{Los incidentes}

\begin{definicion}{Incidente}
Proceso secundario, vinculado al principal, que resuelve una cuestión accesoria (no definitoria del fondo) pero que incide en el avance del proceso principal.
\end{definicion}

\begin{ejemplo}{la excepción de incompetencia}
El demandado plantea al contestar que el juez es incompetente. Entre actor y demandado se genera una controversia que debe resolverse antes de continuar. Si el demandado gana, no gana el juicio: gana esa controversia parcial, que incide en el proceso pero no lo define en cuanto al fondo.
\end{ejemplo}

En épocas de expedientes en papel, el incidente se sustanciaba en un expediente separado, atado con un hilo al principal («expediente por cuerda separada», de donde «ir por cuerda»). Hoy el juez puede optar por no formar uno separado y tramitarlo dentro del principal.

\subsection{El beneficio de litigar sin gastos}\label{sec:litigar-sin-gastos}

En un juicio patrimonial, si el actor no tiene fondos para afrontar las costas, puede presentar un incidente de beneficio de litigar sin gastos, que permite el acceso a la justicia de quien carece de recursos, eximiéndolo en principio de la tasa de justicia inicial y eventualmente de las costas si se concede. Su desarrollo detallado quedó anunciado pero no figura en las fuentes; tampoco se desarrolló la prueba confesional.

\section{Diligencias preliminares}

\begin{definicion}{Diligencia preliminar}
Proceso que se inicia antes de la demanda cuando hay una prueba en peligro de desaparecer o que de otro modo no podría obtenerse, para asegurar la prueba documental necesaria antes de promover la acción principal. También sirve para averiguar datos esenciales desconocidos, como el domicilio del demandado.
\end{definicion}

\begin{advertencia}
La diligencia preliminar no debe confundirse con la prueba anticipada, que es un instituto distinto.
\end{advertencia}

\begin{ejemplo}{la historia clínica}
Quien sabe que sufrió un daño por un error médico ocurrido tiempo atrás, pero no dispone de la historia clínica para confeccionar la demanda, puede pedir al juez, de manera preliminar o conjuntamente con la demanda, que se ordene obtenerla.
\end{ejemplo}

La medida se presenta ante el juzgado que resultaría competente para la futura demanda, según las mismas reglas de competencia territorial. Si el hecho ocurrió en una clínica ubicada en una localidad distinta de la del profesional interviniente, se presenta igualmente ante el juzgado competente para la demanda de fondo.

\section{Medidas cautelares}\label{sec:cautelares}

\begin{definicion}{Medida cautelar}
Instituto procesal que protege un derecho cuando el tiempo hasta la sentencia definitiva podría hacer que deje de ser efectivo, por haber desaparecido, haberse perdido o no poder esperarse razonablemente hasta ese momento.
\end{definicion}

\begin{regla}{los tres requisitos de toda medida cautelar}
\begin{enumerate}
\item \textbf{Peligro en la demora}: riesgo concreto, evidente y actual de que la demora torne ilusorio el derecho.
\item \textbf{Verosimilitud del derecho}: no exige acreditarlo plenamente, sino que resulte creíble y presumible a partir de los elementos aportados, aunque aún no esté reconocido en sentencia.
\item \textbf{Contracautela}: garantía que ofrece quien pide la medida para responder por los daños que pudiera causar si la pidió en exceso o de manera abusiva. Puede ser personal, real o juratoria; los bienes pueden ofrecerse en especie o por un monto, con mayor flexibilidad en asuntos comerciales.
\end{enumerate}
\end{regla}

\begin{ejemplo}{cautelar contra el decreto de aumento de las prepagas}
% TODO: verificar referencia presente en las fuentes originales (el número del decreto no pudo identificarse; se lo menciona de modo genérico).
\textbf{Objeto:} medida cautelar contra la aplicación del decreto, limitada al caso de la persona afectada, sin plantear la inconstitucionalidad general. \textbf{Peligro en la demora:} de demorarse, quien debe afrontar un tratamiento médico urgente no podría seguir pagando su obra social y quedaría indefensa. \textbf{Verosimilitud:} el decreto la afecta, es titular de determinada obra social, padece una enfermedad concreta y sus ingresos no alcanzan para la nueva cuota. \textbf{Contracautela:} tercer y último requisito, una vez acreditados los otros.
\end{ejemplo}

\begin{importante}
La resolución cautelar se dicta habitualmente sin escuchar a la otra parte (\emph{inaudita parte}). Por eso se exige la contracautela: hace responsable a quien pidió la medida si lo hizo en exceso o abusivamente, y la contraparte conserva la posibilidad de recurrirla una vez dictada. Correr traslado antes suele jugar en contra de quien la solicita, pues permite ocultar bienes, modificar la situación de hecho o dilatar.
\end{importante}

\begin{ejemplo}{revinculación: cuando el juez sí corrió traslado}
En un pedido cautelar de revinculación de un padre con sus hijos (un mínimo de contacto mientras tramitaba el proceso principal), el juez se apartó del principio y, atendiendo a la naturaleza de la cuestión y al interés superior del niño, corrió traslado por tres días antes de resolver.
\end{ejemplo}

\begin{cornell}{Incidentes, diligencias y cautelares}
\cq{¿Qué es un incidente?}{Proceso secundario sobre una cuestión accesoria que incide en el principal sin definir el fondo.}
\cq{¿Qué es una diligencia preliminar?}{Proceso previo a la demanda para asegurar prueba en peligro o datos esenciales (p. ej., la historia clínica); se presenta ante el juez competente para la demanda.}
\cq{¿Qué requisitos tiene una medida cautelar?}{Peligro en la demora, verosimilitud del derecho y contracautela.}
\cq{¿Por qué se dicta \emph{inaudita parte}?}{Para evitar la sorpresa frustrada; la contracautela es la contrapartida.}
\cq{¿Traslado o vista?}{Traslado: a las partes; vista: a funcionarios no partes, como el Ministerio Público.}
\cresumen{Los incidentes resuelven cuestiones accesorias; las diligencias preliminares aseguran prueba o datos antes de demandar; las cautelares protegen el derecho durante el juicio exigiendo peligro en la demora, verosimilitud y contracautela, y se dictan por regla sin escuchar a la contraria.}
\end{cornell}

\chapter{Acciones y procesos específicos}

Algunos casos permiten ver cómo se adapta la estrategia procesal a pretensiones particulares: la acción declarativa de inconstitucionalidad, la usucapión y el reclamo por accidente laboral.

\section{La acción meramente declarativa y la inconstitucionalidad}

\begin{definicion}{Acción meramente declarativa}
Acción mediante la cual se solicita al órgano jurisdiccional que aclare o declare una situación jurídica, sin que medie necesariamente una condena de cumplimiento inmediato, para eliminar incertidumbre sobre el alcance, existencia o modalidad de una relación jurídica o la validez de una norma. Es la herramienta idónea para impugnar la constitucionalidad de una norma cuando hay un caso concreto, sin esperar un daño consumado.
\end{definicion}

Requiere un \textbf{caso concreto}: una parte afectada de manera directa, actual y específica; no puede plantearse en abstracto. La aplicación al caso del Consejo Federal Pesquero se estudia en la sección~\ref{sec:pesquero}. Allí la controversia es exclusivamente normativa, por lo que se convierte en cuestión de puro derecho sin etapa probatoria relevante, y no es posible proponer conciliación (sección~\ref{sec:conciliacion-pesquero}).

\section{La usucapión}\label{sec:usucapion}

\begin{definicion}{Usucapión (prescripción adquisitiva)}
Derecho real por el cual el transcurso del tiempo permite que quien poseyó una cosa como dueño pida a un juez el reconocimiento de esa titularidad, aunque no haya sido titular registral originario, por haber ejercido los derechos correspondientes (pagando impuestos y comportándose como propietario) durante el plazo legal.
\end{definicion}

\begin{regla}{arts.~1897 y 1899 del Código Civil y Comercial de la Nación}
Art.~1897: «La prescripción para adquirir es el modo por el cual el poseedor de una cosa adquiere un derecho real sobre ella, mediante la posesión durante el tiempo fijado por la ley.» Art.~1899: «Si no existe justo título o buena fe, el plazo es de veinte años [\ldots].» (Libro Cuarto, Título I, Capítulo~2.)
\end{regla}

Frente a una demanda de usucapión el demandado puede oponer que no se cumplió el plazo, que la posesión no existió como se relata o que fue interrumpida. Si la sentencia es favorable, la titularidad se inscribe en el Registro de la Propiedad Inmueble a nombre del actor, que podrá disponer del bien como si lo hubiese comprado. Es una figura frecuente en el interior (campos, chacras, pequeños terrenos) y menos habitual en la Ciudad de Buenos Aires, por su mayor organización registral.

\begin{ejemplo}{el campo cercano a La Plata: cuando la mediación no alcanza}
Un campo próximo a La Plata era ocupado desde hacía unos cuarenta años por personas que desarrollaban un emprendimiento de cultivo y habían pagado casi todos los impuestos (uno de los elementos que más acredita la usucapión); era prácticamente imposible que los titulares registrales originales siguieran con vida. No correspondía una mediación previa: la usucapión no es un derecho disponible por acuerdo. Aun reuniendo a herederos, transferir la titularidad exigiría otro acto jurídico (donación o venta) que no reemplaza la declaración judicial de los requisitos de posesión y tiempo.
\end{ejemplo}

\begin{table}[H]
\centering
\caption{Usucapión y juicio de escrituración}
\begin{tabularx}{\textwidth}{X X}
\toprule
\textbf{Usucapión} & \textbf{Juicio de escrituración} \\
\midrule
Se funda en la posesión y el tiempo (20 años), exista o no boleto. & Se funda en el boleto; exige acreditar la obligación de escriturar asumida por el vendedor. \\
Depende de la prueba de actos posesorios (impuestos, expensas, uso, testigos). & Depende de la prueba del vínculo contractual y del incumplimiento. \\
Útil cuando el vendedor desapareció, se disolvió o es inhallable. & Requiere, en principio, dirigir la acción contra quien se obligó (o sus sucesores). \\
Se elige si se cuenta con todos los elementos para probar la posesión. & Se elige si la relación contractual y la identidad del obligado son claras. \\
\bottomrule
\end{tabularx}
\end{table}

Con boleto, el actor puede optar entre ambas vías, que pueden conducir al mismo resultado; la elección es estratégica y depende de la documentación y prueba disponibles. La audiencia preliminar de un caso de usucapión se describe en el capítulo~\ref{cap:audiencia}.

\section{El accidente laboral, la ART y la revisión judicial}

\begin{definicion}{Accidente laboral}
Aquel que se produce con ocasión del trabajo o durante el trayecto entre el domicilio y el lugar de trabajo, en cualquier sentido.
\end{definicion}

\begin{ejemplo}{el trabajador del balancín}
Un trabajador se corta un dedo con un balancín y acude a la Comisión Médica, órgano de la Superintendencia de Riesgos del Trabajo, para que un médico fije un porcentaje de incapacidad, que suele ser reducido o directamente no reconocido; inicia entonces un juicio contra la ART reclamando la indemnización que estima que le corresponde.
% TODO: en el apunte original se menciona a la «Superintendencia del Trabajo»; las propias notas ubican a la Comisión Médica dentro de la Superintendencia de Riesgos del Trabajo, según las notas originales.
\end{ejemplo}

\begin{importante}
La decisión de la ART, por ser de carácter administrativo, no cierra el conflicto: toda decisión administrativa debe poder revisarse en sede judicial, de lo contrario se privaría de defensa en juicio. El plazo para impugnarla suele ser breve; si no se ejerce, la resolución queda firme, aunque la posibilidad existió.
\end{importante}

Frente a una resolución desfavorable de la Comisión Médica, se espera a que quede formalizada y, dentro del plazo, se opta por la apelación o por el inicio directo del juicio.
% TODO: verificar consistencia entre fuentes originales (se mencionan alternativamente un plazo de apelación de quince días y de dos días; se indica que ambas vías insumen un inicio de siete u ocho meses).

En el fuero laboral existen audiencias de conciliación, distintas de la audiencia preliminar civil, que algunos jueces convocan aun sin pedido. El juez, que ya conoce el caso y las pruebas, puede orientar a las partes hacia un acuerdo sin adelantar su decisión, mostrando hacia dónde se dirige el proceso, pues un acuerdo suele ser más económico, simple y rápido que continuar hasta la sentencia. La pericia médica en este tipo de casos se describió en el capítulo sobre prueba pericial.

\begin{cornell}{Acciones y procesos específicos}
\cq{¿Para qué sirve la acción meramente declarativa?}{Para eliminar incertidumbre sobre una relación jurídica o la validez de una norma, con un caso concreto.}
\cq{¿Qué es la usucapión y por qué exige juicio?}{Adquisición de un derecho real por posesión durante el tiempo legal (20 años sin justo título o buena fe); no es negociable, requiere declaración judicial.}
\cq{¿Usucapión o escrituración?}{Usucapión: posesión y tiempo; escrituración: boleto e incumplimiento de escriturar. La elección es estratégica.}
\cq{¿Por qué la decisión de la ART es revisable?}{Porque es administrativa y toda decisión administrativa debe poder revisarse judicialmente, dentro del plazo.}
\cresumen{Cada pretensión adapta la estrategia: la acción declarativa exige caso concreto, la usucapión requiere declaración judicial basada en posesión y tiempo, y las decisiones administrativas, como las de la ART, son siempre revisables judicialmente dentro de plazo; en lo laboral, las audiencias de conciliación pueden ahorrar tiempo y costos.}
\end{cornell}

\part{Casos integradores y formación profesional}

\chapter{Casos integradores}

Los casos reales permiten ver cómo se combinan los institutos estudiados. Se presentan tres: el accidente náutico del río Tigre, la restitución internacional de menores con la cámara Gesell, y el Consejo Federal Pesquero.

\section{El accidente náutico en el río Tigre}\label{sec:tigre}

Es un caso real, en trámite, que sirve para aprender a transformar un relato de hechos en la estructura de una demanda de daños y perjuicios. Un kayakista fallece a raíz de la colisión de su embarcación con una lancha en el río Tigre. Importa el monto menos que saber qué rubros se reclaman y, sobre todo, cómo se prueban.
% TODO: verificar consistencia entre fuentes originales (fecha, personas involucradas, días de internación y cantidad y sexo de los tripulantes del kayak difieren entre apuntes; se transcriben los datos tal como figuran en cada uno).

\subsection{Los hechos}

El hecho ocurrió pocos días antes del Carnaval de Río, en febrero de 2019, en el río Tigre, en la zona de la estación fluvial donde salen las lanchas de pasajeros, cerca del puente de la Ruta 60, donde hay numerosos clubes de remo. Una lancha de pasajeros (de las típicas embarcaciones marrones), amarrada junto a otras a la espera de zarpar, inició una maniobra de retroceso para acomodarse. Viajaban dos tripulantes: el \textbf{patrón}, que conduce, y el \textbf{marinero}, que debe ubicarse en la popa durante las maniobras para advertir obstáculos; el marinero estaba adelante, conversando con el patrón, en lugar de estar atrás como vigía.

En simultáneo se aproximaba un kayak doble de un club de remo, con dos remeros (un hombre y una mujer, ambos abogados y profesores) que navegaban correctamente por su derecha, conforme a normas de navegación similares a las de tránsito vehicular. Al retroceder sin control, la lancha colisionó por estribor con la popa del kayak, que volcó. El remero de atrás salió nadando. La remera de adelante fue succionada por la hélice, pues una lancha, a diferencia de un vehículo terrestre, no se detiene de inmediato al colisionar y genera un efecto de succión; su cuerpo fue arrastrado y pasado por encima de la embarcación.

Un catamarán cercano comenzó a tocar la campana y hacer señales de alarma. Una pasajera de otra lancha, al ver que la mujer flotaba sin que nadie actuara, se arrojó al agua y la remolcó hasta que pudo ser subida a una embarcación. Fue trasladada a un hospital, donde permaneció internada dos o tres días y falleció por una sepsis generalizada derivada de las lesiones. El otro tripulante no requirió atención médica de urgencia; un médico, pasajero de otra lancha, le realizó maniobras de reanimación sobre el muelle.
% TODO: verificar consistencia entre fuentes originales (se afirma que el tripulante no requirió atención de urgencia y, a la vez, que recibió maniobras de reanimación; otro apunte menciona cinco días de internación en Pacheco y lesiones leves de la otra persona).

\subsection{Reconstrucción mediante preguntas}

Así debería procederse al entrevistar a un cliente. Se estableció: el hombre salió por sus propios medios; la mujer fue rescatada y estuvo internada hasta fallecer por sepsis; hay por lo menos cinco testigos presenciales clave; la mujer tenía alrededor de cuarenta años, era abogada y ejercía la docencia; los remeros estaban acreditados en el club y para salir con la embarcación deben aprobar un examen de idoneidad; la señora tenía tres hijos, y debe averiguarse de qué trabajaban las víctimas, si estaban casados y convivían.

\subsection{Mecánica del hecho y responsabilidad}

Lo que influye en la mecánica es el procedimiento del marinero y del patrón: realizaban una \textbf{maniobra peligrosa} (ir hacia atrás para acomodarse) sin los recaudos necesarios. Para ella hay que emitir una señal (sonido, luz) y, sobre todo, debe haber un marinero en la popa que vea y avise al conductor, con función preventiva. De haber estado en su puesto, el accidente no habría sucedido.

\begin{importante}
Hechos especialmente relevantes para la demanda: (i) la maniobra peligrosa de retroceso sin señales ni recaudos; (ii) la ausencia del marinero en su puesto de vigía; (iii) que el kayak navegaba correctamente por su derecha; (iv) el efecto de succión de las embarcaciones a motor, que agravó las lesiones.
\end{importante}

En cuanto al \textbf{factor de atribución}, el marinero tiene, en principio, menos responsabilidad que el patrón, que conducía. El kayak era del club de remo; la lancha pertenece a una empresa de transporte fluvial con recorrido regular, como una empresa de colectivos. Los rubros mencionados son daño moral y daño psicológico (capítulo~\ref{cap:danos}). Las contrapartes estaban dificultando el trámite.

\subsection{Demandados, objeto y relato}

Los sujetos demandados, el objeto y los modelos de relato se estudian en el capítulo~\ref{cap:estructura-demanda}: patrón, marinero, empresa y aseguradora citada en garantía (art.~118 Ley de Seguros). Los hechos de la demanda se ordenan distinguiendo los del accidente y los que fundan cada rubro.

\subsection{Elección del foro}

\begin{table}[H]
\centering
\caption{Opciones de jurisdicción según los apuntes}
\begin{tabularx}{\textwidth}{>{\raggedright\arraybackslash}p{0.55\textwidth} X}
\toprule
\textbf{Opción} & \textbf{Fuero y asiento} \\
\midrule
1. Justicia provincial ordinaria & Civil y comercial ordinaria, San Isidro (Provincia de Buenos Aires). \\
2. Justicia federal con asiento en la Provincia & Fuero federal, San Martín. \\
3. Justicia nacional ordinaria & Civil y comercial, Capital Federal. \\
4. Justicia federal con asiento en la Capital & Fuero federal, Capital Federal. \\
\bottomrule
\end{tabularx}
\end{table}

Las aseguradoras tienen domicilio en la Capital, lo que habilita decidir entre justicia nacional y provincial; al ser un hecho ocurrido en Tigre, podría iniciarse ante los tribunales de San Isidro. En el modelo analizado se optó por la justicia nacional, por el domicilio legal de la empresa en la Ciudad. Otro apunte enumera cuatro foros posibles: tribunales nacionales ordinarios (CABA), federales, San Isidro y arbitrales. La opción federal existía por tratarse de un abordaje en un río (Ley de Navegación); tramitaba además un juicio contra la misma empresa en sede federal en San Isidro, y los jueces no lograron ponerse de acuerdo sobre el conflicto de competencia. Sobre el componente emotivo, véase la sección~\ref{sec:emotivo-fuero}; sobre la relación con la causa penal «N.N. s/ homicidio culposo», la sección~\ref{sec:civil-penal}.
% TODO: verificar consistencia entre fuentes originales (el fuero federal figura con asiento en San Martín en un apunte y en San Isidro en otro).

\subsection{Estrategia probatoria}

Se distinguió entre probar el hecho y probar los daños. \textbf{Hecho:} testimonial (personas del catamarán que habrían oído bocinas y pedidos de auxilio, y un testigo que se giró en el momento del impacto); informativa mediante oficios a la municipalidad (cámaras de seguridad), a Gendarmería o Prefectura Naval (cámaras y registros), al organismo de emergencias médicas que atendió a la víctima y a la aseguradora (póliza); documental (libreta de casamiento, acta constitutiva del estudio jurídico y libro contable, habilitación del catamarán ante la Prefectura, carnet náutico del patrón u oficio al club o Prefectura, comprobantes de gastos médicos, expediente penal o IPP); y una pericia mecánica en embarcaciones, que podría resultar sobreabundante. \textbf{Daños:} para daño moral, testigos que den cuenta de la vida de la pareja antes y después; para daño psicológico, pericia psicológica con puntos adecuados y relación causal; para pérdida de chance y lucro cesante, facturaciones del estudio (más que recibos de sueldo, por tratarse de actividad independiente) y testigos socios o compañeros. Los ingresos pueden acreditarse también por informativa (capítulo~\ref{cap:danos}).

En este caso, la responsabilidad del patrón fue caracterizada como subjetiva y la de la empresa como objetiva; los patrón y marinero fueron condenados en sede penal, no la empresa.

\section{La restitución internacional de menores y la cámara Gesell}\label{sec:gesell}

Es un caso real y complejo, usado para ilustrar nulidad, casación y queja (capítulos~\ref{cap:sentencia} y siguientes) y las medidas cautelares.

\begin{ejemplo}{los hechos}
Un matrimonio integrado por una mujer francesa y argentina y sus dos hijos, con residencia alternada entre la Argentina y Madrid, atraviesa una crisis; la madre viajaba con asiduidad con los hijos, pues ambos progenitores tenían poderes para hacerlo. El padre viaja por trabajo a Uruguay y, al regresar, la madre se había ido con los dos hijos sin que pudiera ubicarlos; tras algunos contactos por videollamada, la comunicación se interrumpió. El padre inició en la Argentina un juicio de restitución internacional ante un juzgado civil de familia de la Capital, por ser el domicilio de la familia, y el juez ordenó la restitución. En Francia se sustanciaron varios procesos (un tribunal de primera instancia, donde la familia materna sostenía que el domicilio conyugal era argentino, y una cámara de apelaciones que confirmó el criterio); en ambos un informe concluyó que no podía presumirse abuso del padre hacia uno de los hijos, denunciado por la madre. La Cámara de Casación de París confirmó que los niños debían regresar a la Argentina.
\end{ejemplo}

\begin{ejemplo}{el proceso penal en Francia y el regreso}
La madre formuló además una denuncia penal en Francia. En el sistema francés la policía judicial no cumple funciones de policía de calle sino que asiste al fiscal y al juez; el fiscal concluyó que no hay delito y ordenó el archivo. Ese archivo no equivale al sobreseimiento argentino, que dicta un juez. Resuelto el proceso civil, el padre retiró a los hijos y regresó a Buenos Aires. Una semana después la madre presentó en la Argentina una denuncia penal casi idéntica, con partes editadas de las filmaciones de cámara Gesell hechas en Francia. Una jueza civil argentina ordenó entonces la restitución de los menores a la madre, y le retiraron la custodia al padre. Desde entonces transcurrieron cinco años sin que el padre viera siquiera una fotografía de sus hijos. Respecto de la niña se dictó sobreseimiento total y firme; respecto del hijo varón se elevó a juicio.
\end{ejemplo}

En el incidente de revinculación se ofrecieron todas las garantías posibles (custodia policial, asistentes sociales, acompañamiento psicológico); la contraparte y el Ministerio Público se opusieron, y el juez rechazó la revinculación invocando el interés superior del niño y calificando lo pedido como medida cautelar, sin que estuvieran dadas las condiciones. Frente a ello se formuló la pregunta retórica de la sección~\ref{sec:retorica}.

\subsection{La nueva cámara Gesell y la nulidad}

Como defensor del padre se logró que no permaneciera detenido y que se incorporaran a la causa argentina todos los expedientes franceses, con las filmaciones completas. Se habían realizado tres cámaras Gesell en el exterior, incorporadas por vía diplomática, y dos en el país que no resultaron útiles porque el menor prácticamente no pudo declarar. Faltando quince días para el juicio, la querella presentó certificados médicos y pidió una nueva cámara Gesell; el período de ofrecimiento de prueba ya estaba cerrado. El pedido debía sustanciarse con vista al fiscal, vista al Ministerio Pupilar (por tratarse de un menor) y traslado a la defensa.

Fiscal y Ministerio Pupilar acompañaron el pedido. La defensa se opuso con dos argumentos: que había operado la preclusión, sin justificación para no haberlo pedido antes (en particular al incorporarse las traducciones de las cámaras del exterior), y que los certificados de profesionales designados por la querella no eran imparciales como los del cuerpo médico forense; y señaló el riesgo de revictimización de un menor. El tribunal hizo lugar con una resolución de cuatro líneas fundada solo en la conformidad de fiscalía y defensora de menores. Se descartó la apelación y la revocatoria, y se planteó la \textbf{nulidad} por falta de fundamentación y afectación del derecho de defensa; el tribunal la rechazó considerando que se garantizaba la participación de la defensa en la nueva cámara. Se interpuso casación (sección~\ref{sec:reserva-federal}). Según relato, el juicio oral aún no se había realizado y se agregó un breve intercambio sobre los poderes de ambos progenitores para viajar con los menores y cómo la práctica notarial fue exigiendo una fecha cierta de regreso.
% TODO: verificar consistencia entre fuentes originales (un apunte menciona la intervención de la defensora de menores y otro, del Ministerio Pupilar y de la fiscalía; el cierre sobre los poderes no se desarrolla).

\begin{definicion}{Preclusión en el caso}
Vencida la etapa de ofrecimiento de prueba para el debate, la nueva cámara Gesell ya no podía pedirse, salvo excepciones expresas de la ley (sección~\ref{sec:preclusion}).
\end{definicion}

\section{El Consejo Federal Pesquero}\label{sec:pesquero}

El Consejo Federal Pesquero es un organismo creado por ley en 1998 (Ley 24.922, Ley Federal de Pesca), integrado por representantes de las provincias con costa marina (Buenos Aires, Río Negro, Chubut, Santa Cruz y Tierra del Fuego) y del Ministerio de Economía, Agricultura y Pesca. Determina anualmente qué especies se pescan y asigna cuotas a cada empresa. Una empresa que cumple los requisitos fue excluida de la cuota de merluza (en otro apunte, merluza negra) sin licitación pública, procedimiento transparente que convoca oferentes y adjudica al que mejor cumple en calidad y precio.
% TODO: verificar consistencia entre fuentes originales (un apunte indica que la ley establece cuotas mediante licitación o mecanismo igualitario y que el decreto reglamentario permite otorgarlas por resolución administrativa; otro, que la ley misma otorga la facultad de asignarlas; no se precisa el artículo del decreto).

Los recursos de pesca son bienes de dominio público; su explotación debe garantizar igualdad, no discriminación y transparencia, y el mecanismo viola esos principios porque favorece históricamente a las mismas empresas.

\textbf{Estrategia.} Se inició una acción meramente declarativa de inconstitucionalidad contra la ley y su reglamentación. El objeto no es obtener la cuota sino que se declare la inconstitucionalidad del sistema discrecional, de modo que el Estado deba establecer un mecanismo licitatorio con igualdad de condiciones. Se plantea que la ley del Congreso es inconstitucional porque crea un organismo con facultades de asignación discrecional de recursos públicos, y que la reglamentación del Ejecutivo agrava esa inconstitucionalidad; quien ejerce una actividad lícita tiene derecho a acceder a bienes públicos en igualdad. Se demandó al Estado Nacional (acto legislativo) y al Poder Ejecutivo (reglamentación). El Estado cuestionó su legitimación pasiva, sosteniendo que el legitimado era el Consejo; el juez incorporó al Consejo como \textbf{codemandado}.

\begin{ejercicio}{tarea de análisis}
Contestada la demanda por el Consejo Federal Pesquero, ¿habrá audiencia del art.~360? Justifíquese desde el tipo de controversia.
\end{ejercicio}

\label{sec:conciliacion-pesquero}La controversia es exclusivamente normativa: no hay hechos empíricos controvertidos sobre cómo opera el sistema de cuotas sino si es compatible con la Constitución; es cuestión de puro derecho, sin etapa probatoria relevante, y no cabe conciliación porque, tratándose de la validez constitucional de una norma, no hay margen de acuerdo.

\begin{cornell}{Casos integradores}
\cq{¿Qué hechos del caso Tigre son decisivos?}{Maniobra peligrosa de retroceso, marinero fuera de su puesto de vigía, kayak navegando por su derecha, succión de la hélice.}
\cq{¿Qué foros eran posibles?}{Justicia provincial ordinaria (San Isidro), federal con asiento en la Provincia, nacional ordinaria y federal con asiento en la Capital.}
\cq{¿Por qué el pedido de nueva cámara Gesell era cuestionable?}{Había precluido la prueba y había riesgo de revictimización; el tribunal resolvió sin fundamento propio.}
\cq{¿Qué se persiguió en el caso del Consejo Pesquero?}{Una declaración de inconstitucionalidad del sistema discrecional de cuotas, para obligar a un mecanismo licitatorio.}
\cresumen{Los tres casos combinan competencia, demanda, prueba, nulidad, recursos y cautelares: el accidente del río Tigre integra responsabilidad, rubros y prueba; la restitución internacional con cámara Gesell muestra nulidad y casación por falta de fundamentación y preclusión; el Consejo Pesquero muestra la acción declarativa y la cuestión de puro derecho.}
\end{cornell}

\chapter{El oficio del abogado y su formación}

El conocimiento de las normas, por sí solo, no alcanza. Este capítulo final recoge las reflexiones sobre cómo se aprende y ejerce el oficio, el papel de la inteligencia artificial y las habilidades que perduran.

\section{El proceso como oficio}

En educación se distingue el \textbf{conductismo}, que lleva al estudiante de la mano, del \textbf{constructivismo}, que en lo ideal da herramientas para que construya conocimiento. Pero el traspaso en la música, el derecho o la medicina tiene que ver con la \textbf{transmisión y la experiencia}: se aprende mirando, escuchando e imitando. El proceso es un \textbf{oficio} y, como todo oficio, tiene que ver con las manos, aunque sea intelectual. Hablar solo de lo que dijeron otros autores (Quintiliano, por ejemplo) no deja de ser una cadena de transmisión.

\begin{importante}
El gran tecnicismo del abogado es manejar un proceso: escuchar a quien tuvo o tiene un conflicto y ver cómo llevarlo al tribunal. Es una función intelectual: crear una estructura para llevar adelante una estrategia y convencer al tercero que resolverá, con reglas que se construyen a través del proceso.
\end{importante}

\begin{ejemplo}{el cuaderno que se llevan delante de todos}
Se puede saber que alguien se lleva abiertamente el cuaderno de otra persona; pero hay que encontrar el mecanismo (formarlo, plantearlo, llevarlo) y convencer al juez de que resuelva en consecuencia. Si no, no sirve.
\end{ejemplo}

Son los abogados quienes impulsan los cambios en la práctica: los jueces no modifican la práctica por iniciativa propia. Si la carga dinámica de la prueba se invoca constantemente, argumentando por qué esa parte está en mejor posición de probar, termina consolidándose como jurisprudencia, en especial en materia laboral, donde el código civil procesal se aplica supletoriamente y el empleado suele estar en desventaja informativa. Un pedido «no usual», como una prueba extraordinaria, depende de cómo se lo plantee.

\section{Estudiar demandas y contestaciones}

A un fallo se llega porque hubo quienes litigaron. El fallo muestra lo que dicen los jueces que se dijo, pero no cómo se llegó a él: allí está el razonamiento de cada parte, qué se dijo, qué se probó y qué pensó quien armó el litigio. Conviene estudiar las demandas y contestaciones.

\begin{ejemplo}{un litigio en preparación}
Se prepara un planteo de inconstitucionalidad de la ley de emergencia pesquera, vigente desde hace más de treinta años; hay unos seiscientos empresas pesqueras que no pueden hacer muchas cosas y ninguna lo planteó. Si prospera, se cae la ley y todo se redefine, pero el fallo no mostrará a quien litigó. Causa: «Pro Sur contra Estado Nacional», expediente 50047/25.
% TODO: la numeración del expediente no resulta del todo clara en las notas originales.
\end{ejemplo}

\section{Inteligencia natural e inteligencia artificial}

Antes de recurrir a la inteligencia artificial es necesario tener incorporada una base de inteligencia natural sobre el funcionamiento del proceso; si no se sabe qué preguntar, el riesgo de una respuesta desacertada es extraordinario.

\begin{importante}
La inteligencia artificial es valiosa para quien ya cuenta con criterio propio; cuando no se comprende lo que se produce, juega en contra de la formación. «Es como la gimnasia sin esfuerzo: no existe.» La labor del abogado es intelectual: el entrenamiento consiste en desarrollar la propia manera de pensar y argumentar. Poseer datos vale hoy muy poco, porque el acceso a la información es universal; lo que no se resuelve apretando un botón es cómo manejarla.
\end{importante}

Se plantea el temor de los jóvenes abogados, de dos a cinco años de recibidos, a ejercer de forma independiente frente a la inteligencia artificial; se destaca que, aunque existía hacía tres o cuatro años, recién el año anterior hubo un salto cualitativo y un acceso masivo, de alcance mundial. Se lo compara con la implementación del expediente digital, que dejó fuera a quienes no se adaptaron. Se vincula con la devaluación del título de grado: hoy es necesario avanzar hacia posgrado (una maestría, como mínimo). Se señaló que una parte de los trabajos prácticos no se elaboró con criterio propio, y que quienes egresan sin criterio tendrán mayores dificultades, mientras que cultivarlo es el diferencial.

\begin{ejemplo}{las pruebas PISA}
Las últimas pruebas PISA evidenciaron un deterioro cognitivo en la Argentina, ubicada en el noveno lugar entre dieciséis países de Latinoamérica. Un ítem planteaba decidir entre dos planes de telefonía celular (uno con monto fijo y adicional por gigabyte extra, y otro con costo base menor y costo mayor por gigabyte extra) para un consumo dado: requería solo sumas y restas, y la mayoría no pudo resolverlo. Las diferencias de comprensión lectora entre quienes sostienen hábitos de estudio profundos y quienes no resultarán significativas aun con el mismo título.
\end{ejemplo}

Se identifica además el uso intensivo de redes sociales (contenidos breves de consumo inmediato) como factor de atención fragmentada, en detrimento de la lectura extensa. La disponibilidad inmediata del conocimiento, como tener un gimnasio en casa, puede generar la comodidad de postergarlo; no resulta claro si es la necesidad la que impulsa a superarse o la abundancia la que perjudica. Existe también un margen de compromiso individual, pues la suma de decisiones individuales conforma la sociedad.
% TODO: contenido incompleto o ambiguo en las notas originales (no se desarrolla el alcance del «compromiso individual»).

\section{Habilidades que perduran}

¿Qué sirve para hoy y probablemente dentro de diez años? Buena parte de las formaciones útiles hace diez o veinte años quedaron obsoletas y el conocimiento memorizado es relativo: «apretamos un botón y la información está». Lo que persiste es el manejo criterioso de la información, no la acumulación de datos: hacer la pregunta correcta, verificar los datos antes de actuar, no dar por cierto lo que solo parece serlo, y aplicar la norma con método y en el momento oportuno.

\begin{cornell}{El oficio del abogado}
\cq{¿Por qué el proceso es un oficio?}{Porque se aprende por transmisión y experiencia, con reglas operativas, aunque sea una tarea intelectual.}
\cq{¿Por qué estudiar demandas y contestaciones?}{Porque allí está el razonamiento de las partes que el fallo no muestra.}
\cq{¿Qué lugar tiene la inteligencia artificial?}{Es valiosa para quien tiene criterio propio; sin él juega en contra: no hay gimnasia sin esfuerzo.}
\cq{¿Quién impulsa los cambios en la práctica judicial?}{Los abogados, planteando constantemente una necesidad (por ejemplo, la carga dinámica).}
\cresumen{El abogado ejerce un oficio intelectual que se aprende practicando y observando cómo razonan las partes. Las herramientas tecnológicas no reemplazan el criterio propio: lo que perdura es manejar la información con método, hacer la pregunta correcta y verificar antes de actuar.}
\end{cornell}

\appendix
\chapter{Normativa y referencias citadas en las fuentes}

El siguiente cuadro reúne las normas citadas, con la indicación del tema que ilustran. Las referencias se conservan tal como figuran en las fuentes, sin completarlas ni corregirlas.

\begin{longtable}{>{\raggedright\arraybackslash}p{0.34\textwidth} >{\raggedright\arraybackslash}p{0.58\textwidth}}
\toprule
\textbf{Norma} & \textbf{Tema} \\
\midrule
\endfirsthead
\toprule
\textbf{Norma} & \textbf{Tema} \\
\midrule
\endhead
\bottomrule
\endfoot
Constitución Nacional, arts.~14, 17, 116 y 117 & Salud y propiedad (reserva del caso federal); competencia federal. \\
Ley 48, art.~14 & Reserva del caso federal. \\
CPCCN, arts.~90 y 91 & Intervención de terceros (enunciados). \\
CPCCN, art.~133 & Resoluciones que se notifican por cédula o personalmente. \\
CPCCN, arts.~160 a 163 & Resoluciones judiciales; fundamentación de sentencias. \\
CPCCN, art.~310 & Caducidad de instancia (seis meses). \\
CPCCN, art.~330 & Contenido de la demanda. \\
CPCCN, art.~347 & Excepciones previas (texto abreviado). \\
CPCCN, art.~356, inc.~1 & Reconocimiento o negativa de hechos y documentos (texto abreviado). \\
CPCCN, art.~357 & Reconvención (texto abreviado). \\
CPCCN, art.~359 & Cuestión de puro derecho. \\
CPCCN, art.~360 & Audiencia preliminar. \\
CPCCN, art.~363 & Notificación de la audiencia; hechos nuevos. \\
CPCCN, art.~377 & Carga de la prueba (texto abreviado). \\
CPCCN, arts.~388 y 389 & Exhibición de documentos; documentos en poder de terceros. \\
CPCCN, art.~400 & Oficios. \\
CPCCN, art.~432 & Caducidad de la prueba testimonial por falta de notificación. \\
CPCCN, arts.~457, 458 y 459 & Prueba pericial, perito único, puntos de pericia. \\
CPCCN, art.~626, inc.~3 & Incapacidad e inhabilitación (remisión del art.~458). \\
CPCCN, art.~139 (según un apunte) & Cuestión de puro derecho. % TODO: verificar referencia presente en las fuentes originales.
\\
Código Procesal Penal, art.~125 (según las fuentes) & Fundamentación bajo pena de nulidad. \\
Código Civil y Comercial, arts.~1897 y 1899 & Usucapión. \\
Ley de Seguros 17.418, art.~118 & Citación en garantía. \\
Ley 24.922 & Ley Federal de Pesca; Consejo Federal Pesquero. \\
\end{longtable}

\end{document}
